% =====================================================================
% 第十三讲 (Exposé XIII) 集合层与非交换群层的上同调真性
% (Propreté cohomologique des faisceaux d'ensembles et des faisceaux
%  de groupes non commutatifs)
% M. Raynaud
% 原书页码: pp. 259-322
% =====================================================================
\chapter{集合层与非交换群层的上同调真性}

\noindent Mme M. Raynaud\footnote{按 A. Grothendieck 的未发表笔记.}

本讲旨在用平展上同调推广 IX 和 X 的某些结果. 还说明如何把 SGA 5 II 中对这类层仍有意义的结果推广到未必交换的群层. 假定已知 SGA 4 中阐述的平展上同调概念.

主要结果 (\cref{thm:XIII.2.4}) 给出非紧合态射 $f \colon U \to S$ 的重要例子，它在"维数 $\leqslant 1$ 时上同调真"，i.e. 对 $U$ 上某些 (平展拓扑意义的) 群层 $F$，$f_* F$ 和 $R^1 f_* F$ 的构造与任何基变换 $S' \to S$ 交换. 这个性质确实被 $S$ 上紧合概形 $X$ 的、相对 $S$ 为正规交叉除子 $D$ 的余集的开集 $U$ 满足，至少当对 $F$ 要求它常有限、阶与 $S$ 的剩余特征互素时. 若不再假设 $F$ 的阶与 $S$ 的剩余特征互素，有类似结果，只要把 $R^1 f_* F$ 换成子层 $R^1_t f_* F$ (只考虑 $F$ 下"相对 $S$ 在 $X$ 上温和分歧"的挠子). 特别这可以用来证明可分闭域上紧合光滑代数曲线去掉有限个闭点后所得曲线的\emph{温和分歧基本群}拓扑有限型 (\cref{cor:XIII.2.12}).

第 4 号致力于同伦正合序列和 Künneth 公式.

最后附录给出 X.3.6 中证明的 Abhyankar 引理的有用变体.

\setcounter{section}{-1}

\section{叠理论的回顾}

下面将用 \cite{exp13-1} 和 \cite{exp13-2} 中阐述的叠理论. 限于概形的平展位点. 给定概形 $X$，用 $X_{\mathrm{ét}}$ 记 $X$ 的平展位点. 回忆 $X$ 上的叠 $F$ 是 $X_{\mathrm{ét}}$ 上的纤维范畴\index[term]{die@叠}，使对每个平展于 $X$ 上的概形 $X'$ 和纤维 $F_{X'}$ 的每对对象 $x$，$y$，预层 $\mathrm{Hom}_{X'}(x, y)$ 是层，且使对每个平展满射态射 $X'' \to X'$，$F_{X''}$ 的配备相对 $X'' \to X'$ 的下降数据的每个对象是 $F_{X'}$ 的某个对象的逆像.

用 $F(X')$ 记 $F/X'$ 的笛卡儿截面的范畴. 更一般地，若 $\mathbf{Sch}_X$ 是配备平展拓扑的 $X$ 上方概形范畴，叠 $F$ 可以延拓成 $\mathbf{Sch}_X$ 上的叠 $F$，对每个态射 $f \colon X' \to X$，仍用 $F(X')$ 记该叠 $F$ 在 $X'$ 上方的笛卡儿截面范畴.

\emph{叠丛}是满足下列条件的叠\index[term]{diecong@叠丛}：对每个平展于 $X$ 上的概形 $X'$ 和纤维 $F_{X'}$ 的每对对象 $x$，$y$，$x$ 到 $y$ 的每个态射是同构，$x$ 和 $y$ 局部同构，且 $X_{\mathrm{ét}}$ 中使 $F_{X'}$ 非空的对象 $X'$ 的集是 $X_{\mathrm{ét}}$ 的加细. 例如群层下的挠子叠是叠丛，而且它有笛卡儿截面. 反之有截面的叠丛，i.e. 存在 $F_X$ 的对象 $x$，等价于自同构群层 $\mathrm{Aut}_X(x)$ 下的挠子叠.

有叠 $F$ 的子叠丛和\emph{极大子叠丛}的显然概念. 给定 $F(X)$ 的笛卡儿截面 $x$，存在唯一的 $F$ 的极大子叠丛\index[term]{jidazidiecong@极大子叠丛} $G_x$，使 $x$ 经 $G_x$ 分解. 称 $G_x$ 为 $x$ 生成的子叠丛；按定义它是\emph{平凡叠丛}\index[term]{pingfandiecong@平凡叠丛}. 由
\[
S_F(X') = \{F|_{X'}\index[notations]{$S_F$} \text{ 的极大子叠丛}\}
\]
定义的预层 $S_F$ 是层，称为 $F$ 的\emph{极大子叠丛层}\index[term]{jidazidiecongceng@极大子叠丛层}. 设 $\mathcal{O}$ 是由
\[
\mathcal{O}(X') = \{F_{X'} \text{ 的对象模同构的类}\}
\]
定义的预层. 把 $F_{X'}$ 的每个对象 $x$ 配给 $x$ 生成的 $F|_{X'}$ 的极大子叠丛，得到态射
\[
\mathcal{O} \longrightarrow S_F;
\]
由 \cite[III 2.1.4]{exp13-2}，该态射使 $S_F$ 成为与 $\mathcal{O}$ 关联的层.

称叠 $F$ 可构\index[term]{kegoudie@可构叠} (resp. ind-$L$-有限\index[term]{ind-L-youxiandie@ind-$L$-有限叠}，$L$ 是素数集)，若对每个平展于 $X$ 上的概形 $X'$ 和 $F_{X'}$ 的每个对象 $x$，层 $\mathrm{Aut}_{X'}(x)$ 如此 \cite[VII 2.2.1]{exp13-2}. 称叠 1-可构\index[term]{1-kegoudie@1-可构叠}，若它可构且极大子叠丛层可构.

\section{上同调真性}

\refstepcounter{theorem}\noindent\textbf{\thesection.\arabic{theorem}.}\quad 设 $S$ 是概形，$f \colon X \to Y$ 是 $S$-概形的态射. 若 $S'$ 是 $S$-概形，考虑下列所有方形笛卡儿的图式：
\begin{equation}
\begin{tikzcd}[row sep=large, column sep=large]
X \arrow[d,"{f}"] & X' \arrow[l,"h"] \arrow[d,"{f'}"] \\
Y \arrow[d] & Y' \arrow[l,"g"] \arrow[d] \\
S & S' \arrow[l]
\end{tikzcd}
\tag{1.0.1}
\label{eq:XIII.1.0.1}
\end{equation}
若 $Y_1$ 是 $Y$ 上方的平展概形，令 $X_1 = X \times_Y Y_1$，$Y'_1 = Y' \times_Y Y_1$，考虑笛卡儿方形
\begin{equation}
\begin{tikzcd}[row sep=large, column sep=large]
X_1 \arrow[d,"{f_1}"] & X'_1 \arrow[l,"{h_1}"] \arrow[d,"{f'_1}"] \\
Y_1 & Y'_1 \arrow[l,"{g_1}"]
\end{tikzcd}
\tag{1.0.2}
\label{eq:XIII.1.0.2}
\end{equation}

\begin{definition}\label{def:XIII.1.1}
设 $F$ 是 $X$ 上的叠. 称 $(F, f)$ 相对 $S$ 在维数 $\leqslant -1$\index[term]{shangtongdiaozhen@上同调真} (resp. $\leqslant 0$，resp. $\leqslant 1$) 时上同调真，若对每个 $S$-概形 $S'$，典范函子 (由叠的逆像的万有性质以显然方式定义)：
\[
g^{*} f_* F \longrightarrow f'_* h^{*} F \quad \text{(参见 1.0.1)}
\]
忠实 (resp. 忠实全，resp. 范畴等价).
\end{definition}

若在 $S$ 上无混淆可能，特别当 $S = Y$ 时，说上同调真而不再说相对 $S$ 上同调真.

\refstepcounter{theorem}\noindent\textbf{\thesection.\arabic{theorem}.}\quad 设 $F$ 是 $X$ 上的集合层；设 $\Phi$ 是与 $F$ 关联的离散范畴叠，i.e. 纤维在平展于 $X$ 上的每个概形 $X_1$ 上方是对象集为 $F(X_1)$ 的离散范畴的叠. 称 $(F, f)$ 相对 $S$ 在维数 $\leqslant -1$ (resp. $\leqslant 0$) 时上同调真，若 $(\Phi, f)$ 相对 $S$ 在维数 $\leqslant 0$ (resp. $\leqslant 1$) 时上同调真.

典范态射
\[
(1.2.1) \quad g^{*} f_* F \longrightarrow f'_* h^{*} F
\]
经离散范畴叠的过渡给出典范态射
\[
g^{*} f_* \Phi \longrightarrow f'_* h^{*} \Phi.
\]
从而说 $(F, f)$ 相对 $S$ 在维数 $\leqslant -1$ (resp. $\leqslant 0$) 时上同调真，等价于说对每个 $S$-概形 $S'$，态射 (1.2.1) 单射 (resp. 双射).

\refstepcounter{theorem}\noindent\textbf{\thesection.\arabic{theorem}.}\quad 设 $F$ 是 $X$ 上的群层，$\Phi$ 是 $X$ 上群 $F$ 的挠子叠 \cite[II 2.3.2]{exp13-1}. 称 $(F, f)$ 相对 $S$ 在维数 $\leqslant -1$ (resp. $\leqslant 0$，resp. $\leqslant 1$) 时上同调真，若 $(\Phi, f)$ 相对 $S$ 在维数 $\leqslant -1$ (resp. $\leqslant 0$，resp. $\leqslant 1$) 时上同调真. 上同调真性条件可以如下明确.

\begin{proposition}\label{prop:XIII.1.3.1}
记号同 \eqref{eq:XIII.1.0.1} 和 \eqref{eq:XIII.1.0.2}. 设 $F$ 是 $X$ 上的群层. 用 $F'$ (resp. $F_1$，resp. $F'_1$，等等) 记 $F$ 在 $X'$ (resp. $X_1$，resp. $X'_1$，等等) 上的逆像. 则下列条件等价：
\begin{enumerate}[i)]
\item $(F, f)$ 相对 $S$ 在维数 $\leqslant -1$ (resp. $\leqslant 0$，resp. $\leqslant 1$) 时上同调真.
\item 对每个态射 $S' \to S$、每个平展于 $Y$ 上方的概形 $Y_1$ 和 $X_1$ 上群 $F_1$ 的每个挠子 $P$，若 $P F_1$ 表示 $F_1$ 被 $P$ 扭转的群 \cite[II 4.1.2.3]{exp13-1}，典范态射
\[
a_0 \colon g_{1*}(f_{1*}(P F_1)) \longrightarrow f'_{1*}(P' F'_1)
\]
单射 (resp. $a_0$ 双射且典范态射
\[
a_1 \colon g^{*}(R^1 f_* F) \longrightarrow R^1 f'_* F'
\]
单射，resp. $a_0$ 和 $a_1$ 双射).
\item[(ii bis)] 对每个态射 $S' \to S$、每个平展于 $Y$ 上方的概形 $Y_1$、$X_1$ 上群 $F_1$ 的每个挠子 $P$ 和 $P F_1$ 下的每个挠子 $R$，典范态射
\[
\alpha_0 \colon g_{1*}(f_{1*} R) \longrightarrow f'_{1*} R'
\]
单射 (resp. $\alpha_0$ 双射，resp. 态射 $\alpha_0$ 和
\[
\alpha_1 \colon g_{1*}(R^1 f_{1*}(P F_1)) \longrightarrow R^1 f'_{1*}(P' F'_1)
\]
双射).
\end{enumerate}
\end{proposition}

\begin{proof}
(i)$\Rightarrow$(ii bis). 由 \cite[II 4.2.5]{exp13-1}，群 $P F_1$ 的每个挠子 $R$ 形如 $R = Q \wedge^{F_1} P^{\circ}$，其中 $Q$ 是群 $F_1$ 的挠子，$P^{\circ}$ 是 $P$ 的逆挠子. 有 $R' \simeq Q' \wedge^{F'_1} P'^{\circ}$. 设 $\Phi$ 是 $F$ 下的挠子叠，设 $x$，$y$ (resp. $x'$，$y'$) 是纤维范畴 $(g^{*} f_* \Phi)_{Y'_1}$ (resp. $(f'_* \Phi')_{Y'_1}$) 的对应 $P$，$Q$ (resp. $P'$，$Q'$) 的对象. 有关系
\[
Q \wedge^{F_1} P^{\circ} \simeq \mathrm{Hom}_{F_1}(P, Q),
\]
由此得出典范同构
\[
\mathrm{Hom}_{Y'_1}(x, y) \simeq g_{1*} f_{1*}(Q \wedge^{F_1} P^{\circ}), \quad \mathrm{Hom}_{Y'_1}(x', y') \simeq f'_{1*}(Q' \wedge^{F'_1} P'^{\circ}).
\]
从而态射 $\alpha_0$ 等同于态射
\[
\mathrm{Hom}_{Y'_1}(x, y) \longrightarrow \mathrm{Hom}_{Y'_1}(x', y'),
\]
由此得出：若 $(F, f)$ 相对 $S$ 在维数 $\leqslant -1$ 时上同调真，则 $\alpha_0$ 单射，若 $(F, f)$ 在维数 $\leqslant 0$ 时上同调真，则 $\alpha_0$ 双射.

现在假设 $(F, f)$ 相对 $S$ 在维数 $\leqslant 1$ 时上同调真，i.e. 典范态射
\[
\phi \colon g^{*} f_* \Phi \longrightarrow f'_* \Phi'
\]
是等价. 设 $G$ 是叠 $f_* \Phi$ 的极大子叠丛层 \cite[III 2.1.8]{exp13-1}；有同构 $G \simeq R^1 f_* F$. 由于 $g^{*} G$ 是 $g^{*} f_* \Phi$ 的极大子叠丛层 \cite[III 2.1.5.5]{exp13-2}，态射 $\alpha_1$ 由 $\phi|_{Y'_1}$ 取极大子叠丛层得到，故是同构.

(ii bis)$\Rightarrow$(ii). 只需证明 $\alpha_0$ 双射时 $a_1$ 单射. 设 $Y'_1$ 是平展于 $Y'$ 上方的概形，$s$ 和 $t$ 是 $g^{*}(R^1 f_* F)(Y'_1)$ 的在 $R^1 f'_* F'(Y'_1)$ 中有同样像的两个元素，证明 $s = t$. 该断言局部于 $Y'_1$ 的平展拓扑，考虑到 $g^{*}(R^1 f_* F)$ 的逆像定义，可假设 $Y'_1$ 是平展于 $Y$ 上方的概形 $Y_1$ 的逆像，且 $s$ 和 $t$ 来自 $X_1$ 上的挠子 $P$ 和 $Q$. 对 $s$ 和 $t$ 的假设于是意味着 $P$ 和 $Q$ 在 $X'_1$ 上的逆像 $P'$ 和 $Q'$ 在 $Y'_1$ 的平展拓扑下局部同构. 令 $R = \mathrm{Hom}_{F_1}(P, Q)$，态射
\[
g_{1*} f_{1*} R \longrightarrow f'_{1*} R'
\]
双射的事实证明 $P$ 和 $Q$ 在 $Y_1$ 的平展拓扑下局部同构，故 $s = t$.

(ii)$\Rightarrow$(i). 为证明 $\phi$ 忠实 (resp. 忠实全)，只需证明：若 $Y_1$ 是平展于 $Y$ 上的概形，$P$，$Q$ 是 $X_1$ 上群 $F_1$ 的两个挠子，$x$，$y$ (resp. $x'$，$y'$) 是 $(g^{*} f_* \Phi)_{Y'_1}$ (resp. $(f'_* \Phi')_{Y'_1}$) 的对应 $P$，$Q$ (resp. $P'$，$Q'$) 的对象，则态射
\[
a \colon \mathrm{Hom}(x, y) \longrightarrow \mathrm{Hom}(x', y')
\]
单射 (resp. 双射). $a$ 等同于典范态射
\[
H^0(Y'_1, g_{1*} f_{1*}(Q \wedge^{F_1} P^{\circ})) \longrightarrow H^0(Y'_1, f'_{1*}(Q' \wedge^{F'_1} P'^{\circ})).
\]
若 $\mathrm{Hom}(x, y) \neq \varnothing$，则 $Q \wedge^{F_1} P^{\circ}$ 是 $P F_1$ 下在 $Y_1$ 上局部平凡的挠子；由此得出 $f_{1*}(Q \wedge^{F_1} P^{\circ})$ 是 $f_{1*}(P F_1)$ 下的挠子，且 $g_{1*} f_{1*}(Q \wedge^{F_1} P^{\circ})$ 是平凡挠子. 态射 $a$ 于是等同于典范态射
\[
H^0(Y_1, g_{1*} f_{1*}(P F_1)) \longrightarrow H^0(Y'_1, f'_{1*}(P' F'_1)).
\]
若 $\mathrm{Hom}(x', y') \neq \varnothing$ 且 $a_1$ 单射也如此，因为此时 $Q' \wedge^{F'_1} P'^{\circ}$ 平凡，且由 $a_1$ 的单射得出 $P$ 和 $Q$ 在 $Y_1$ 上局部同构. 由此得出：$a_0$ 单射 (resp. $a_0$ 双射且 $a_1$ 单射) 时 $\phi$ 忠实 (resp. 忠实全).

剩下证明 $a_0$ 和 $a_1$ 双射时函子 $\phi$ 本质满射. 设 $Y''$ 是平展于 $Y'$ 上方的概形，$X'' = X' \times_{Y'} Y''$，设 $P''$ 是 $X''$ 上群 $F'' = F'|_{X''}$ 的挠子. 将证明存在 $(g^{*} f_* \Phi)_{Y''}$ 的元素 $x$，使它在 $(f'_* \Phi')_{Y''}$ 中的像同构于 $P''$. 设 $p''$ 是 $P''$ 的类. 由 $a_1$ 满射得出可以找到平展满射态射 $Y''_1 \to Y''$、平展态射 $Y_1 \to Y$ (使有态射 $Y''_1 \to Y'_1$) 和 $X_1$ 上群 $F_1$ 的挠子 $P_1$，使 $P_1$ 在 $X''_1$ 上的逆像 $P''_1$ 同构于 $P''$ 的逆像. 用 $\phi$ 忠实全的事实，看到 $(g^{*} f_* \Phi)_{Y''_1}$ 中对应 $P''_1$ 的对象 $x_1$ 配备相对 $Y''_1 \to Y''$ 的下降数据，故来自 $(g^{*} f_* \Phi)_{Y''}$ 的元素 $x$. 由于 $x$ 在 $(f'_* \Phi')_{Y''}$ 中的像是 $P''$，这证明 $\phi$ 本质满射，完成证明.
\end{proof}

\begin{example}\label{exa:XIII.1.4}
设 $f \colon X \to Y$ 是紧合态射. 由 \cite[VII 2.2.2]{exp13-2} 得出，对 $X$ 上每个 ind-有限叠 $F$，对 $(F, f)$ 相对 $Y$ 在维数 $\leqslant 1$ 时上同调真. 特别，对 $X$ 上每个集合层 (resp. 每个群层，resp. 每个 ind-有限群层) $F$，$(F, f)$ 在维数 $\leqslant 0$ (resp. $\leqslant 0$，resp. $\leqslant 1$) 时上同调真.
\end{example}

\begin{remark}\label{rem:XIII.1.5}
a) 设 $F$ 是 $X$ 上使 $(F, f)$ 相对 $S$ 在维数 $\leqslant -1$ (resp. $\leqslant 0$) 时上同调真的群层. 若把 $F$ 看作集合层，$(F, f)$ 相对 $S$ 在维数 $\leqslant -1$ (resp. $\leqslant 0$) 时上同调真，但反之不真.

例如取 $Y$ 为严格局部离散赋值环的谱，闭点为 $t$，生成点为 $s$，$f \colon X \to Y$ 为闭纤维为空的非空 $Y$-概形，$F$ 为 $X$ 上的非平凡常群层，$P$ 为使 $H^0(X_s, P F|_{X_s}) = 1$ 的 $F$ 下挠子. 则把 $P F$ 看作集合层时，$(P F, f)$ 相对 $Y$ 在维数 $\leqslant -1$ 时上同调真. 把 $P F$ 看作群层时，有同构 $(P F)^{\circ P} \simeq F$；由于典范态射
\[
H^0(X, F) \longrightarrow H^0(X_t, F|_{X_t}) = 1
\]
不单射，这证明 $(P F, f)$ 相对 $Y$ 在维数 $\leqslant -1$ 时不真.

b) 假设 $f$ 相容 (i.e. 拟紧且拟分离). 设 $F$ 是 $X$ 上的叠. 对 $Y'$ 的每个几何点 $y$，用 $Y$ (resp. $Y'$) 记 $Y$ (resp. $Y'$) 在 $y$ 处的严格局部化，令 $X = X \times_Y Y$，$X' = X' \times_{Y'} Y'$，等等. 为使 $(F, f)$ 相对 $S$ 在维数 $\leqslant -1$ (resp. $\leqslant 0$，resp. $\leqslant 1$) 时上同调真，必要且充分条件是：对每个 $S$-概形 $S'$ 和 $Y'$ 的每个几何点 $y$，典范函子
\[
F(X) \longrightarrow F'(X')
\]
忠实 (resp. 忠实全，resp. 范畴等价).

确实，若 $S'$ 是 $S$-概形，为使函子
\[
g^{*} f_* F \longrightarrow f'_* F'
\]
忠实 (resp. 忠实全，resp. 等价)，必要且充分条件是它诱导的 $Y'$ 的各几何点 $y'$ 上的纤维函子如此 \cite[III 2.1.5.9]{exp13-2}. 该断言于是由相容态射的叠的直像的几何纤维计算 \cite[VII 2.1.5]{exp13-2} 得出.

c) 设 $F$ 是 $X$ 上的叠. $(F, f)$ 相对 $S$ 在维数 $\leqslant -1$ (resp. $\leqslant 0$，resp. $\leqslant 1$) 时上同调真的事实局部于 $Y$ 的平展拓扑.

设 $S'$ 是 $S$-概形，$F'$ 是 $F$ 在 $X'$ 上的逆像 (参见 \eqref{eq:XIII.1.0.1}). 若 $(F, f)$ 相对 $S$ 在维数 $\leqslant 1$ 时上同调真，则 $(F', f')$ 也如此. 但 $(F, f)$ 相对 $S$ 在维数 $\leqslant -1$ (resp. $\leqslant 0$) 时上同调真不蕴含 $(F', f')$ 如此.

例如设 $S'$ 是离散赋值环，$f' \colon E_{S'} \to S'$ 是 $S'$ 上方的仿射空间，$x$ 是 $E_{S'}$ 的位于 $S'$ 的生成点上方的一个闭点，$F'$ 是 $E_{S'}$ 上的集合层，它在 $E_{S'} - \{x\}$ 上的限制是单元常层，在 $x$ 上方的几何点处的纤维有两个元素. 则 $(F', f')$ 相对 $S'$ 在维数 $\leqslant -1$ 时不真. 设 $S = S'[Z]$，$f \colon E_S \to S$ 是 $S$ 上的仿射空间，$T$ 是 $X = E_S$ 的不与闭集 $Z = 0$ 相交、且使 $f(T)$ 含 $S$ 的生成点的闭部分. 设 $G$ 是 $F'$ 在 $X$ 上的逆像，$F$ 是把 $G|_{X - T}$ 用空集延拓得到的 $X$ 上集合层. 则 $(F, f)$ 相对 $S$ 在维数 $\leqslant -1$ 时上同调真，但经由 $Z = 0$ 定义的基变换 $S' \to S$ 后不再如此.

d) 设 $F$ 是 $X$ 上使 $(F, f)$ 相对 $Y$ 在维数 $\leqslant -1$ (resp. $\leqslant 0$，resp. $\leqslant 1$) 时上同调真的叠. 则对 $Y$ 的每个几何点 $y$，典范函子
\[
(f_* F)_y \longrightarrow F(X_y)
\]
忠实 (resp. 忠实全，resp. 范畴等价).
\end{remark}

\begin{proposition}\label{prop:XIII.1.6}
设 $f \colon X \to Y$ 和 $g \colon Y \to Z$ 是两个 $S$-态射，$\Phi$ 是 $X$ 上的叠.
\begin{enumerate}[1)]
\item 假设 $(\Phi, f)$ 和 $(f_* \Phi, g)$ 相对 $S$ 在维数 $\leqslant -1$ (resp. $\leqslant 0$，resp. $\leqslant 1$) 时上同调真. 则 $(\Phi, gf)$ 也如此.
\item 假设 $(\Phi, gf)$ 相对 $S$ 在维数 $\leqslant -1$ 时\emph{上同调真} (resp. $(\Phi, gf)$ 相对 $S$ 在维数 $\leqslant 0$ 时上同调真且 $(\Phi, f)$ 相对 $S$ 在维数 $\leqslant -1$ 时上同调真，resp. $(\Phi, gf)$ 相对 $S$ 在维数 $\leqslant 1$ 时上同调真且 $(\Phi, f)$ 相对 $S$ 在维数 $\leqslant 0$ 时上同调真). 则 $(f_* \Phi, g)$ 相对 $S$ 在维数 $\leqslant -1$ (resp. $\leqslant 0$，resp. $\leqslant 1$) 时上同调真.
\end{enumerate}
\end{proposition}

对每个 $S$-概形 $S'$，考虑下列所有方形笛卡儿的图式：
\begin{equation}
\begin{tikzcd}[row sep=large, column sep=large]
X' \arrow[r,"{f'}"] \arrow[d,"{h}"] & Y' \arrow[r,"{g'}"] \arrow[d,"{k}"] & Z' \arrow[r] \arrow[d,"{m}"] & S' \arrow[d] \\
X \arrow[r,"{f}"] & Y \arrow[r,"{g}"] & Z \arrow[r] & S.
\end{tikzcd}
\tag{1.6.1}
\label{eq:XIII.1.6.1}
\end{equation}
典范态射
\[
m^{*}(g_* f_* \Phi) \longrightarrow g'_* f'_*(h^{*} \Phi)
\]
等同于典范态射
\[
i \quad j
\]
的复合：
\[
m^{*}(g_* f_* \Phi) \xrightarrow{i} g'_*(k^{*} f_* \Phi) \xrightarrow{j} g'_* f'_*(h^{*} \Phi).
\]
1) 假设蕴含 $i$ 和 $j$ 忠实 (resp. 忠实全，resp. 等价)；故 $ji$ 如此.

2) 假设蕴含 $ji$ 忠实 (resp. $ji$ 忠实全且 $j$ 忠实，resp. $ji$ 等价且 $j$ 忠实全)；由此得出 $i$ 忠实 (resp. 忠实全，resp. 等价).

\begin{corollary}\label{cor:XIII.1.7}
设 $f \colon X \to Y$ 和 $g \colon Y \to Z$ 是两个 $S$-态射，$F$ 是 $X$ 上的群层. 假设 $(F, gf)$ 相对 $S$ 在维数 $\leqslant -1$ 时上同调真 (resp. $(F, gf)$ 相对 $S$ 在维数 $\leqslant 0$ 时上同调真且 $(F, f)$ 相对 $S$ 在维数 $\leqslant -1$ 时上同调真). 则 $(f_* F, g)$ 相对 $S$ 在维数 $\leqslant -1$ (resp. $\leqslant 0$) 时上同调真.
\end{corollary}

沿用 \eqref{eq:XIII.1.6.1} 的记号，对每个平展于 $Y$ 上方的概形 $Y_1$，用 $f_1$，$F_1$ 记 $f$，$F$ 经态射 $Y_1 \to Y$ 的逆像. 设 $\Phi$ 是 $F$ 下的挠子叠，$\Psi$ 是 $f_* F$ 下的挠子叠. 有典范函子
\[
\phi \colon \Psi \longrightarrow f_* \Phi,
\]
它把平展于 $Y$ 上的每个概形 $Y_1$ 和 $Y_1$ 上群 $f_{1*} F_1$ 的每个挠子 $P$ 配给由 $f_1^{*} P$ 经结构群扩张 $f_1^{*} f_{1*} F_1 \to F_1$ 得到的 $X_1$ 上的挠子 $\widetilde{P}$. 函子 $\phi$ 忠实全. 确实，若 $P$ 和 $Q$ 是 $Y_1$ 上群 $f_{1*} F_1$ 的两个挠子，有典范态射
\[
\mathrm{Isom}_{f_{1*} F_1}(P, Q) \longrightarrow f_{1*}(\mathrm{Isom}_{F_1}(\widetilde{P}, \widetilde{Q}))
\]
它局部如此故是同构. 由此得出典范态射
\[
\mathrm{Isom}_{f_{1*} F_1}(P, Q) \longrightarrow \mathrm{Isom}_{F_1}(\widetilde{P}, \widetilde{Q})
\]
是同构，故 $\phi$ 忠实全.

有交换图式
\[
\begin{tikzcd}[row sep=large, column sep=large]
g'_* k^{*} \Psi \arrow[r,"{i}"] \arrow[d,"{g'_* k^{*}(\phi)}"] & m^{*}(g_* \Psi) \arrow[d,"{m^{*} g_*(\phi)}"] \\
g'_* k^{*}(f_* \Phi) \arrow[r,"{j}"] & m^{*}(g_* f_* \Phi),
\end{tikzcd}
\]
其中 $i$ 和 $j$ 是基变换态射. 由 \ref{prop:XIII.1.6} 2) 得出 $j$ 忠实 (resp. 忠实全). 由于 $g'_* k^{*}(\phi)$ 和 $m^{*} g_*(\phi)$ 忠实全，由上面图式得出 $i$ 忠实 (resp. 忠实全).

\begin{corollary}\label{cor:XIII.1.8}
设 $f \colon X \to Y$ 是相容 $S$-态射，$g \colon Y \to Z$ 是紧合 $S$-态射，$\Phi$ 是 $X$ 上的 ind-有限叠 \cite[VII 2.2.1]{exp13-2}. 假设 $(\Phi, f)$ 相对 $S$ 在维数 $\leqslant -1$ (resp. $\leqslant 0$，resp. $\leqslant 1$) 时上同调真，则 $(\Phi, gf)$ 也如此.
\end{corollary}

由于 $f$ 相容，$f_* \Phi$ 是 ind-有限\emph{叠} (SGA 4 IX 1.6 (ii)). 该推论于是由 \ref{prop:XIII.1.6} 1) 和 \ref{exa:XIII.1.4} 得出.

\begin{corollary}\label{cor:XIII.1.9}
设 $f \colon X \to Y$ 是整 $S$-态射，$g \colon Y \to Z$ 是 $S$-态射. 若 $F$ 是 $X$ 上的集合层，为使 $(f_* F, g)$ 相对 $S$ 在维数 $\leqslant -1$ (resp. $\leqslant 0$) 时上同调真，必要且充分条件是 $(F, gf)$ 如此. 若 $F$ 是 $X$ 上的群层，为使 $(f_* F, g)$ 相对 $S$ 在维数 $\leqslant -1$ (resp. $\leqslant 0$，resp. $\leqslant 1$) 时上同调真，必要且充分条件是 $(F, gf)$ 如此.
\end{corollary}

集合层情形的断言由 \ref{prop:XIII.1.6} 和 $(F, f)$ 相对 $S$ 在维数 $\leqslant 0$ 时上同调真的事实得出. 设 $F$ 是 $X$ 上的群层，$\Phi$ 是 $F$ 下的挠子叠. 由 SGA 4 VIII 5.8，$F$ 下的每个挠子局部平凡于 $Y$ 上. 由此得出叠 $f_* \Phi$ 等价于 $f_* F$ 下的挠子叠，等价由把平展于 $Y$ 上的每个概形 $Y_1$ 和 $X_1 = X \times_Y Y_1$ 上群 $F|_{X_1}$ 的每个挠子 $P$ 配给群 $f_* F|_{Y_1}$ 的挠子 $f_* P$ 得到. 由于 $(F, f)$ 相对 $S$ 在维数 $\leqslant 1$ 时上同调真，该推论于是由 \ref{prop:XIII.1.6} 得出.

\begin{definition}\label{def:XIII.1.10}
1.10.1. 设 $E$ 是范畴，考虑图式
\[
\Phi \xrightarrow{p} \Phi_1 \xrightarrow[p_2]{p_1} \Phi_2,
\]
其中 $\Phi$，$\Phi_1$，$\Phi_2$ 是 $E$ 上方的纤维范畴，箭是纤维范畴态射，设
\[
a \colon p_1 p \xrightarrow{\sim} p_2 p
\]
是函子同构.

称上图式正合，若满足条件 a)：对 $\Phi$ 的每对对象 $x$，$y$ 和满足 $p_1(f) = p_2(f)$ (借 $a$ 等同 $p_1 p$ 和 $p_2 p$) 的态射 $f \colon p(x) \to p(y)$，存在唯一态射 $g \colon x \to y$ 使 $p(g) = f$.

1.10.2. 考虑图式
\[
\Phi \xrightarrow{p} \Phi_1 \xrightarrow[p_1, p_2]{} \Phi_2 \xrightarrow[p_{12}, p_{23}, p_{31}]{} \Phi_3,
\]
其中 $\Phi$，$\Phi_i$ ($1 \leqslant i \leqslant 3$) 是 $E$ 上的纤维范畴，箭是纤维范畴态射. 假设给定函子同构
\[
a \colon p_1 p \xrightarrow{\sim} p_2 p
\]
\[
a_1 \colon p_{31} p_2 \xrightarrow{\sim} p_{12} p_1, \quad a_2 \colon p_{12} p_2 \xrightarrow{\sim} p_{23} p_1, \quad a_3 \colon p_{23} p_2 \xrightarrow{\sim} p_{31} p_1,
\]
使下列图式交换：
\[
\begin{tikzcd}[row sep=large, column sep=large]
p_{23} p_1 p \arrow[r,"{\mathrm{id}.a}"] \arrow[d,"{a_2 . \mathrm{id}}"] & p_{23} p_2 p \arrow[r,"{a_3 . \mathrm{id}}"] & p_{31} p_1 p \arrow[d,"{\mathrm{id}.a}"] \\
p_{12} p_2 p & p_{12} p_1 p \arrow[l,"{\mathrm{id}.a}"'] & p_{31} p_2 p \arrow[l,"{a_1 . \mathrm{id}}"']
\end{tikzcd}
\]
等同 $p_1 p$ 和 $p_2 p$，$p_{31} p_2$ 和 $p_{12} p_1$，等等.

称上图式正合，若满足条件：
a) 与 1.10.1 的条件 a) 类似.
b) 对 $\Phi_1$ 的每个对象 $x_1$ 和满足
\[
(1.10.2.1) \quad p_{23}(u) p_{31}(u) = p_{12}(u)^{-1}
\]
的同构 $u \colon p_1(x_1) \xrightarrow{\sim} p_2(x_1)$，存在 $\Phi$ 的对象 $x$ 和同构 $i \colon p(x) \xrightarrow{\sim} x_1$，使图式
\begin{equation}
\begin{tikzcd}[row sep=large, column sep=large]
p_1 p(x) \arrow[r,"{p_1(i)}"] \arrow[d] & p_1(x_1) \arrow[d,"{u}"] \\
p_2 p(x) \arrow[r,"{\sim}","{p_2(i)}"'] & p_2(x_1)
\end{tikzcd}
\tag{1.10.2.2}
\label{eq:XIII.1.10.2.2}
\end{equation}
交换.
\end{definition}

1.10.3. 以显然方式定义范畴 $E$ 上方纤维范畴的正合图式态射概念.

1.10.4. 特别将用 $E$ 是位点、$\Phi$，$\Phi_i$ ($1 \leqslant i \leqslant 3$) 是 $E$ 上的叠的情形中的正合图式概念.

设 $f \colon E \to E'$ 是位点态射，
\[
(1.10.4.1) \quad \Phi \xrightarrow{p} \Phi_1 \xrightarrow[p_2]{p_1} \Phi_2 \xrightarrow[p_{12}, p_{23}, p_{31}]{} \Phi_3
\]
是 $E$ 上叠的正合图式. 经直像得到显然正合的图式
\[
f_* \Phi \longrightarrow f_* \Phi_1 \rightrightarrows f_* \Phi_2 \Rrightarrow f_* \Phi_3.
\]
若 $h \colon E'' \to E$ 是位点态射，同样有正合图式
\[
h^{*} \Phi \xrightarrow{p''} h^{*} \Phi_1 \xrightarrow[p''_1, p''_2]{} h^{*} \Phi_2 \xrightarrow[p''_{12}, p''_{23}, p''_{31}]{} h^{*} \Phi_3.
\]
先验证 1.10.2 的条件 a). 设 $F''$ 是 $E''$ 的对象，$x''$ 和 $y''$ 是 $(h^{*} \Phi)_{F''}$ 的两个对象，$x''_1$ 和 $y''_1$ 是它们在 $h^{*} \Phi_1$ 中的像，$x''_2$，$y''_2$ 是它们在 $h^{*} \Phi_2$ 中的像. 设 $u''_1 \colon x''_1 \to y''_1$ 是满足 $p''_1(u''_1) = p''_2(u''_1)$ 的态射，证明 $u''_1$ 来自唯一态射 $u'' \colon x'' \to y''$. 问题局部于 $F''$，可假设有 $E$ 的对象 $F_1$，$F''$ 到 $F_1$ 经 $h$ 的逆像 $F''_1$ 的态射，$\Phi_{F_1}$ 的在 $F''$ 上的逆像是 $x''$ 和 $y''$ 的对象 $x$，$y$. 设 $x_1$，$y_1$ (resp. $x_2$，$y_2$) 是 $x$，$y$ 在 $\Phi_1$ (resp. $\Phi_2$) 中的像. 可假设 $u''_1$ 来自满足 $p_1(u_1) = p_2(u_1)$ 的态射 $u_1 \colon x_1 \to y_1$. 由 (1.10.4.1) 的正合性，得到唯一态射 $u \colon x \to y$，它的 $h$ 逆像是所求态射 $u''$.

1.10.2 的条件 b) 类似验证. 设 $x''_1$ 是 $(h^{*} \Phi_1)_{F''}$ 的对象，$u'' \colon p''_1(x''_1) \to p''_2(x''_1)$ 是满足关系
\[
p''_{23}(u'') p''_{31}(u'') = p''_{12}(u'')^{-1}
\]
的态射，证明存在 $(h^{*} \Phi)_{F''}$ 的对象 $x''$ 和同构 $i'' \colon p''(x'') \simeq x''_1$，使类似 \eqref{eq:XIII.1.10.2.2} 的图式交换. 由于问题局部于 $F''$，可假设有对象 $F_1$、如上 $F'' \to F''_1$ 的态射和 $(\Phi_1)_{F_1}$ 的在 $(h^{*} \Phi_1)_{F''}$ 中的逆像是 $x''_1$ 的对象 $x_1$. 同样可假设 $u''$ 来自满足 \eqref{eq:XIII.1.10.2.2} 的态射 $u \colon p_1(x_1) \to p_2(x_1)$. $\Phi_{F_1}$ 的对象 $x$ (它的 $h$ 逆像是回答问题的元素 $x''$) 的存在于是由 (1.10.4.1) 的正合性得出.

\begin{example}\label{exa:XIII.1.11}
1) 设 $f \colon X_1 \to X$ 是变动的概形上平展层范畴的下降态射 (例如\emph{万有浸没态射} (SGA 4 VIII 9.3)). 设 $X_2 = X_1 \times_X X_1$，$g \colon X_2 \to X$ 是典范投影，$F$ 是 $X$ 上的集合层. 由 loc. cit. 得出有集合层的正合序列
\[
(1.11.1) \quad F \longrightarrow f_* f^{*} F \rightrightarrows g_* g^{*} F.
\]
若 $\Phi$ 是与 $F$ 关联的离散范畴叠，$\Phi_3$ 是 $X$ 上的终叠，i.e. 所有纤维约化为一个元素 (唯一态射是恒等态射) 的叠，说序列 (1.11.1) 正合等价于说叠图式
\[
\Phi \longrightarrow f_* f^{*} \Phi \rightrightarrows g_* g^{*} \Phi \Rrightarrow \Phi_3
\]
如此.

2) 设 $f \colon X_1 \to X$ 是变动的概形上平展层范畴的有效下降态射 (例如紧合满射态射，或整满射态射，或局部有限表现忠实平坦态射 (SGA 4 VIII 9.4)). 设 $X_2 = X_1 \times_X X_1$，$g \colon X_2 \to X$ 是典范投影，$X_3 = X_1 \times_X X_1 \times_X X_1$，$h \colon X_3 \to X$ 是典范态射. 设 $\Phi$ 是 $X$ 上的叠，$\Phi_1 = f_* f^{*} \Phi$，$\Phi_2 = g_* g^{*} \Phi$，$\Phi_3 = h_* h^{*} \Phi$. 则有正合图式
\[
\Phi \longrightarrow \Phi_1 \rightrightarrows \Phi_2 \Rrightarrow \Phi_3,
\]
其中箭是投影关联的典范态射.

确实把 $\Phi$ 看作配备平展拓扑的 $X$ 上方概形范畴 $\mathbf{Sch}_X$ 上的叠. 由 \cite[VII 2.2.8]{exp13-2}，$\Phi$ 对 $\mathbf{Sch}_X$ 上使覆盖态射为平展层范畴的有效下降态射的最粗拓扑也是叠. 上图式的正合性由此立即得出.
\end{example}

\begin{proposition}\label{prop:XIII.1.12}
设 $S$ 是概形，$f \colon X \to Y$ 是 $S$-态射.
\begin{enumerate}[1)]
\item 设
\[
\Phi \longrightarrow \Phi_1 \rightrightarrows \Phi_2
\]
是 $X$ 上叠的正合图式. 假设 $(\Phi_1, f)$ 相对 $S$ 在维数 $\leqslant 0$ 时上同调真，$(\Phi_2, f)$ 相对 $S$ 在维数 $\leqslant -1$ 时上同调真. 则 $(\Phi, f)$ 相对 $S$ 在维数 $\leqslant 0$ 时上同调真.
\item 设
\[
\Phi \longrightarrow \Phi_1 \rightrightarrows \Phi_2 \Rrightarrow \Phi_3
\]
是 $X$ 上叠的正合图式. 假设 $(\Phi, f)$ 相对 $S$ 在维数 $\leqslant 1$ 时上同调真，$(\Phi_2, f)$ 相对 $S$ 在维数 $\leqslant 0$ 时上同调真，$(\Phi_3, f)$ 相对 $S$ 在维数 $\leqslant -1$ 时上同调真. 则 $(\Phi, f)$ 相对 $S$ 在维数 $\leqslant 1$ 时上同调真.
\end{enumerate}
\end{proposition}

对每个 $S$-概形 $S'$，考虑下列所有方形笛卡儿的交换图式：
\[
\begin{tikzcd}[row sep=large, column sep=large]
X' \arrow[r,"{f'}"] \arrow[d,"{h}"] & Y' \arrow[r] \arrow[d,"{g}"] & S' \arrow[d] \\
X \arrow[r,"{f}"] & Y \arrow[r] & S.
\end{tikzcd}
\]
证明 2)，1) 的证明类似. 由于直像和逆像函子把叠的正合图式变成正合图式 (1.10.4)，有叠的正合图式态射：
\[
\begin{tikzcd}[row sep=large, column sep=large]
g^{*} f_* \Phi \arrow[r,"{\pi}"] \arrow[d,"{\phi}"] & g^{*} f_* \Phi_1 \arrow[r,"{\pi_1}","{\pi_2}"'] \arrow[d,"{\phi_1}"] & g^{*} f_* \Phi_2 \arrow[r, Rightarrow] \arrow[d,"{\phi_2}"] & g^{*} f_* \Phi_3 \arrow[d,"{\phi_3}"] \\
f'_* h^{*} \Phi \arrow[r] & f'_* h^{*} \Phi_1 \arrow[r, Rightarrow] & f'_* h^{*} \Phi_2 \arrow[r, Rightarrow] & f'_* h^{*} \Phi_3.
\end{tikzcd}
\]
由假设 $\phi_1$ 是范畴等价，$\phi_2$ 忠实全，$\phi_3$ 忠实. 由前面图式得出 $\phi$ 是等价.

\begin{proposition}\label{prop:XIII.1.13}
设 $f \colon X \to Y$ 是 $S$-态射.
\begin{enumerate}[1)]
\item 设
\[
F \longrightarrow G \rightrightarrows H
\]
是 $X$ 上集合层的正合图式. 假设 $(G, f)$ 相对 $S$ 在维数 $\leqslant 0$ 时上同调真，$(H, f)$ 相对 $S$ 在维数 $\leqslant -1$ 时上同调真. 则 $(F, f)$ 相对 $S$ 在维数 $\leqslant 0$ 时上同调真.
\item 设 $F \to G$ 是 $X$ 上群层的单态射. 若 $Y_1$ 是平展于 $Y$ 上的概形，令 $X_1 = Y_1 \times_Y X$，用 $f_1$ (resp. $F_1$，resp. $G_1$) 记 $f$ (resp. $F$，resp. $G$) 在 $Y_1$ 上的逆像 (参见 1.0.2). 假设 $(G, f)$ 相对 $S$ 在维数 $\leqslant 0$ (resp. $\leqslant 1$) 时上同调真，且对每个平展于 $Y$ 上的概形 $Y_1$ 和 $G_1$ 下的每个挠子 $Q$，$(Q/F_1, f_1)$ 相对 $S$ 在维数 $\leqslant -1$ (resp. $\leqslant 0$) 时上同调真. 则 $(F, f)$ 相对 $S$ 在维数 $\leqslant 0$ (resp. $\leqslant 1$) 时上同调真.
\item 设 $F \to G$ 是 $X$ 上群层的单态射. 假设 $(F, f)$ 相对 $S$ 在维数 $\leqslant 1$ 时上同调真，$(G, f)$ 相对 $S$ 在维数 $\leqslant 0$ 时上同调真. 则对 $G$ 下的每个挠子 $Q$，$(Q/F, f)$ 相对 $S$ 在维数 $\leqslant 0$ 时上同调真.
\end{enumerate}
\end{proposition}

\begin{proof}
1) 设 $\Phi$ (resp. $\Phi_1$，resp. $\Phi_2$) 是与 $F$ (resp. $G$，resp. $H$) 关联的离散范畴叠，$\Phi_3$ 是 $X$ 上的终叠. 有正合图式
\[
\Phi \longrightarrow \Phi_1 \rightrightarrows \Phi_2 \Rrightarrow \Phi_3.
\]
由假设 $(\Phi_1, f)$ 相对 $S$ 在维数 $\leqslant 1$ 时上同调真，$(\Phi_2, f)$ 相对 $S$ 在维数 $\leqslant 0$ 时上同调真 (1.2). 由于 $(\Phi_3, f)$ 显然相对 $S$ 在维数 $\leqslant -1$ 时上同调真，由 \ref{prop:XIII.1.12} 得出 $(\Phi, f)$ 相对 $S$ 在维数 $\leqslant 1$ 时上同调真，i.e. $(F, f)$ 相对 $S$ 在维数 $\leqslant 0$ 时上同调真.

2) 先证明：若 $(G, f)$ 相对 $S$ 在维数 $\leqslant 0$ 时上同调真，且诸 $(Q/F_1, f_1)$ 相对 $S$ 在维数 $\leqslant -1$ 时上同调真，则 $(F, f)$ 相对 $S$ 在维数 $\leqslant 0$ 时上同调真. 由 \ref{prop:XIII.1.3.1} 只需证明：对每个平展于 $Y$ 上方的概形 $Y_1$ 和 $X_1$ 上群 $F_1$ 的每个挠子 $P$，把 $P F_1$ 看作集合层时 $(P F_1, f_1)$ 相对 $S$ 在维数 $\leqslant 0$ 时上同调真，且典范态射
\[
d \colon g^{*}(R^1 f_* F) \longrightarrow R^1 f'_* F'
\]
单射. 第一个断言由 1) 立即得出，因为若 $Q$ 表示由 $P$ 经结构群扩张 $F_1 \to G_1$ 得到的挠子，有同构
\[
Q^{G_1}/P F_1 \xrightarrow{\sim} Q/F_1.
\]
证明 $d$ 单射. 只需证明：若 $Y_1$ 是平展于 $Y$ 上方的概形，$P$ 和 $\widetilde{P}$ 是 $F_1$ 下的两个在 $X'_1$ 上的逆像 $P'$ 和 $\widetilde{P}'$ 同构的挠子，则 (必要时作 $Y_1$ 的平展满射扩张) $P$ 和 $\widetilde{P}$ 变成同构. 选同构 $p' \colon P' \xrightarrow{\sim} \widetilde{P}'$. 设 $Q$ (resp. $\widetilde{Q}$) 是由 $P$ (resp. $\widetilde{P}$) 经结构群扩张 $F_1 \to G_1$ 得到的挠子. $Q$ (resp. $\widetilde{Q}$) 在 $X'_1$ 上的逆像 $Q'$ (resp. $\widetilde{Q}'$) 由 $P'$ (resp. $\widetilde{P}'$) 经结构群扩张 $F'_1 \to G'_1$ 得到；设 $q' \colon Q' \xrightarrow{\sim} \widetilde{Q}'$ 是由 $p'$ 同样得到的同构. 由于 $(G, f)$ 相对 $S$ 在维数 $\leqslant 0$ 时上同调真，可假设 (必要时作 $Y_1$ 的平展满射扩张) $q'$ 是同构 $q \colon Q \xrightarrow{\sim} \widetilde{Q}$ 的像. 与挠子 $P$ (resp. $\widetilde{P}$) 关联 $Q/F_1$ 的截面 $x$ (resp. $\widetilde{Q}/\widetilde{F}_1$ 的截面 $\widetilde{x}$)，且为使 $P$ 和 $\widetilde{P}$ 同构，必要且充分条件是在同构 $Q \xrightarrow{\sim} \widetilde{Q}$ 使导出的同构
\[
e \colon H^0(X_1, Q/F_1) \longrightarrow H^0(X_1, \widetilde{Q}/\widetilde{F}_1)
\]
把 $x$ 变成 $\widetilde{x}$. 取同构 $q$. 截面 $e(x)$ 和 $\widetilde{x}$ 在 $H^0(X'_1, \widetilde{Q}'/\widetilde{F}'_1)$ 中有同样的像. 由于 $(\widetilde{Q}/\widetilde{F}_1, f_1)$ 相对 $S$ 在维数 $\leqslant -1$ 时上同调真，必要时作 $Y_1$ 的平展满射扩张，有 $e(x) = \widetilde{x}$，这证明 $d$ 单射.

为完成证明，剩下证明：若 $(G, f)$ 相对 $S$ 在维数 $\leqslant 1$ 时上同调真，且对每个平展于 $Y$ 上方的概形 $Y_1$ 和 $X_1$ 上群 $F_1$ 的每个挠子 $Q$，$(Q/F_1, f_1)$ 相对 $S$ 在维数 $\leqslant 0$ 时上同调真，则态射 $d$ 满射. 设 $P'$ 是 $X'_1$ 上群 $F'_1$ 的挠子，$Q'$ 是由 $P'$ 经结构群扩张得到的 $G'_1$ 下挠子. $P'$ 的给出等价于 $Q'$ 和 $H^0(X'_1, Q'/F'_1)$ 的截面 $x'$ 的给出. 由态射
\[
g^{*}(R^1 f_* G) \longrightarrow R^1 f'_* G'
\]
的满射得出，必要时作 $Y_1$ 的平展满射扩张，存在 $G_1$ 下的挠子 $Q$，使它在 $X'_1$ 上的逆像同构于 $Q'$. 用 $(Q/F_1, f_1)$ 相对 $S$ 在维数 $\leqslant 0$ 时上同调真的事实，同样可假设存在 $H^0(X_1, Q/F_1)$ 的元素 $x$，使它在 $H^0(X'_1, Q'/F'_1)$ 中的像是 $x'$. $Q$ 和 $x$ 的给出确定 $F_1$ 下的挠子 $P$，使它在 $X_1$ 上的逆像同构于 $P'$，这证明 $d$ 满射.

3) 证明 $(Q/F, f)$ 相对 $S$ 在维数 $\leqslant -1$ 时上同调真，i.e. 对每个 $S$-概形 $S'$ 和每个平展于 $Y$ 上方的概形 $Y_1$，若 $x$，$\widetilde{x}$ 是 $H^0(X_1, Q_1/F_1)$ 的在 $H^0(X'_1, Q'_1/F'_1)$ 中的像 $x'$，$\widetilde{x}'$ 相等的两个元素，则 (作 $Y_1$ 的满射扩张后) $x = \widetilde{x}$. 与 $x$ (resp. $\widetilde{x}$) 关联 $F_1$ 下的挠子 $P$ (resp. $\widetilde{P}$)，使 $Q_1$ 由 $P$ (resp. $\widetilde{P}$) 经结构群扩张 $F_1 \to G_1$ 得到. 由关系 $x' = \widetilde{x}'$ 得出同构 $u' \colon P' \xrightarrow{\sim} \widetilde{P}'$，使它在 $Q'_1$ 上诱导的同构是恒等. 由于 $(F, f)$ 相对 $S$ 在维数 $\leqslant 0$ 时上同调真，由此得出 (作 $Y_1$ 的平展满射扩张后) 有提升 $u'$ 的同构 $u \colon P \to \widetilde{P}$；$(G, f)$ 相对 $S$ 在维数 $\leqslant -1$ 时上同调真的事实蕴含 $x = \widetilde{x}$.

证明 $(Q/F, f)$ 相对 $S$ 在维数 $\leqslant 0$ 时上同调真. 设 $Y''$ 是平展于 $Y'$ 上的概形，$x''$ 是 $H^0(X'', Q''/F'')$ 的元素. 与 $x''$ 关联 $X''$ 上群 $F''$ 的挠子 $P''$. 由于 $(F, f)$ 相对 $S$ 在维数 $\leqslant 1$ 时上同调真，可以找到平展满射态射 $Y''_1 \to Y''$ 和 $Y_1 \to Y$，使有态射 $Y''_1 \to Y'_1$，以及 $X_1$ 上群 $F_1$ 的挠子 $P$，使 $P$ 在 $X''_1$ 上的逆像同构于 $P''$ 的逆像. 由 $(G, f)$ 相对 $S$ 在维数 $\leqslant 0$ 时上同调真的事实得出，甚至可以选 $Y''_1$ 和 $Y_1$，使由 $P$ 经结构群扩张 $F_1 \to G_1$ 得到的挠子同构于 $Q_1$；与 $P$ 对应 $H^0(X_1, Q_1/F_1)$ 的元素 $x$，使它在 $H^0(X''_1, Q''_1/F''_1)$ 中的像同构于 $x''$ 的逆像，这完成证明.
\end{proof}

\begin{proposition}\label{prop:XIII.1.14}
设 $f \colon X \to S$ 是 $S$-概形，$F$ 是 $X$ 上的集合层或群层 (resp. ind-$L$-群层，$L$ 是素数集). 假设 $F$ 局部常，$(F, f)$ 在维数 $\leqslant 0$ (resp. $\leqslant 1$) 时上同调真，且 $f$ 局部 0-无环 (resp. 对 $L$ 局部 1-非球) (SGA 4 XV 1.11). 则对 $S$ 的几何点的每个特殊化 $s_1 \to s_2$，特殊化态射 (SGA 4 VIII 7.1)
\[
a_0 \colon (f_* F)_{s_2} \longrightarrow (f_* F)_{s_1}
\]
是同构，且若 $F$ 是群层，态射
\[
a_1 \colon (R^1 f_* F)_{s_2} \longrightarrow (R^1 f_* F)_{s_1}
\]
单射 (resp. 态射 $a_0$ 和 $a_1$ 是同构).
\end{proposition}

证明照抄 SGA 4 XVI 2.3 的证明，把"紧合"换成"上同调真"即可.

\begin{corollary}\label{cor:XIII.1.15}
设 $f \colon X \to S$ 是态射，$\Phi$ 是 $X$ 上的叠，$L$ 是素数集. 假设对每个平展于 $X$ 上的概形 $X_1$ 和 $\Phi_{X_1}$ 的每对对象 $x$，$y$，层 $\mathrm{Hom}_{X_1}(x, y)$ 局部常，层 $\mathrm{Aut}_{X_1}(x)$ 是局部常 ind-$L$-群，$\Phi$ 的极大子叠丛层 $S_\Phi$ \cite[III 2.1.7]{exp13-1} 局部常. 假设 $(\Phi, f)$ 在维数 $\leqslant 1$ 时上同调真，$f$ 对 $L$ 局部 1-非球. 则对 $S$ 的几何点的每个特殊化 $s_1 \to s_2$，特殊化态射
\[
a \colon (f_* \Phi)_{s_2} \longrightarrow (f_* \Phi)_{s_1}
\]
是范畴等价.
\end{corollary}

设 $S_1$ (resp. $S_2$) 是 $S$ 在 $s_1$ (resp. $s_2$) 处的严格局部化，$X_2$，$\Phi_2$ (resp. $X_1$，$\Phi_1$) 是 $X_2$，$\Phi_2$ 在 $S_2$ 上 (resp. $X_1$，$\Phi_1$ 在 $S_1$ 上) 的逆像，考虑笛卡儿方形
\[
\begin{tikzcd}[row sep=large, column sep=large]
X_1 \arrow[d,"{f_1}"] & X_2 \arrow[l,"h"] \arrow[d,"{f_2}"] \\
S_1 & S_2 \arrow[l,"g"]
\end{tikzcd}
\]
需证明函子
\[
\phi \colon \Phi_2(X_2) \longrightarrow \Phi_1(X_1)
\]
是等价. 函子 $\phi$ 忠实全；设 $x$，$y$ 是 $(\Phi_2)_{X_2}$ 的两个对象；典范态射
\[
\mathrm{Hom}_{X_2}(x, y) \longrightarrow \mathrm{Hom}_{X_1}(\phi(x), \phi(y))
\]
等同于典范态射
\[
H^0(X_2, \mathrm{Hom}_{X_2}(x, y)) \longrightarrow H^0(X_1, h^{*}(\mathrm{Hom}_{X_2}(x, y))).
\]
由 \ref{prop:XIII.1.14} 该态射是同构.

证明 $\phi$ 是等价. 设 $x_1$ 是 $\Phi_1(X_1)$ 的对象，$G_1$ 是 $x_1$ 生成的 $\Phi_1$ 的极大子叠丛. 态射
\[
H^0(X_2, S_{\Phi_2}) \longrightarrow H^0(X_1, h^{*}(S_{\Phi_2})) = H^0(X_1, S_{\Phi_1})
\]
双射，故存在 $\Phi_2$ 的极大子叠丛 $G_2$，使 $h^{*} G_2$ 同构于 $G_1$. 于是只需证明函子
\[
G_2 \longrightarrow h_* h^{*} G_2
\]
是等价. 但这种形式的问题局部于 $X_2$ 的平展拓扑，可假设 $G_2$ 是 $G_2$ 的某个对象的自同构群下的挠子叠丛，此时断言由 \ref{prop:XIII.1.14} 得出.

\begin{corollary}\label{cor:XIII.1.16}
假设同 1.14. 再假设 $F$ 是集合层 (resp. ind-$L$-群层)，且 $f_* F$ (resp. $R^1 f_* F$) 可构，则 $f_* F$ (resp. $R^1 f_* F$) 局部常.
\end{corollary}

该推论由 \ref{prop:XIII.1.14} 借助 SGA 4 IX 2.11 得出.

\begin{remark}\label{rem:XIII.1.17}
回忆 $f$ 局部 0-无环的条件在 $f$ 平坦、纤维可分、$X$ 和 $Y$ 局部诺特时满足 (SGA 4 XV 4.1)，$f$ 对 $L$ 局部 1-非球的条件在 $f$ 光滑、$L$ 为与 $S$ 的剩余特征不同的素数集时满足 (SGA 4 XV 2.1).
\end{remark}

\section{上同调真性的特殊情形：相对正规交叉除子}

2.0. 设 $R$ 是分式域 $K$ 的离散赋值环，$L$ 是 $K$-平展代数；则 $L$ 是有限个域 $L_i$ 的直积，其中 $L_i$ 是 $K$ 的平展扩张. 若 $L'_i$ 表示 $L_i$ 在 $L_i$ 的代数闭包中生成的 伽罗瓦扩张，称 $L$ 在 $R$ 上温和分歧，若诸 $L'_i$ 是 X 3 意义下的温和分歧扩张，i.e. $L'_i|K$ 的某个惯性群 $I_i$ 的阶与 $R$ 的剩余特征 $p$ 互素.

已知 $I_i$ 在一切情形都是与 $p$ 互素的循环群被 $p$-群扩张得到的. (当 $R$ 的剩余扩张可分时，这由 \cite[第四章 命题 7 推论 4]{exp13-5} 得出. loc. cit. 中给出的证明按下列方式推广到一般情形. 沿用 loc. cit. 的假设和记号但不假设剩余扩张可分. 设 $H_i$ 是惯性群 $G_0$ 的由使 $A_L$ 的每个幺模化元 $\pi$ 满足 $s\pi/\pi \in U^i$ 的元素 $s$ 组成的子群. 验证 $G_0/H_1$ 是与 $p$ 互素的阶群，且对 $i \geqslant 1$，$H^i/H^{i+1}$ 是 $p$-群，由此得到所述结果.)

2.0.1. 当 $R$ 严格局部时，有下列简单刻画：$K$-代数 $L$ 在 $R$ 上温和分歧，当且仅当诸 $[L_i : K]$ 与 $p$ 互素. 而且，若 $L$ 在 $R$ 上温和分歧，诸 $L_i$ 是 $K$ 的循环扩张. 确实，当 $R$ 严格局部时，$I_i$ 等于 $L'_i$ 在 $K$ 上的 伽罗瓦群. 正如刚才所回忆，$I_i$ 是与 $p$ 互素的循环群被 $p$-群扩张得到的. 若假设 $L'_i$ 在 $R$ 上温和分歧，$I_i$ 是与 $p$ 互素的阶循环群. 由此得出 $[L_i : K]$ 与 $p$ 互素且 $L_i = L'_i$. 反之，若 $[L_i : K]$ 与 $p$ 互素，$I_i$ 不能含非平凡正规 $p$-子群；故 $I_i$ 是与 $p$ 互素的阶循环群，这证明 $L_i$ 在 $R$ 上温和分歧.

2.0.2. 设 $R$ 是分式域 $K$ 的离散赋值环，$L$ 是 $K$-平展代数，$R$ 是 $R$ 的严格局部化，$K$ 是它的分式域，$L = L \otimes_K K$. 则 $L$ 在 $R$ 上温和分歧\index[term]{wenhefenqidaishu@温和分歧代数}当且仅当 $L$ 在 $R$ 上温和分歧. 归结到 $L$ 是域的情形. 设 $L = \prod_i L_i$，其中诸 $L_i$ 是 $K$ 的扩张域；若 $L'$ 是 $L$ 生成的 伽罗瓦扩张，且 $L' = L' \otimes_K K$，同样有 $L'$ 分解成域之积 $L' = \prod_j L'_j$，且每个 $L_i$ 是某个 $L'_j$ 的子扩张. 由于 $L'$ 是 $K$ 的 伽罗瓦扩张，诸 $L'_j$ 是 $K$ 的 伽罗瓦扩张. 假设 $L$ 在 $R$ 上温和分歧；由于 $L'_j|K$ 的 伽罗瓦群同构于 $L'|K$ 的惯性群，诸 $L'_j$ 在 $R$ 上温和分歧，故诸 $L_i$ 也如此. 反之，假设 $L$ 在 $R$ 上温和分歧. 对每个 $j$，设 $v_j$ 是 $L'_j$ 的延拓 $K$ 的赋值的离散赋值，仍用 $v_j$ 记它在 $L'$ 上诱导的赋值. $j$ 变动时，$v_j$ 遍历 $L'$ 的延拓 $K$ 的赋值的诸赋值. 设 $G = \mathrm{Gal}(L'|K)$，$H = \mathrm{Gal}(L'|L)$，$I_j$ 是 $L'|K$ 在 $v_j$ 处的惯性群，$J_j$ 是 $L'|L$ 在 $v_j$ 处的惯性群. 群 $I_j$ 是与 $p$ 互素的循环群被 $p$-群 $P_j$ 扩张得到的. 由于诸 $L_i$ 在 $R$ 上温和分歧，$I_j/J_j$ 与 $p$ 互素，故有 $P_j \subset J_j$. 从而 $H$ 含所有 $P_j$，故也含 $j$ 变动时诸 $P_j$ 生成的群 $P$. 但 $P$ 在 $G$ 中不变，因为 $G$ 的内自同构把诸 $I_j$ 彼此变换，故也把诸 $P_j$ 彼此变换. 由此得出 $P$ 是 $H$ 的在 $G$ 中正规的子群，故由于 $L'$ 是 $L$ 生成的 伽罗瓦扩张，有 $P = 1$，这证明 $L$ 在 $R$ 上温和分歧.

更一般地设 $R \to R'$ 是离散赋值环的态射，使 $R$ 的幺模化元 $\pi$ 的像是 $R'$ 的幺模化元 $\pi'$，且剩余扩张 $k(R')$ 是 $k(R)$ 的可分扩张. 设 $K$ 是 $R$ 的分式域，$K'$ 是 $R'$ 的分式域，$L$ 是 $K$-平展代数，$L' = L \otimes_K K'$. 则 $L$ 在 $R$ 上温和分歧当且仅当 $L'$ 在 $R'$ 上温和分歧. 可假设 $R$ 和 $R'$ 严格局部. 由 2.0.1 只需证明 $L$ 是域时 $L'$ 也是. 设 $\widetilde{R}$ 是 $R$ 在 $L$ 中的正规化，$\widetilde{\pi}$ 是 $\widetilde{R}$ 的幺模化元，$\widetilde{R}' = \widetilde{R} \otimes_R R'$. 由于扩张 $k(\widetilde{R})|k(R)$ 根式，且扩张 $k(R')|k(R)$ 平展，$k(\widetilde{R}) \otimes_{k(R)} k(R')$ 是域 [EGA IV 4.3.2 和 4.3.5]. 这证明 $\widetilde{R}'$ 是局部环，且由于 $\pi$ 在 $R'$ 中的像是 $\pi'$，有 $k(\widetilde{R}') = \widetilde{R}'/(\widetilde{\pi})$；从而 $\widetilde{R}'$ 是离散赋值环 \cite[第一章 §2 命题 2]{exp13-5}，故 $L'$ 是域.

2.0.3. 经归结到严格局部情形，看到\emph{温和分歧代数}的子代数和代数，两个温和分歧代数的张量积温和分歧，温和分歧代数经离散赋值环扩张后仍温和分歧，经温和分歧扩张后变成温和分歧的代数和代数温和分歧.

\refstepcounter{theorem}\noindent\textbf{\thesection.\arabic{theorem}.}\quad 设 $X$ 是 $S$-概形，$D$ 是 $X$ 上 $\geqslant 0$ 的除子. 回忆 (SGA 5 II 4.2) 称 $D$ 相对 $S$ 严格正规交叉\index[term]{yangezhengguijiaochachuzi@严格正规交叉除子}，若存在 $\Gamma(X, \mathcal{O}_X)$ 的元素的有限族 $(f_i)_{i \in I}$，使 $D = \sum_{i \in I} \mathrm{div}(f_i)$，且满足下列条件：

2.1.0. 对 $\mathrm{Supp}\, D$ 的每个点 $x$，$X$ 在 $x$ 处平展于 $S$ 上，且若用 $I(x)$ 记使 $f_i(x) = 0$ 的 $i \in I$ 的集，子概形 $V((f_i)_{i \in I(x)})$ 在 $S$ 上光滑且在 $X$ 中余维数 $\mathrm{card}\, I(x)$.

称除子 $D$ 相对 $S$ 正规交叉\index[term]{zhengguijiaochachuzi@正规交叉除子}，若它在 $X$ 上对平展拓扑局部严格正规交叉.

设 $D$ 是相对 $S$ 正规交叉除子. 令 $Y = \mathrm{Supp}\, D$，$U = X - Y$，用 $i \colon U \to X$ 记典范浸入. 对 $S$ 的每个几何点 $s$ 和几何纤维 $Y_s$ 的每个极大点 $y$，环 $R = \mathcal{O}_{X_s, y}$ 是离散赋值环.

在本号余下部分，将用下列技术定义：

\begingroup
\addtocounter{theorem}{-1}
\renewcommand{\thetheorem}{\thesection.1.1}
\begin{definition}\label{def:XIII.2.1.1}
设 $F$ 是 $U$ 上的集合层. 称 $F$ 在 $X$ 上 (沿 $D$) 相对 $S$ 温和分歧\index[term]{wenhefenqiceng@温和分歧层}，若对 $S$ 的每个几何点 $s$ 满足下列条件：
对 $Y_s$ 的每个极大点 $y$，$F$ 在 $\mathcal{O}_{X_s}$ 的分式域 $K$ 上的限制可由 $K$-平展代数 $L$ 的谱表示，后者在 $\mathcal{O}_{X_s, y}$ 上温和分歧.
\end{definition}
\endgroup

最常无混淆可能时，在术语中省略 $D$ 的提及.

\begingroup
\addtocounter{theorem}{-1}
\renewcommand{\thetheorem}{\thesection.1.2}
\begin{definition}\label{def:XIII.2.1.2}
若 $F$ 是 $U$ 上在 $X$ 上相对 $S$ 温和分歧的群层，用
\[
H^1_t(U, F)
\]
记 $H^1(U, F)$ 的由在 $X$ 上相对 $S$ 温和分歧的 $F$ 下挠子的类组成的子集\index[notations]{$H_t^1(U, F)$}.
\end{definition}
\endgroup

设
\[
\begin{tikzcd}[row sep=large, column sep=large]
U \arrow[r,"i"] \arrow[dr,"g"] & X \arrow[d,"f"] \\
& T
\end{tikzcd}
\]
是 $S$-概形的交换图式，$i$ 如 2.1 中；用
\[
R^1_t g_* F
\]
记与预层 $T' \mapsto H^1_t(U', F)$ 关联的\index[notations]{$R_t^1 g_* F$} $T$ 上层，其中 $T'$ 遍历 $T$ 上方的平展概形，$U' = U \times_T T'$；$R^1_t g_* F$ 是 $R^1 g_* F$ 的子层.

注意 $g$ 相容、$t$ 是 $T$ 的几何点、$T$ 是 $T$ 在 $t$ 处的严格局部化、$U = U \times_T T$ 时，有同构
\[
(2.1.2.1) \quad (R^1_t g_* F)_t \simeq H^1_t(U, F)
\]

2.1.3. 设 $\mathcal{C}_t((U, X)/S)$ 或简记 $\mathcal{C}_t$\index[notations]{$\mathcal{C}_t$} 是 $U$ 的在 $X$ 上相对 $S$ 温和分歧的平展覆盖范畴. 假设 $U$ 连通，$a$ 是 $U$ 的几何点. 设 $\Gamma_t$ 是把 $U$ 的在 $X$ 上相对 $S$ 温和分歧的平展覆盖 $U'$ 配给 $U'$ 的位于 $a$ 上方的几何点集的函子. 由 2.0 得出对 $(\mathcal{C}_t, \Gamma_t)$ 满足 V 4 的公理 $(G_1)$ 至 $(G_6)$. 从而 $\Gamma_t$ 可被 pro-对象表示，称之为 $(U, X)$ 在 $a$ 处带点、相对 $S$ 的\emph{温和分歧万有覆盖}\index[term]{wenhefenqiwanyoufugai@温和分歧万有覆盖}. 温和分歧万有覆盖的 $U$-自同构群的逆群称为温和分歧基本群\index[term]{wenhefenqijibenqun@温和分歧基本群}，记为
\[
\pi^t_1((U, X)/S, a)\index[notations]{$\pi_1^t((U, X)/S, a)$} \quad \text{或简记} \quad \pi^t_1(U, a) \quad \text{甚至} \quad \pi^t_1(U).
\]
它显然是基本群 $\pi_1(U, a)$ 的商 (\cref{prop:V.6.9}).

2.1.4. 设 $F$ 是 $U$ 上的群层，$P$ 是群 $F$ 的右挠子，$Q$ 是群 $F$ 的左挠子，假设 $P$ 和 $Q$ 在 $X$ 上相对 $S$ 温和分歧. 则 $P \wedge^F Q$ 也如此. 归结到证明：若 $R$ 是分式域 $K$ 的离散赋值环，$F$ 是 $K$ 上有限平展群概形，$P$ 和 $Q$ 是在 $R$ 上温和分歧的 $F$ 下两个挠子，则 $P \wedge^F Q$ 也如此. 而 $T = P \wedge^F Q$ 是 $P \times_K Q$ 的商. 若 $L$，$M$，$N$ 分别表示表示 $T$，$P$，$Q$ 的 $K$-代数，则 $L$ 是 $M \otimes_K N$ 的子代数，由 2.0.3 得出 $L$ 在 $R$ 上温和分歧.

由前面得出：若 $F$ 是 $U$ 上的群层，且存在 $F$ 下的在 $X$ 上相对 $S$ 温和分歧的挠子，则 $F$ 在 $X$ 上相对 $S$ 温和分歧. 确实 $P$ 的逆挠子 $P^{\circ}$ 在 $X$ 上相对 $S$ 温和分歧，因为作为集合层它同构于 $P$. 若 $P F$ 是 $F$ 被 $P$ 扭转的群，有同构
\[
F \simeq P^{\circ} \wedge^F P,
\]
从而 $F$ 在 $X$ 上相对 $S$ 温和分歧.

如前看到，若 $F \to F'$ 是 $U$ 上两个在 $X$ 上相对 $S$ 温和分歧的群层的态射，$P$ 是在 $X$ 上相对 $S$ 温和分歧的 $F$ 下挠子，则由 $P$ 经结构群扩张 $F \to F'$ 得到的挠子 $P'$ 在 $X$ 上相对 $S$ 温和分歧.

特别，典范态射
\[
H^1(U, F) \longrightarrow H^1(U, F')
\]
经限制到 $H^1_t(U, F)$ 给出典范态射
\[
H^1_t(U, F) \longrightarrow H^1_t(U, F')
\]

2.1.5. 设 $S' \to S$ 是态射，用 $U'$ (resp. $X'$，等等) 记 $U$ (resp. $X$，等等) 在 $S'$ 上的逆像. 若 $F$ 是 $U$ 上在 $X$ 上相对 $S$ 温和分歧的集合层，由\cref{def:XIII.2.1.1} 和 2.0.3 得出 $F'$ 在 $X'$ 上相对 $S'$ 温和分歧.

若 $F$ 是 $U$ 上的群层，$F$ 下在 $X$ 上相对 $S$ 温和分歧的挠子在 $S'$ 上的逆像是 $F'$ 下在 $X'$ 上相对 $S'$ 温和分歧的挠子. 特别有典范函子
\[
(2.1.5.1) \quad \mathcal{C}_t((U, X)/S) \longrightarrow \mathcal{C}_t((U', X')/S').
\]
假设 $U$ 和 $U'$ 连通，$a$ 是 $U$ 的几何点，$a'$ 是 $U'$ 的位于 $a$ 上方的几何点；由前面得出典范态射
\[
(2.1.5.2) \quad \pi^t_1(U', a') \longrightarrow \pi^t_1(U, a).
\]
若 $S' \to S$ 是态射，$h \colon T' \to T$ 是典范投影，态射
\[
h^{*}(R^1 g_* F) \longrightarrow R^1 g'_* F'
\]
经限制给出典范态射
\[
(2.1.5.3) \quad h^{*}(R^1_t g_* F) \longrightarrow R^1_t g'_* F'
\]

2.1.6. 设 $F$ 是 $U$ 上在 $X$ 上相对 $S$ 温和分歧的群层. 记号同 \ref{def:XIII.2.1.2}，有典范正合序列：
\[
(2.1.6.1) \quad 1 \longrightarrow H^1(X, i_* F) \longrightarrow H^1_t(U, F) \longrightarrow H^0(X, R^1_t i_* F)
\]
\[
1 \longrightarrow R^1 f_*(i_* F) \longrightarrow R^1_t g_* F \longrightarrow f_*(R^1_t i_* F).
\]
第一个由正合序列 (SGA 4 III 3.2) 得出：
\[
1 \longrightarrow H^1(X, i_* F) \longrightarrow H^1(U, F) \longrightarrow H^0(X, R^1 i_* F).
\]
只需证明 $H^1(X, i_* F)$ 在 $H^1(U, F)$ 中的像实际包含在 $H^1_t(U, F)$ 中，且 $H^1_t(U, F)$ 在 $H^0(X, R^1 i_* F)$ 中的像实际包含在 $H^0(X, R^1_t i_* F)$ 中. $i_* F$ 下的挠子在 $U$ 上的逆像是 $i^{*} i_* F$ 下的挠子，后者显然在 $X$ 上相对 $S$ 温和分歧，故经结构群扩张 $i^{*} i_* F \to F$ 后仍如此，这证明箭 $H^1(X, i_* F) \to H^1_t(U, F)$ 存在. $H^1_t(U, F)$ 在 $H^0(X, R^1 i_* F)$ 中的像包含在 $H^0(X, R^1_t i_* F)$ 中的事实由 $R^1_t i_* F$ 的定义立即得出. 这证明第一个正合序列存在，第二个由第一个经局部化得出.

\refstepcounter{theorem}\noindent\textbf{\thesection.\arabic{theorem}.}\quad 沿用 2.1 的记号. 当 $U$ 上的叠 $\Phi$ 局部于 $X$ 上和 $S$ 上 (对平展拓扑) 是 $S$ 上叠 $\Psi$ 的逆像时，定义它的温和分歧对象概念.

先设 $G$ 是 $U$ 上的叠丛，假设给定平展满射态射 $S_1 \to S$、平展满射态射 $X_2 \to X \times_S S_1$、$S_1$ 上的平凡叠丛 $H$ 和同构
\[
G|_{U_2} \longrightarrow H|_{U_2},
\]
其中 $U_1 = U \times_X X_1$，$U_2 = U \times_X X_2$. 选 $H|_{X_2}$ 的平凡化时，上面同构把 $G|_{U_2}$ 等同于群层 $F$ 下的挠子叠. 称 $G_U$ 的元素 $x$ 在 $X$ 上相对 $S$ 温和分歧，若 $x$ 在 $U_2$ 上的限制是在 $X$ 上相对 $S$ 温和分歧的挠子. 由 2.1.4 这个概念不依赖 $H|_{X_2}$ 的平凡化方式.

现在设 $\Phi$ 是 $U$ 上的叠，假设给定平展满射态射 $S_1 \to S$、平展满射态射 $X_2 \to X \times_S S_1$、$S_1$ 上的叠 $\Psi$ 和同构
\[
i \colon \Phi|_{U_2} \longrightarrow \Psi|_{U_2}.
\]
设 $x$ 是 $\Phi_U$ 的元素，$G_x$ 是 $x$ 生成的 $\Phi$ 的极大子叠丛 \cite[III 2.1.7]{exp13-1}，$S_\Phi$ 是 $\Phi$ 的极大子叠丛层. 同构 $i$ 诱导同构
\[
S_\Phi|_{U_2} \longrightarrow S_\Psi|_{U_2}.
\]
由 \ref{cor:XIII.5.7} 得出，必要时把 $S_1$ 换成平展满射扩张，有 $\Psi$ 的唯一极大子叠丛 $H$ (可假设平凡)，使 $i$ 定义同构
\[
G_x|_{U_2} \longrightarrow H|_{U_2}.
\]
称元素 $x$ 在 $X$ 上相对 $S$ 温和分歧\index[term]{diedewenhefenqiyuansu@叠的温和分歧元素}，若它作为配备上述同构的 $G_x$ 的元素如此.

2.2.1. 设 $\Phi$ 是 $U$ 上的叠，局部于 $X$ 上和 $S$ 上是 $S$ 上叠的逆像，设
\[
\begin{tikzcd}[row sep=large, column sep=large]
U \arrow[r,"i"] \arrow[dr,"g"] & X \arrow[d,"f"] \\
& T
\end{tikzcd}
\]
是如 \ref{def:XIII.2.1.2} 的图式. 对每个平展于 $T$ 上的概形 $T'$，若 $U' = U \times_T T'$，考虑 $(g_* \Phi)_{T'} = \Phi_{U'}$ 的由 $\Phi_{U'}$ 的在 $X$ 上相对 $S$ 温和分歧的元素组成的子集 $(g_{*t} \Phi)_{T'}$\index[notations]{$(g_{*t}\Phi)_{T'}$}. 称 $\Phi$ 经 $g$ 的温和分歧直像\index[term]{diedewenhefenqizhixiang@叠的温和分歧直像}，记为
\[
g_{*t} \Phi
\]
是 $g_* \Phi$ 的由平展于 $T$ 上方的概形 $T'$ 上方的对象为 $(g_{*t} \Phi)_{T'}$ 的元素组成的全子范畴. $g_{*t} \Phi$ 显然是 $g_* \Phi$ 的子叠.

2.2.2. 经归结到挠子叠的情形，看到若 $\Phi$ 是 $U$ 上局部于 $S$ 和 $X$ 的平展拓扑是 $S$ 上叠的逆像的叠，典范态射
\[
h^{*}(g_* \Phi) \longrightarrow g'_* \Phi'
\]
经限制给出典范态射
\[
h^{*}(g_{*t} \Phi) \longrightarrow g'_{*t} \Phi'
\]

\begin{remark}\label{rem:XIII.2.3}
a) 若 $F$ 是 $U$ 上局部常可构集合层，为使 $F$ 在 $X$ 上相对 $S$ 温和分歧，只需 \ref{def:XIII.2.1.1} 的条件对 $S$ 的位于 $S$ 的极大点上方的几何点满足. 可假设除子 $D$ 严格正规交叉. 层 $F$ 可由 $U$ 的平展覆盖 $V$ 表示. 若 $s$ 是 $S$ 的几何点，$y$ 是 $Y_s$ 的极大点，记 $S$ 为 $S$ 在 $s$ 处的严格局部化，$X$ 为 $X$ 在 $y$ 处的严格局部化，$U = U \times_X X$，$V = V \times_X X$. 若 \ref{def:XIII.2.1.1} 的条件对位于 $S$ 的极大点上方的几何点满足，由下面 \ref{prop:XIII.5.5} 得出 $V$ 是在 $X$ 上相对 $S$ 温和分歧的 $U$ 的覆盖；从而 $V$ 是在 $X$ 上相对 $S$ 温和分歧的 $U$ 的平展覆盖.

b) 设 $F$ 是 $U$ 上在 $X$ 上相对 $S$ 温和分歧的群层. 若 $s$ 是 $S$ 的几何点，$y$ 是 $Y_s$ 的极大点，记 $K$ 为 $\mathcal{O}_{X_s, y}$ 的分式域. 假设对每个点 $s$ 和每个点 $y$，谱表示 $F|_K$ 的 $K$-代数 $L$ 的秩与 $\mathcal{O}_{X_s}$ 的剩余特征 $p$ 互素. 有时滥用语言说 $F$ 与 $S$ 的剩余特征互素\index[term]{yushengyutezhenghusu@与剩余特征互素}. 此时，$F$ 下的每个挠子 $P$ 在 $X$ 上相对 $S$ 温和分歧.

设 $R$ 是 $\mathcal{O}_{X_s, y}$ 在 $y$ 处的严格局部化，$K$ 是它的分式域，$F$ 是 $F$ 在 $K$ 上的逆像. 证明可假设 $F$ 常. 由于 $F$ 在 $X$ 上相对 $S$ 温和分歧，$F$ 可由 $K$-代数 $L = \prod L_i$ 的谱表示，其中诸 $L_i$ 是 $K$ 的与 $p$ 互素的阶扩张. 可以找到 $K$ 的与 $p$ 互素的阶扩张 $K'$，使 $F|_{K'}$ 是常层. 由 2.0.3，为证明 $P|_K$ 在 $R$ 上温和分歧，只需看 $P|_{K'}$ 在 $R$ 在 $K'$ 中的整闭包上温和分歧，归结到 $F$ 常的情形. 此后假设 $F$ 常. 表示 $P|_K$ 的 $K$-代数 $H$ 是彼此同构的 $K$ 的扩张 $H_i$ 的积. 由于 $H$ 的秩与 $p$ 互素，$[H_1 : K]$ 也如此，这证明 $H$ 在 $X$ 上相对 $S$ 温和分歧.

c) 设 $X$ 是正则概形，$D$ 是 $X$ 的\emph{正规交叉除子} (SGA 5 I 3.1.5)，$U = X - \mathrm{Supp}\, D$，$F$ 是 $U$ 上的集合层. 若 $y$ 是 $\mathrm{Supp}\, D$ 的极大点，记 $K$ 为 $\mathcal{O}_{X, y}$ 的分式域. 称 $F$ 相对 $D$ 温和分歧\index[term]{wenhefenqi@温和分歧}，若对 $\mathrm{Supp}\, D$ 的每个极大点 $y$，$F|_K$ 可由在 $\mathcal{O}_{X, y}$ 上温和分歧的 $K$-代数表示.
\end{remark}

\begin{theorem}\label{thm:XIII.2.4}
设 $f \colon X \to S$ 是 $S$-概形，$D$ 是 $X$ 上相对 $S$ 正规交叉 (2.1) 的除子，$Y = \mathrm{Supp}\, D$，$U = X - Y$，$i \colon U \to X$ 是典范浸入. 设 $F$ 是 $U$ 上满足下列条件之一的集合层 (resp. 群层)：
\begin{enumerate}[a)]
\item $F$ 局部于 $X$ 和 $S$ 的平展拓扑是 $S$ 上集合层 (resp. 可构群层) 的逆像.
\item $F$ 在 $U$ 上局部常可构，且在 $X$ 上相对 $S$ 温和分歧.
\end{enumerate}
则有下列结论：
\begin{enumerate}[1)]
\item $(F, i)$ 相对 $S$ 在维数 $\leqslant 0$ 时上同调真 (resp. 对每个态射 $h \colon S' \to S$，若 $i' \colon U' \to X'$ 是 $i$ 在 $S'$ 上的逆像，$F' = F|_{U'}$，$k = h_{(X)}$，典范态射
\[
\Psi \colon k^{*}(R^1_t i_* F) \longrightarrow R^1_t i'_* F'
\]
是同构).
若 $F$ 是与 $S$ 的剩余特征互素 (2.3.b)) (在 $X$ 上相对 $S$ 温和分歧) 的群层，则 $(F, i)$ 相对 $S$ 在维数 $\leqslant 1$ 时上同调真.
\item 若 $F$ 是可构集合层 (resp. 群层)，$i_* F$ (resp. $R^1_t i_* F$) 可构.
\end{enumerate}
\end{theorem}

\begin{proof}
对每个 $S$-概形 $S'$，考虑下列所有方形笛卡儿的图式：
\[
\begin{tikzcd}[row sep=large, column sep=large]
U \arrow[d,"{i}"] & U' \arrow[l,"g"] \arrow[d,"{i'}"] \\
X \arrow[d,"{f}"] & X' \arrow[l,"k"] \arrow[d,"{f'}"] \\
S & S' \arrow[l,"h"]
\end{tikzcd}
\]
由于问题局部于 $X$ 的平展拓扑，可假设 $D$ 是相对 $S$ \emph{严格正规交叉除子} (2.1)；而且，把 $X$ 限制到 $Y$ 的邻域后可假设 $X$ 在 $S$ 上光滑.

**\ref{thm:XIII.2.4} 1) 的证明**

2.4.1. 满足 a) 的集合层的情形. 可假设 $F = g^{*} G$，其中 $G$ 是 $S$ 上的层. 由 (SGA 4 XVI 3.2) 得出典范态射
\[
(2.4.1.1) \quad f^{*} G \longrightarrow i_* F
\]
是同构. 对每个 $S$-概形 $h \colon S' \to S$，同样有同构 $f'^{*} G' \to i'_* F'$；从而典范态射
\[
\phi \colon k^{*}(i_* F) \longrightarrow i'_* F'
\]
等同于自然同构
\[
k^{*} f^{*} G \simeq f'^{*} G'.
\]

2.4.2. 满足 b) 的集合层的情形. 需证明 $\phi$ 是同构，只需看它在 $X'$ 的每个几何点 $x'$ 处如此. 设 $S$ (resp. $X$，resp. $S'$，resp. $X'$) 是 $S$ (resp. $X$，resp. $S'$，resp. $X'$) 在 $x'$ 处的严格局部化，令 $U = U_{(X)}$，$U' = U_{(X')}$，等等. 态射 $\phi_x$ 等同于典范态射
\[
\phi \colon H^0(U, F) \longrightarrow H^0(U', F')
\]
可以找到 $U$ 的 \ref{prop:XIII.5.4} 中出现类型的主覆盖 $V$，使 $F$ 在 $V$ 上的逆像是取值 $C$ 的常层. 若 $\Pi$ 是 $V$ 在 $U$ 上的 伽罗瓦群，$\Pi$ 作用在 $F|_V$ 上，有
\[
(2.4.2.1) \quad H^0(U, F) \simeq H^0(V, C_V)^{\Pi}
\]
其中第二个成员表示 $H^0(V, C_V)$ 的在 $\Pi$ 下不变的元素的集. 由于 $V' = V \times_U U'$ 是 $U'$ 的群 $\Pi' \simeq \Pi$ 的主覆盖，看到态射 $\phi$ 由典范态射
\[
H^0(V, C_V) \longrightarrow H^0(V', C_{V'})
\]
取 $\Pi$ 的不变元得到. 由于 $V$ 和 $V'$ 连通 (\cref{prop:XIII.5.4})，该态射从而 $\phi$ 是同构.

注意若 $F$ 还是群层，$P$ 是 $U$ 上在 $X$ 上相对 $D$ 温和分歧的群 $F$ 的挠子，由前面证明和 2.2 得出 $(P F, i)$ 相对 $S$ 在维数 $\leqslant 0$ 时上同调真.

2.4.3. 群层的情形. 为证明 $\Psi$ 是同构，只需证明对 $Y'$ 的每个几何点 $y'$，态射
\[
\Psi_{y'} \colon (k^{*}(R^1_t i_* F))_{y'} \longrightarrow (R^1_t i'_* F')_{y'}
\]
是同构. 由 \ref{def:XIII.2.1.2}，$\Psi_{y'}$ 等同于典范态射
\[
\Psi \colon H^1_t(U, F) \longrightarrow H^1_t(U', F').
\]
设 $\widetilde{U}$ 是 $U$ 的温和分歧万有覆盖 (2.1.3)，$\widetilde{F}$ 是 $F$ 在 $\widetilde{U}$ 上的逆像. 由情形 a) 的 \ref{cor:XIII.5.7} 和情形 b) 的 \ref{prop:XIII.5.5} 得出，$H^1_t(U, F)$ 等同于 $H^1(U, F)$ 的在 $\widetilde{U}$ 上的逆像平凡的 $F$-挠子类组成的子集. 另一方面，经典推理 (参见 IX 5，第 231 页) 表明在 $\widetilde{U}$ 上的逆像平凡的 $H^1(U, F)$ 的元素集等同于 $H^1(\pi^t_1(U), H^0(\widetilde{U}, \widetilde{F}))$. 由此得到典范同构
\[
(2.4.3.1) \quad H^1_t(U, F) \xrightarrow{\sim} H^1(\pi^t_1(U), H^0(\widetilde{U}, \widetilde{F})).
\]
从而态射 $\Psi$ 等同于典范态射
\[
H^1(\pi^t_1(U), H^0(\widetilde{U}, \widetilde{F})) \longrightarrow H^1(\pi^t_1(U'), H^0(\widetilde{U}', \widetilde{F}')).
\]
证明该态射是同构. 态射 $\pi^t_1(U') \to \pi^t_1(U)$ 由 \ref{cor:XIII.5.6} 是同构，态射 $H^0(\widetilde{U}, \widetilde{F}) \to H^0(\widetilde{U}', \widetilde{F}')$ 也如此. 确实，在情形 b) 中显然，因为 $F$ 常且 $\widetilde{U}$ 和 $\widetilde{U}'$ 连通. 在情形 a) 中，设 $G$ 是 $S$ 上使 $F = g^{*} G$ 的可构群层；由于态射 $\widetilde{U} \to S$ 和 $\widetilde{U}' \to S'$ 0-无环 (\cref{cor:XIII.5.7})，有
\[
H^0(\widetilde{U}, \widetilde{F}) \simeq H^0(S, G) \simeq H^0(S', G') \simeq H^0(\widetilde{U}', \widetilde{F}'),
\]
这蕴含 $\Psi$ 是同构. \ref{thm:XIII.2.4} 1) 的最后一个断言由前面，考虑到 \ref{rem:XIII.2.3} b) 得出.

**\ref{thm:XIII.2.4} 2) 的证明** — 满足 a) 的可构集合层情形由 (2.4.1.1) 立即得出. 设 $F = g^{*} G$ 是满足 a) 的群层，其中 $G$ 可构；可假设 $S$ 仿射；设 $(S_j)_{j \in J}$ 是 $S$ 的既约闭子概形的有限族，其并覆盖 $S$，使 $G$ 在 $S_j$ 上的逆像是局部常层. 考虑到 \ref{thm:XIII.2.4} 1)，为建立 $R^1_t i_* F$ 可构，只需看它在基变换 $S_j \to S$ ($j \in J$) 后如此. 归结到 $F$ 局部常的情形 b).

此后假设 $F$ 是满足 b) 的集合层或群层. 由于问题局部于 $X$ 的平展拓扑，可假设 $X$ 在 $S$ 上有限表现，经极限过渡可假设 $X$ 和 $S$ 诺特.

设 $D = \sum_{1 \leqslant i \leqslant r} \mathrm{div}\, f_i$，其中对 $\mathrm{Supp}\, D$ 的每个点 $x$，若 $I(x)$ 是使 $f_i(x) = 0$ 的 $i$ 的集，子概形 $V((f_i)_{i \in I(x)})$ 在 $S$ 上光滑且在 $X$ 中余维数 $\mathrm{card}\, I(x)$. 设 $P$ 是 $[1, r]$ 的部分集，对每个 $I \in P$，令
\[
X_I = \Big(\bigcap_{i \in I} V(f_i)\Big) \cap \Big(\bigcap_{i \notin I} X_{f_i}\Big).
\]
设 $z$ 是 $X_I$ 的点. 先限制到 $z$ 的平展邻域，可以找到 $X$ 的含 $z$ 的开集 $W$ 和 $U \cap W$ 的在 $W$ 上相对 $S$ 温和分歧、\ref{rem:XIII.5.6.1} 中考虑类型的主覆盖 $V$，使 $F$ 在 $V$ 上的逆像是取值 $C$ 的常层. 设 $\pi$ 是该覆盖 $V$ 的 伽罗瓦群. 对 $X_I$ 的每个几何点 $x$，由 (2.4.2.1) 有
\[
H^0(U, F) \simeq H^0(V, C_V)^{\pi}.
\]
由此得出 $i_* F|_{X_I \cap W}$ 局部常 (SGA 4 IX 2.13)，从而 $i_* F$ 可构.

最后证明 $F$ 是局部常群层时 $R^1_t i_* F|_{X_I}$ 可构. 若 $x$ 是 $X$ 的几何点，在 (2.4.3.1) 中得到表达式
\[
H^1_t(U, F) \xrightarrow{\sim} H^1(\pi^t_1(U), H^0(\widetilde{U}, \widetilde{F})).
\]
若 $p$ 是 $X$ 的剩余特征，由 \ref{cor:XIII.5.6} 有 $\pi^t_1(U) = \prod_{\ell \neq p} \mathbb{Z}_\ell(1)^{\mathrm{card}\, I}$. 设 $L$ 是整除有限群 $H^0(\widetilde{U}, \widetilde{F})$ 的阶的素数集，令 $K = \prod_{\ell \in L - \{p\} \cap L} \mathbb{Z}_\ell(1)^{\mathrm{card}\, I}$. 由 \cite[第一章 §5 例 2]{exp13-4} 有
\[
H^1_t(U, F) \simeq H^1(K, H^0(\widetilde{U}, \widetilde{F})).
\]
由于 $K$ 拓扑有限型且 $H^0(\widetilde{U}, \widetilde{F})$ 有限，首先得出层 $R^1_t i_* F|_{X_I}$ 的纤维有限. 另一方面，$L$ 不依赖点 $x$. 对每个 $q \in L$，设 $X_{I,q}$ 是 $X_I$ 的方程 $q = 0$ 的闭集，$X_{I'}$ 是 $X_I$ 的诸 $X_{I,q}$ 的并的余集的开集. 则 $R^1_t i_* F|_{X_{I,q}}$ 和 $R^1_t i_* F|_{X_{I'}}$ 局部常；确实，$X_{I,q}$ (resp. $X_{I'}$) 的几何点的特殊化箭对群 $K$ (\cref{rem:XIII.5.6.1}) 从而对集合 $H^1_t(U, F)$ 诱导同构，且可以应用 SGA 4 IX 2.13.
\end{proof}

\begin{corollary}\label{cor:XIII.2.5}
设 $f \colon X \to S$ 是态射，$D$ 是 $X$ 上相对 $S$ 正规交叉 (2.1) 的除子，$Y = \mathrm{Supp}\, D$，$U = X - Y$，$i \colon U \to X$ 是典范浸入. 设 $\Phi$ 是 $U$ 上的叠，假设给定平展满射态射 $S_1 \to S$ 和 $X_2 \to X \times_S S_1$，$S_1$ 上的叠 $\Psi$ 和同构 $\Phi|_{U_2} \simeq \Psi|_{U_2}$ (参见 2.2).
则对每个态射 $h \colon S' \to S$，若 $k = h_{(X)}$，典范函子
\[
\phi \colon k^{*} i_* \Phi \longrightarrow i'_* \Phi'
\]
忠实全. 若 $\Psi$ 可构，典范函子
\[
\psi \colon k^{*} i_{*t} \Phi \longrightarrow i'_{*t} \Phi'
\]
是范畴等价.
而且，若叠 $\Psi$ 可构 (resp. $\Psi$ 1-可构 (0) 且 $S$ 局部诺特)，$i_* \Phi$ 可构 (resp. $i_{*t} \Phi$ 1-可构).
\end{corollary}

证明 $\phi$ 忠实全. 只需看对 $X'$ 的每个几何点 $x'$，函子
\[
\phi \colon \Phi(U) \longrightarrow \Phi'(U')
\]
如此 (沿用 2.4.2 的记号). 设 $a$，$b$ 是 $\Phi(U)$ 的两个元素，$a'$，$b'$ 是它们在 $\phi$ 下的像. 由于态射 $U \to S$ 局部 0-无环 (\cref{cor:XIII.5.7})，有同构
\[
H^0(U, S_\Phi) \xrightarrow{\sim} H^0(S, S_\Psi).
\]
从而 $a$ 和 $b$ 由 $\Psi$ 的元素经逆像得到，故 $F = \mathrm{Hom}_U(a, b)$ 也如此. 由于层 $F' = \mathrm{Hom}_{U'}(a', b')$ 是 $F$ 在 $U'$ 上的逆像，由 2.4.1 得出典范态射
\[
H^0(U, F) \longrightarrow H^0(U', F')
\]
是同构，这证明 $\phi$ 忠实全.

证明对 $X'$ 的每个几何点 $x'$，函子
\[
\psi \colon i_{*t} \Phi(X') \longrightarrow i'_{*t} \Phi'(X')
\]
是等价. 由前面 $\psi$ 忠实全. 证明 $\psi$ 本质满射. 设 $a'$ 是 $\Phi'(U')$ 的在 $X'$ 上相对 $S'$ 温和分歧 (2.2) 的元素，证明它是 $\Phi(U)$ 的温和分歧元素的像. 由 \ref{thm:XIII.2.4} 1) 得出典范态射
\[
H^0(U, S_\Phi) \longrightarrow H^0(U', S_{\Phi'})
\]
是同构. 设 $G'$ 是 $a'$ 生成的 $\Phi'$ 的极大子叠丛；则存在 $\Phi$ 的极大子叠丛 $G$ (某个 $S$ 上叠的逆像)，使
\[
m^{*} G \simeq G',
\]
其中 $m$ 是态射 $U' \to U$. 典范函子
\[
(\ast) \quad k^{*} i_{*t} G \longrightarrow i'_{*t} G'
\]
是等价，因为 $G$ 等同于来自 $S$ 的可构群层下的挠子叠丛，且可以应用 \ref{thm:XIII.2.4} 1). 由 $(\ast)$ 得出存在 $G(U)$ 的在 $X$ 上相对 $S$ 温和分歧的元素 $a$，使它在 $U'$ 上的逆像是 $a'$，这证明 $\psi$ 是等价.

$\Psi$ 可构时 $i_* \Phi$ 也可构；确实 $i_* \Phi$ 的对象 $x$ 局部于 $S$ 的平展拓扑是 $\Psi$ 的对象 $y$ 的逆像；由 2.4.1.1 得出 $\mathrm{Aut}(x)$ 是 $\mathrm{Aut}(y)$ 的逆像，故可构. 最后，$\Psi$ 1-可构时由下面 \ref{cor:XIII.6.3} 得出 $i_{*t} \Phi$ 也 1-可构.

\begin{corollary}\label{cor:XIII.2.6}
记号同 2.4. 假设 $S$ 在使 $Y_s \neq \varnothing$ 的每个点 $s$ 处特征零. 则若 $F$ 是 $U$ 上局部常可构群层 (resp. $U$ 上局部于 $X$ 和 $S$ 是 $S$ 上\emph{可构叠}的逆像的可构叠)，$(F, i)$ 相对 $S$ 在维数 $\leqslant 1$ 时上同调真.
\end{corollary}

由于 $U$ 上每个可构集合层在 $X$ 上相对 $S$ 温和分歧，该推论由 \ref{thm:XIII.2.4} (resp. \ref{cor:XIII.2.5}) 得出.

\begin{corollary}\label{cor:XIII.2.7}
记号同 \ref{thm:XIII.2.4}，但再给定 $S$-概形 $T$、紧合态射 $p \colon X \to T$，并假设 $X$ 和 $T$ 在 $S$ 上有限表现；设 $q = pi$. 设 $F$ 是 $U$ 上满足 \ref{thm:XIII.2.4} 的条件 a) 或 b) 之一的可构集合层 (resp. 满足 \ref{thm:XIII.2.4} 的条件 a) 或 b) 之一的群层，resp. $U$ 上局部于 $X$ 和 $S$ 是 $S$ 上可构叠 $G$ 的逆像的叠). 则有下列结论：
\begin{enumerate}[1)]
\item $(F, q)$ 相对 $S$ 在维数 $\leqslant 0$ 时上同调真 (resp. 对每个态射 $h \colon S' \to S$，若 $m = h_{(T)}$，典范态射：
\[
\Theta \colon m^{*}(R^1_t q_* F) \longrightarrow R^1_t q'_* F'
\]
是同构，resp. 对每个态射 $S' \to S$，典范态射
\[
\xi \colon m^{*}(q_{*t} F) \longrightarrow q'_{*t} F'
\]
是等价).
\item 层 $q_* F$ (resp. 层 $R^1_t q_* F$，resp. 叠 $q_{*t} F$) 可构. 在最后情形，若假设 $S$ 局部诺特且 $G$ 1-可构，$q_{*t} F$ 也 1-可构.
\end{enumerate}
\end{corollary}

第一部分由 \ref{thm:XIII.2.4}，\ref{cor:XIII.2.5} 和 \ref{cor:XIII.1.8} 的证明立即得出. 证明 2). 若 $F$ 是 $U$ 上满足 \ref{thm:XIII.2.4} a) 或 \ref{thm:XIII.2.4} b) 的可构集合层，由 \ref{thm:XIII.2.4} 2) 得出 $i_* F$ 可构；故 $q_* F = p_*(i_* F)$ 也可构 (SGA 4 XIV 1.1).

设 $F$ 是 $U$ 上满足 \ref{thm:XIII.2.4} a) 或 \ref{thm:XIII.2.4} b) 的可构群层，证明 $R^1_t q_* F$ 可构. 经极限过渡 (EGA IV 8.10.5 和 17.7.8) 并用 1)，可假设 $S$ 诺特. 设 $\Phi$ 是 $X$ 上的叠，它在每个平展于 $X$ 上的概形 $X'$ 处的纤维由 $U' = U \times_X X'$ 上在 $X'$ 上相对 $S$ 温和分歧的群 $F|_U$ 的挠子组成；故有
\[
S(i_{*t} \Phi) \simeq R^1_t i_* F
\]
且由 \ref{thm:XIII.2.4} 2) 该层可构. 由下面 \ref{cor:XIII.6.3} 得出 $S(p_* \Phi)$ 可构，i.e. $R^1_t q_* F$ 可构.

最后，若 $F$ 是 $U$ 上局部于 $X$ 和 $S$ 是 $S$ 上可构叠的逆像的叠，$i_{*t} F$ 可构，故 $q_{*t} F = p_*(i_{*t} F)$ 也可构. 若还假设 $S$ 局部诺特且 $S_G$ 可构，由 \ref{cor:XIII.6.3} 得出 $S(i_{*t} F)$ 可构；故 $S(q_* i_{*t} F)$，i.e. $S(q_{*t} F)$ 也可构 (\cref{cor:XIII.6.2}).

\begin{corollary}\label{cor:XIII.2.8}
设
\[
\begin{tikzcd}[row sep=large, column sep=large]
U \arrow[r,"i"] \arrow[dr,"g"] & X \arrow[d,"f"] \\
& S
\end{tikzcd}
\]
是概形的交换图式，其中 $U$ 是 $X$ 中相对 $S$ 正规交叉除子的余集的开集，$f$ 是有限表现紧合态射. 设 $L$ 是素数集. 假设 $g$ 局部 0-无环 (resp. 对 $L$ 局部 1-非球). 则若 $F$ 是 $U$ 上局部常可构、在 $X$ 上相对 $S$ 温和分歧的集合层 (resp. $L$-群层)，$f_* F$ (resp. $R^1_t f_* F$) 局部常可构，且 $(F, f)$ 相对 $S$ 在维数 $\leqslant 0$ 时上同调真 (resp. $R^1_t f_* F$ 的构造与任何基变换 $S' \to S$ 交换). 在非 resp. 情形，若 $F$ 是群层，对 $S$ 的几何点的每个特殊化 $s_1 \to s_2$，特殊化态射
\[
(R^1_t f_* F)_{s_2} \longrightarrow (R^1_t f_* F)_{s_1}
\]
单射.
\end{corollary}

该推论由 \ref{cor:XIII.2.6} 和 \ref{prop:XIII.1.14} (resp. $R^1_t f_* F$ 的 \ref{prop:XIII.1.14} 类似物，按 loc. cit. 证明) 立即得出.

\begin{corollary}\label{cor:XIII.2.9}
设
\[
\begin{tikzcd}[row sep=large, column sep=large]
U \arrow[r,"i"] \arrow[dr,"g"] & X \arrow[d,"f"] \\
& S
\end{tikzcd}
\]
是概形的交换图式，其中 $U$ 是 $X$ 中相对 $S$ 正规交叉除子的余集的开集，$f$ 是有限表现紧合光滑态射. 设 $L$ 是与 $S$ 的剩余特征不同的素数集. 设 $F$ 是 $U$ 上局部常可构、在 $X$ 上相对 $S$ 温和分歧的 $L$-群层. 则 $R^1 f_* F$ 局部常可构，且 $(F, f)$ 相对 $S$ 在维数 $\leqslant 1$ 时上同调真.
\end{corollary}

该推论由 \ref{cor:XIII.2.8} 和 $R^1_t f_* F = R^1 f_* F$ (2.3 b)) 的事实得出.

\refstepcounter{theorem}\noindent\textbf{\thesection.\arabic{theorem}.}\quad 若 $U$ 是连通概形，$a$ 是 $U$ 的几何点，$L$ 是素数集，记
\[
(2.10.0) \quad \pi^L_1(U, a)\index[notations]{$\pi_1^L(U, a)$}
\]
为 $\pi_1(U, a)$ 的阶的所有素因子都在 $L$ 中的有限商的投射极限.

将定义基本群的特殊化态射，推广 X.2.

设 $g \colon U \to S$ 是纤维几何连通的相容态射 (resp. $g = f i$ 形式的态射，其中 $f \colon X \to S$ 是有限表现紧合态射，$i \colon U \to X$ 是使 $U$ 为 $X$ 中相对 $S$ 正规交叉除子的余集的开浸入 (参见 \ref{cor:XIII.2.8})). 设 $L$ 是素数集，在非 resp. 情形假设对每个有限常 $L$-群 $C$，$(C_U, g)$ 相对 $S$ 在维数 $\leqslant 1$ 时上同调真. 设 $s_1 \to s_2$ 是 $S$ 的几何点的特殊化态射，$S$ 是 $S$ 在 $s_2$ 处的严格局部化，$U = U \times_S S$. 有交换图式
\[
\begin{tikzcd}[row sep=large, column sep=large]
U_{s_1} \arrow[d] & U \arrow[l,"{h_1}"] \arrow[r,"{h_2}"] \arrow[d,"{g}"] & U_{s_2} \arrow[d] \\
s_1 & S \arrow[l] \arrow[r] & s_2.
\end{tikzcd}
\]
若 $a_1$ 是 $U_{s_1}$ 的几何点，$a_2$ 是 $U_{s_2}$ 的几何点，态射 $h_1$ 和 $h_2$ 定义典范态射
\[
\pi_1 \colon \pi^L_1(U_{s_1}, a_1) \longrightarrow \pi^L_1(U, a_1), \quad \pi_2 \colon \pi^L_1(U_{s_2}, a_2) \longrightarrow \pi^L_1(U, a_2)
\]
(resp. $\pi_1 \colon \pi^t_1(U_{s_1}, a_1) \longrightarrow \pi^t_1(U, a_1)$，$\pi_2 \colon \pi^t_1(U_{s_2}, a_2) \longrightarrow \pi^t_1(U, a_2)$) (V 7 和 2.1.5.2). 上同调真性假设 (resp. \ref{cor:XIII.2.8}) 证明 $\pi_2$ 是同构. 若选 $a_1$ 到 $a_2$ 的道路类，得到同构
\[
\pi_{12} \colon \pi^L_1(U, a_1) \xrightarrow{\sim} \pi^L_1(U, a_2) \quad \text{(resp. } \pi_{12} \colon \pi^t_1(U, a_1) \xrightarrow{\sim} \pi^t_1(U, a_2)\text{)}
\]
从而有态射 $\pi = \pi_2^{-1} \pi_{12} \pi_1$
\[
\pi \colon \pi^L_1(U_{s_1}, a_1) \longrightarrow \pi^L_1(U_{s_2}, a_2) \quad \text{(resp. } \pi \colon \pi^t_1(U_{s_1}, a_1) \longrightarrow \pi^t_1(U_{s_2}, a_2)\text{)}.
\]
换 $a_1$ 到 $a_2$ 的道路类，等价于用 $\pi^L_1(X_{s_2}, a_2)$ (resp. $\pi^t_1(X_{s_2}, a_2)$) 的内自同构修改 $\pi$. 称上述定义的一个态射为与特殊化态射 $s_1 \to s_2$ 关联的基本群特殊化态射，简记为
\[
\pi \colon \pi^L_1(X_{s_1}) \longrightarrow \pi^L_1(X_{s_2}) \quad \text{(resp. } \pi \colon \pi^t_1(X_{s_1}) \longrightarrow \pi^t_1(X_{s_2})\text{)}.
\]

\begin{lemma}\label{lem:XIII.2.11}
设 $f \colon X \to S$ 是有限表现紧合态射，$D$ 是 $X$ 上相对 $S$ 正规交叉除子，$Y = \mathrm{Supp}\, D$，$U = X - Y$，$i \colon U \to X$ 是典范态射，$s_1 \to s_2$ 是 $S$ 的几何点的特殊化态射，$y_1$ 是 $Y_{s_1}$ 的几何点，$y_2$ 是 $Y_{s_2}$ 的几何点，使 $y_1$ 在 $X$ 上的投影 $z_1$ 是 $y_2$ 的投影 $z_2$ 的一般化. 设 $I_{y_1}$ 是 $\pi^t_1(U_{s_1})$ 在 $y_1$ 处的惯性子群. 则 $I_{y_1}$ 经特殊化态射
\[
\pi \colon \pi^t_1(U_{s_1}) \longrightarrow \pi^t_1(U_{s_2})
\]
的像是 $\pi^t_1(U_{s_2})$ 在 $y_2$ 处的惯性子群.
\end{lemma}

设 $X$ (resp. $\widetilde{X}$) 是 $X$ 在 $y_2$ (resp. $y_1$) 处的严格局部化，$U = U \times_X X$ (resp. $\widetilde{U} = U \times_X \widetilde{X}$). 有典范态射 $\widetilde{U} \to U$，由 \ref{def:XIII.1.10} 得出有交换图式
\[
\begin{tikzcd}[row sep=large, column sep=large]
\pi^t_1(\widetilde{U}_{s_1}) \arrow[r,"{\pi'}"] \arrow[d] & \pi^t_1(U_{s_2}) \arrow[d] \\
\pi^t_1(U_{s_1}) \arrow[r,"{\pi}"] & \pi^t_1(U_{s_2}),
\end{tikzcd}
\]
其中 $\pi'$ 是典范态射 $\pi^t_1(\widetilde{U}_{s_1}) \to \pi^t_1(U_{s_1})$ 与特殊化态射的复合. 由于 $\pi^t_1(\widetilde{U}_{s_1})$ (resp. $\pi^t_1(U_{s_1})$) 是 $\pi^t_1(U_{s_1})$ 在 $y_1$ 处 (resp. $\pi^t_1(U_{s_2})$ 在 $y_2$ 处) 的惯性群，只需证明 $\pi'$ 满射. 但这由 \ref{cor:XIII.5.6} 中得到的表达式得出.

\begin{corollary}\label{cor:XIII.2.12}
设 $X$ 是特征 $p \geqslant 0$ 的可分闭域 $k$ 上亏格 $g$ 的紧合光滑连通曲线. 设 $U$ 是从 $X$ 去掉 $n$ 个互异闭点 $a_1, \ldots, a_n$ 得到的开集. 则温和分歧基本群 $\pi^t_1(U)$ (2.1.3) 可由 $2g + n$ 个元素 $x_i$，$y_i$，$\sigma_j$ ($1 \leqslant i \leqslant g$，$1 \leqslant j \leqslant n$) 生成，使 $\sigma_j$ 是对应 $a_j$ 的惯性群的生成元，且满足关系
\[
(\ast) \quad \prod_{1 \leqslant i \leqslant g} (x_i y_i x_i^{-1} y_i^{-1}) \cdot \prod_{1 \leqslant j \leqslant n} \sigma_j = 1.
\]
对每个阶与 $p$ 互素、由满足关系 $(\ast)$ 的元素 $\overline{x_i}$，$\overline{y_i}$，$\overline{\sigma_j}$ 生成的有限群 $G$，存在 $U$ 的群 $G$ 的平展覆盖，对应把 $x_i$，$y_i$，$\sigma_j$ 分别送到 $\overline{x_i}$，$\overline{y_i}$，$\overline{\sigma_j}$ 的同态 $\pi^t_1(U) \to G$. 换言之，若 $p'$ 表示与 $p$ 不同的素数集，$\pi^{p'}_1(U)$ 是由生成元 $x_i$，$y_i$，$\sigma_j$ 只服从关系 $(\ast)$ 生成的 pro-$p'$-群.
\end{corollary}

\begin{proof}
可假设 $k$ 代数闭. 先假设 $k$ 特征零. 则存在 $k$ 的在 $\mathbb{Q}$ 上超越次数有限的代数闭子扩张 $k'$，使 $X$ 由定义在 $k'$ 上的紧合光滑曲线 $X'$ 经标量扩张得到，且可假设点 $a_1, \ldots, a_n$ 来自 $X'$ 的有理点 $a'_1, \ldots, a'_n$. 由于 $k'$ 在 $\mathbb{Q}$ 上超越次数有限，可以找到 $k'$ 到复数域 $\mathbb{C}$ 的嵌入；设 $\widetilde{U} = U' \times_{k'} \mathbb{C}$. 设 $k''$ 是 $k$ 的代数闭扩张，使有 $k$ 和 $\mathbb{C}$ 到 $k''$ 的 $k'$-态射. 若 $g' \colon U' \to k'$ 是结构态射，$F$ 是 $U''$ 上有限常群层，由 \ref{cor:XIII.2.9} 得出特殊化态射
\[
(R^1 g'_* F')_{\mathbb{C}} \longrightarrow (R^1 g'_* F')_{k''} \longleftarrow (R^1 g'_* F')_k
\]
是同构. 用基本群的语言，这证明有同构 (模内自同构确定)
\[
\pi_1(U) \xrightarrow{\sim} \pi_1(\widetilde{U}),
\]
显然该同构把 $X' - U'$ 的某点处的惯性群变成同一点处的惯性群. 可假设 $k = \mathbb{C}$. 此时由 \emph{Riemann 存在定理} (\cref{cor:XII.5.2}) 得出基本群 $\pi_1(U)$ 正是 $\pi_1(U)$ 的关联解析空间的拓扑基本群关于有限指数子群拓扑的完备化. 后者可以用超越方式计算 \cite[第 7 章 §47]{exp13-3}；它可由 $2g + n$ 个元素 $x_i$，$y_i$，$\sigma_j$ 生成，其中 $\sigma_j$ 是中心在 $a_j$ 的小圆盘的局部基本群 $\pi_1(D_j)$ 的生成元的像，i.e. 对应点 $a_j$ 的惯性群的生成元，这些元素只满足关系 $(\ast)$.

若 $k$ 的特征 $p > 0$，可以找到剩余域 $k$、分式域 $K$ 特征零的完备离散赋值环 $A$，和 $S = \mathrm{Spec}\, A$ 上连通、紧合且光滑的概形 $X_1$，使 $X_1 \times_S \mathrm{Spec}\, k \simeq X$ (\cref{cor:III.7.4}). 诸点 $a_j$ 提升为 $X_1$ 在 $S$ 上方的截面 $s_j$；设 $Y_{1j}$ 是底集为 $s_j(S)$ 的既约闭子概形，$Y_1$ 是诸 $Y_{1j}$ 的并，$U_1 = X_1 - Y_1$，$g_1 \colon U_1 \to S$ 是结构态射. 设 $K$ 是 $K$ 的代数闭扩张，$U = U_1 \times_S K$. 若 $C$ 是有限常群，由 \ref{cor:XIII.2.8} 得出特殊化态射
\[
(R^1 g_{1*} C_{U_1})_k \longrightarrow (R^1 g_{1*} C_{U_1})_K
\]
单射，且 $C$ 的阶与 $p$ 互素时甚至双射. 用基本群的语言，这意味着特殊化态射 (\cref{def:XIII.1.10})
\[
\pi \colon \pi_1(U) \longrightarrow \pi^t_1(U)
\]
满射，且特殊化态射
\[
\pi^{p'}_1(U) \longrightarrow \pi^{p'}_1(U)
\]
双射. 最后，若 $x_i$，$y_i$，$\sigma_j$ 是 $\pi_1(U)$ 的生成元，其中 $\sigma_j$ 是对应点 $b_j = Y_{1j}(K)$ 的惯性群的生成元，则由 \ref{exa:XIII.1.11}，$\pi(\sigma_j)$ 是对应 $a_j$ 的惯性群的生成元，这完成证明.
\end{proof}

\begin{remark}[2003 年 (MR) 增补]
设 $k$ 是特征 $p > 0$ 的代数闭域，$X$ 是 $k$ 上亏格 $g$ 的紧合光滑连通代数曲线，$U$ 是 $X$ 的仿射开集，即 $X$ 的 $r \geqslant 1$ 个有理点的余集. 有基本群 $\pi_1(U)$，它的商 $\pi^t_1(U)$ 分类 $U$ 的在 $X - U$ 的点处温和分歧的有限平展覆盖，以及 $\pi^t_1(U)$ 的商 $\pi^{p'}_1(U)$ 分类 $U$ 的 伽罗瓦群阶与 $p$ 互素的 Galois 平展覆盖.

在 (Coverings of algebraic curves，Amer. J. Math. 79 (1957)，p. 825--856) 中，S. Abhyankar 对 $U$ 的有限平展覆盖的 伽罗瓦群的有限群结构提出若干猜想.

关于阶与 $p$ 互素、且是 $U$ 的连通平展覆盖的 伽罗瓦群的有限群的猜想，已在\cref{cor:XIII.2.12} 中证明并精确化. 在探讨 $\pi_1(U)$ 的有限商之前，先给出这个群的"大小"的若干指示. 设 $X = \mathrm{Spec}(A)$. 已知 $\mathrm{Hom}_{\mathrm{cont}}(\pi_1(U), \mathbb{Z}/p\mathbb{Z})$ 由 Artin-Schreier 理论描述 (\cref{cor:XI.6.9}). 该群同构于 $A/\wp(A)$，其中 $\wp$ 是把 $a$ 送到 $a^p - a$ 的 $A$ 到 $A$ 的映射. 先假设 $U$ 是仿射直线 $\mathbf{A}^1$，环为 $k[T]$. 设 $E$ 是 $k[T]$ 的形如 $\sum_i a_i T^i$ (其中 $p$ 整除 $i$ 时 $a_i = 0$) 的元素集. 立即看到复合映射 $E \to k[T] \to k[T]/\wp k[T]$ 双射. 于是系数 $a_i$ ($(i, p) = 1$) 表现得像参数化仿射直线的 $p$ 阶循环覆盖的空间的坐标. 特别由此得出 $\pi_1(\mathbf{A}^1)$ 不拓扑有限型，且若作代数闭域的扩张 $k \to k'$，自然态射 $\pi_1(\mathbf{A}^1 \times_k k') \to \pi_1(\mathbf{A}^1)$ (满射) 不双射，与紧合情形相反. 在一般情形，曲线 $U$ 实现为仿射直线上的有限概形，得出同样的现象对 $\pi_1(U)$ 发生 (而且更一般地对维数 $\geqslant 1$ 的任何连通有限型仿射 $k$-概形 $V$ 的 $\pi_1(V)$ 发生).

于是，若 $G$ 是有限群，记 $G^{(p')}$ 为 $G$ 的阶与 $p$ 互素的最大商群. 为使有限群 $G$ 是 $\pi_1(U)$ 的拓扑商，必须 $G^{(p')}$ 是 $\pi^{p'}_1(U)$ 的拓扑商，而原则上可以借助\cref{cor:XIII.2.12} 回答这个条件. 在前述文章中，Abhyankar 猜想这个必要条件也充分. 该猜想已被 M. Raynaud 在仿射直线的情形证明，被 D. Harbater 在一般情形证明 (Invent. Math. 116 (1994)，p. 425--462 和 117，p. 1--25). 例如，在仿射直线的情形 $\pi^{p'}_1(\mathbf{A}^1) = 1$，得出有限群 $G$ 是 $\mathbf{A}^1$ 的连通平展覆盖的 伽罗瓦群，当且仅当 $G^{(p')} = 1$，i.e. 当且仅当 $G$ 由它的 Sylow $p$-子群生成. 从而任何阶为 $p$ 的倍数的有限单群都合适.
\end{remark}

\section{上同调真性与一般局部无环性}

\begin{theorem}\label{thm:XIII.3.1}
设 $S$ 是生成点为 $s$ 的不可约概形，$X$ 和 $Y$ 是两个有限表现 $S$-概形，$f \colon X \to Y$ 是 $S$-态射. 对每个 $S$-概形 $S'$，用 $Y'$，$X'$，等等记 $Y$，$X$，等等经态射 $S' \to S$ 的逆像. 有下列性质：
\begin{enumerate}[1)]
\item a) 可以找到 $S$ 的非空开集 $S'$，使对 $X'$ 上每个有限常集合层 $F'$，$f'_* F'$ 可构，且 $(F', f')$ 相对 $S'$ 在维数 $\leqslant 0$ 时上同调真.
\item b) 设 $F$ 是 $X$ 上的可构集合层. 则可以找到 $S$ 的非空开集 $S'$ (依赖 $F$)，使 $f'_* F'$ 可构，且 $(F', f')$ 相对 $S'$ 在维数 $\leqslant 0$ 时上同调真.
\item 假设 $k(s)$ 的代数闭包 $k$ 上维数 $\leqslant \dim X_s$ 的有限型概形强可消解奇点 (SGA 5 I 3.1.5). 则还有下列性质：
\item a) 可以找到 $S$ 的非空开集 $S'$，使对 $X'$ 上每个与 $S$ 的剩余特征互素的阶有限常群层 $F'$，若 $\Phi'$ 是 $F'$ 下的挠子叠，$f'_* \Phi'$ 1-可构，且 $(F', f')$ 相对 $S'$ 在维数 $\leqslant 1$ 时上同调真.
\item b) 设 $L$ 是与 $S$ 的剩余特征不同的素数集，$\Phi$ 是 $X$ 上 (0) 的 1-可构 ind-$L$-叠，使对每个平展于 $X$ 上的概形 $X_1$ 和 $\Phi_{X_1}$ 的每对对象 $x$，$x_1$，层 $\mathrm{Hom}_{X_1}(x, x_1)$ 可构. 再假设 $S$ 局部诺特. 则可以找到 $S$ 的非空开集 $S'$，使 $f'_* \Phi'$ 1-可构，对 $(f'_* \Phi')_{Y_1}$ 的每对对象 $y$，$y_1$，$\mathrm{Hom}_{Y_1}(y, y_1)$ 可构，且 $(\Phi', f')$ 相对 $S'$ 在维数 $\leqslant 1$ 时上同调真.
\end{enumerate}
\end{theorem}

\begin{proof}
可假设 $S$ 仿射；由 SGA 4 VIII 1.1，可假设 $S$ 整；最后经极限过渡，可假设 $S$ 是 $\mathbb{Z}$ 上有限型代数的谱；特别 $S$ 此时诺特. 由于问题局部于 $Y$，可假设 $Y$ 仿射. 而且只需在有限扩张 $S' \to S$ 后证明定理，其中 $S'$ 是整概形，$S' \to S$ 由平展态射和有限根式满射态射复合而成.

**1) 常集合层的情形**

1) 1. 归结到 $X$ 在 $S$ 上正规的情形. 设 $X_{1s}$ 是 $(X_s)_{\mathrm{r\acute{e}d}}$ 的正规化；把 $S$ 限制到非空开集并作 $S$ 的根式扩张后，可假设 $X_{1s}$ 来自 $S$ 上正规的概形 $X_1$，且态射 $X_{1s} \to X_s$ 来自有限满射态射 $p \colon X_1 \to X$ (EGA IV 8.8.2 和 9.6.1). 假设定理对 $fp$ 已证. 把 $S$ 限制到开集后，可假设对 $X$ 上每个常集合层 $F$，$(p_* F, fp)$ 相对 $S$ 在维数 $\leqslant 0$ 时上同调真，且 $f_* p_*(p^{*} F)$ 可构. 由 \ref{cor:XIII.1.9}，$(p^{*} p_* F, f)$ 相对 $S$ 在维数 $\leqslant 0$ 时上同调真. 态射
\[
F \longrightarrow p_* p^{*} F = G
\]
是单态射. 已经由此得出，$f_* F$ 是 $f_* G$ 的子层，故 $f_* F$ 可构 (SGA 4 IX 2.9 (ii))，且 $(F, f)$ 相对 $S$ 在维数 $\leqslant -1$ 时上同调真.

设 $X_2 = X_1 \times_X X_1$，$q \colon X_2 \to X$ 是典范态射. 由 \ref{exa:XIII.1.11} 1)，有正合序列
\[
F \longrightarrow G \rightrightarrows q_* q^{*} F.
\]
由刚才对 $fq$ 代替 $f$ 证明的，可假设 (把 $S$ 限制到非空开集后) 对 $X$ 上每个常集合层 $F$，$(q^{*} F, fq)$ 相对 $S$ 在维数 $\leqslant -1$ 时上同调真，故 $(q_* q^{*} F, f)$ 相对 $S$ 在维数 $\leqslant -1$ 时上同调真. 由 \ref{prop:XIII.1.13} 1) 得出 $(F, f)$ 相对 $S$ 在维数 $\leqslant 0$ 时上同调真.

1) 2. 归结到 $X$ 在 $S$ 上正规仿射的情形.

设 $U_s$ 是 $X_s$ 的稠密仿射开集. 把 $S$ 限制到非空开集后，可假设 $U_s \to X_s$ 提升为相对 $S$ 概形支配的开浸入 $i \colon U \to X$ (EGA IV 8.9.1). 由于态射 $X \to S$ 正规，由 SGA 2 XIV 1.18 有：
\[
\mathrm{prof}^{\mathrm{ét}}_{S-U}(X) \geqslant 2;
\]
从而对 $X$ 上每个常层 $F$，典范态射
\[
F \longrightarrow i_* i^{*} F
\]
是同构. 由此得出，若假设定理对 $fi$ 和 $i$ 已证，则把 $S$ 限制到非空开集后，$(i_* F, i)$ 和 $(i_* F, fi)$ 相对 $S$ 在维数 $\leqslant 0$ 时上同调真. 故 $(F, f)$ 也如此 (1.6 2)). 又由于 $f_* F = (fi)_*(i_* F)$ 可构，这完成归结.

1) 3. 证明结束.

可假设 $S$ 正规 (EGA IV 7.8.3). 可以找到 $X_s$ 的紧化：
\[
\begin{tikzcd}[row sep=large, column sep=large]
X_s \arrow[r,"{j_s}"] \arrow[dr,"{f_s}"] & P_s \arrow[d,"{g_s}"] \\
& Y_s
\end{tikzcd}
\]
其中 $j_s$ 是支配开浸入，$g_s$ 是紧合态射；作 $k(s)$ 的根式扩张并把 $P_s$ 换成它的正规化 (不改变 $X_s$) 后，可假设 $P_s$ 几何正规. 把 $S$ 限制到非空开集并作根式满射扩张后，可假设上图式来自图式
\[
\begin{tikzcd}[row sep=large, column sep=large]
X \arrow[r,"{j}"] \arrow[dr,"{f}"] & P \arrow[d,"{g}"] \\
& Y
\end{tikzcd}
\]
其中 $P$ 是 $S$ 上正规概形，$j$ 是相对 $S$ 概形支配的开浸入，$g$ 是紧合态射 (EGA IV 6.9.1，9.9.4 和 9.6.1). 对 $X$ 上取值 $C$ 的每个有限常集合层 $F$，$j_* F$ 是取值 $C$ 的常层 (SGA 4 2.14.1)，经任何基变换 $S' \to S$ 后也如此；由此得出 $(F, j)$ 相对 $S$ 在维数 $\leqslant 0$ 时上同调真. 故 $(F, f)$ 也如此，因为 $g$ 紧合 (\cref{cor:XIII.1.8}). 由于 $g$ 紧合，$f_* F = g_* C_P$ 可构，这完成 1) a) 的证明.

**2) 可构集合层的情形**

设 $F$ 是 $X$ 上的可构集合层. 由 SGA 4 IX 2.14 (ii)，可以找到有限族态射 $p_i \colon Z_i \to X$，以及在每个 $Z_i$ 上取值 $C_i$ 的有限常集合层 $C_i$，使有单态射
\[
j \colon F \longrightarrow \prod_i p_{i*} C_i = G.
\]
由 1) a)，可假设 (把 $S$ 限制到非空开集后) 诸 $(C_i, fp_i)$ 相对 $S$ 在维数 $\leqslant 0$ 时上同调真，诸 $f_* p_{i*} C_i$ 可构. 已经得出 $(G, f)$ 相对 $S$ 在维数 $\leqslant 0$ 时上同调真 (\cref{cor:XIII.1.9})，故 $(F, f)$ 相对 $S$ 在维数 $\leqslant -1$ 时上同调真，且 $f_* F$ 可构. 设 $K$ 是余直和 $K = G \amalg_F G$；由于 $F$ 和 $G$ 可构，$K$ 也可构. 由前面得出，把 $S$ 限制到非空开集后可假设 $(K, f)$ 相对 $S$ 在维数 $\leqslant -1$ 时上同调真. 由 \ref{prop:XIII.1.13} 1) 得出 $(F, f)$ 相对 $S$ 在维数 $\leqslant 0$ 时上同调真.

**3) 常群层的情形**

若 $F$ 是 $X$ 上的常群层，记 $\Phi$ 为 $F$ 下的挠子叠.

3) 1. 先证明，把 $S$ 限制到非空开集后，对 $X$ 上每个与 $S$ 的剩余特征互素的阶有限常群层 $F$，$(F, f)$ 相对 $S$ 在维数 $\leqslant 0$ 时上同调真，且 $f_* \Phi$ 可构.

为此归结到 $X$ 在 $S$ 上光滑的情形. 作 $k(s)$ 的有限扩张 (可把它看作平展扩张与根式扩张的复合) 后，可以找到紧合满射态射 $p_s \colon X_{1s} \to X_s$，其中 $X_{1s}$ 是与 $X_s$ 同维数、在 $S$ 上光滑的概形，且把 $S$ 限制到非空开集后可假设 $p_s$ 来自紧合满射态射 $p \colon X_1 \to X$，其中 $X_1$ 在 $S$ 上光滑 (EGA IV 9.6.1 和 12.1.6). 设 $X_2 = X_1 \times_X X_1$，$q \colon X_2 \to X$ 是典范态射. 有 $X$ 上叠的正合图式
\[
\Phi \longrightarrow p_* p^{*} \Phi \rightrightarrows q_* q^{*} \Phi
\]
(1.11 2)). 假设定理在光滑情形已证，首先看到可假设 $f_* p_* p^{*} \Phi$ 可构；故 $f_* \Phi$ 也可构 (下面 \ref{lem:XIII.3.1.1}). 而且由 \ref{prop:XIII.1.6} 2)，可假设 $(p_* p^{*} \Phi, f)$ 相对 $S$ 在维数 $\leqslant 0$ 时上同调真；由此得出 $(\Phi, f)$ 相对 $S$ 在维数 $\leqslant -1$ 时上同调真. 可假设 $(q_* q^{*} \Phi, f)$ 相对 $S$ 在维数 $\leqslant -1$ 时上同调真，这蕴含 $(\Phi, f)$ 相对 $S$ 在维数 $\leqslant 0$ 时上同调真 (1.12 1)).

随后如 1) 2 归结到 $X$ 在 $S$ 上光滑仿射的情形. 设
\[
\begin{tikzcd}[row sep=large, column sep=large]
X \arrow[r,"{i}"] \arrow[dr,"{f}"] & P \arrow[d,"{q}"] \\
& Y
\end{tikzcd}
\]
是 $X$ 的紧化，其中 $i$ 是支配开浸入，$q$ 是紧合态射. 由于 $\dim P_s = \dim X_s$，可以对 $P_s$ 应用奇点消解假设. 作 $S$ 的平展扩张和根式扩张后，可以找到紧合态射 $r \colon Z \to P$，其中 $Z$ 在 $S$ 上光滑，$r^{-1}(X) \simeq X$，且 $r^{-1}(X)$ 是 $Z$ 中相对 $S$ 正规交叉除子的余集. 由 \ref{rem:XIII.2.3} b)，$F$ 下的每个挠子相对 $S$ 在 $Z$ 上温和分歧. 由 \ref{cor:XIII.2.7} 得出 $(F, f)$ 相对 $S$ 在维数 $\leqslant 0$ 时上同调真，这证明我们的断言.

3) 2. 归结到 $X$ 在 $S$ 上光滑的情形.

作 $k(s)$ 的有限扩张后，可以找到紧合满射态射 $p_s \colon X_{1s} \to X$，其中 $X_{1s}$ 在 $s$ 上光滑，且可假设 $p_s$ 来自紧合满射态射 $p \colon X_1 \to X$，其中 $X_1$ 在 $S$ 上光滑. 假设定理对 $fp$ 已证，对 $f$ 证明. 设 $F$ 是 $X$ 上与 $S$ 的剩余特征互素的阶有限常群层，$\Phi$ 是 $F$ 下的挠子叠. 设 $X_2 = X_1 \times_X X_1$，$X_3 = X_1 \times_X X_1 \times_X X_1$，$q \colon X_2 \to X$，$r \colon X_3 \to X$ 是典范态射. 由 \ref{exa:XIII.1.11} 2)，有叠的正合图式
\[
\Phi \longrightarrow p_* p^{*} \Phi \rightrightarrows q_* q^{*} \Phi \Rrightarrow r_* r^{*} \Phi.
\]
为证明 $(\Phi, f)$ 相对 $S$ 在维数 $\leqslant 1$ 时上同调真，只需证明 $(p_* p^{*} \Phi, f)$ 如此，$(q_* q^{*} \Phi, f)$ 相对 $S$ 在维数 $\leqslant 0$ 时上同调真，$(r_* r^{*} \Phi, f)$ 相对 $S$ 在维数 $\leqslant -1$ 时上同调真 (1.12 2)). 由上面 3) 1，可假设对每个有限常群层 $F$，$(q^{*} \Phi, fq)$，$(q^{*} \Phi, q)$，$(r^{*} \Phi, fr)$，$(r^{*} \Phi, r)$ 相对 $S$ 在维数 $\leqslant 0$ 时上同调真. 由 \ref{prop:XIII.1.6} 2) 得出 $(q_* q^{*} \Phi, f)$ 和 $(r_* r^{*} \Phi, f)$ 相对 $S$ 在维数 $\leqslant 0$ 时上同调真. 假设定理在光滑情形已证，$(p_* \Phi, fp)$ 和 $(p_* \Phi, p)$ 从而 $(p_* p^{*} \Phi, f)$ (1.6 2)) 相对 $S$ 在维数 $\leqslant 1$ 时上同调真. 这确实证明 $(\Phi, f)$ 相对 $S$ 在维数 $\leqslant 1$ 时上同调真.

而且由假设 $f_* p_* p^{*} \Phi$ 1-可构；由 \ref{thm:XIII.3.1} 可假设 $f_* q_* q^{*} \Phi$ 可构；由下面 \ref{lem:XIII.3.1.1} 得出 $f_* \Phi$ 1-可构.

3) 3. 归结到 $X$ 在 $S$ 上光滑仿射的情形.

由 3) 2，可假设 $X$ 在 $S$ 上光滑. 设 $(U_i)_{i \in I}$ 是 $X$ 的仿射开集的有限覆盖，$X_1$ 是诸 $U_i$ 的直和，$p \colon X_1 \to X$ 是典范态射. 由于 $p$ 是变动的概形上有限型平展层范畴的有效下降态射，如 3) 2 看到，若假设定理对诸 $U_i$，i.e. 对 $X_1$ 已证，则对 $X$ 也真.

3) 4. $X$ 在 $S$ 上光滑仿射的情形.

如 3) 1 看到，把 $S$ 限制到非空开集并作 $S$ 的平展满射扩张和根式满射扩张后，可以找到交换图式
\[
\begin{tikzcd}[row sep=large, column sep=large]
X \arrow[r,"{i}"] \arrow[dr,"{f}"] & P \arrow[d,"{q}"] \\
& Y
\end{tikzcd}
\]
其中 $P$ 在 $S$ 上光滑，$X$ 是 $P$ 中相对 $S$ 正规交叉除子的余集，$g$ 是紧合态射. 若 $F$ 是与 $S$ 的剩余特征互素的阶常层，$F$ 下的每个挠子相对 $S$ 在 $P$ 上温和分歧 (2.3 b)). $(F, f)$ 相对 $S$ 在维数 $\leqslant 1$ 时上同调真且 $f_* \Phi$ 可构的事实于是由 \ref{cor:XIII.2.7} 得出.

**4) 2) b) 的证明**

4) 1. $\Phi$ 是叠丛的情形.

可以找到有限型平展满射态射 $p \colon X_1 \to X$，使 $p^{*} \Phi$ 是平凡叠丛. 经下降，如 3) 2，看到只需对 $X_1$，$X_1 \times_X X_1$ 和 $X_1 \times_X X_1 \times_X X_1$ 证明定理. 归结到 $\Phi$ 是可构群层 $F$ (其纤维阶与 $S$ 的剩余特征互素) 下的挠子叠丛的情形.

由 SGA 4 IX 2.14，可以找到有限态射 $p_i \colon Z_i \to X$ 的有限族，以及对每个 $i$，与 $S$ 的剩余特征互素的阶有限常群层 $C_i$，使有单态射
\[
j \colon F \longrightarrow \prod_i p_{i*} C_i = G.
\]
设 $\Phi_i$ 是 $C_i$ 下的挠子叠，$\Psi$ 是 $G$ 下的挠子叠. 由 2) a)，把 $S$ 限制到非空开集后可假设诸 $(C_i, fp_i)$ 相对 $S$ 在维数 $\leqslant 1$ 时上同调真，诸叠 $f_* p_{i*} \Phi_i$ 1-可构. 由 \ref{cor:XIII.1.9} 得出诸 $(p_{i*} C_i, f)$ 相对 $S$ 在维数 $\leqslant 1$ 时上同调真；故 $(G, f)$ 也如此. 而且，由于 $p_{i*} \Phi_i$ 等价于群 $p_{i*} C_i$ 下的挠子叠 (SGA 4 VIII 5.8)，看到 $f_* \Psi$ 1-可构.

由于 $R^1 f_* G$ 可构，可以找到由有限型平展 $Y$-概形 $T$ 表示的层和满态射
\[
a \colon T \longrightarrow R^1 f_* G
\]
(SGA 4 IX 2.7)；且可假设 $T(T)$ 的恒等截面的像由 $X \times_Y T = X_T$ 上群 $G|_{X_T}$ 的挠子 $Q$ 定义. 设 $f_T \colon X_T \to T$ 是典范态射，$F_T = F|_{X_T}$，等等. 由 1) b) 可假设 (把 $S$ 限制到非空开集后) $(Q/F_T, f_T)$ 相对 $S$ 在维数 $\leqslant 0$ 时上同调真，且 $f_{T*}(Q/F_T)$ 可构. 由 \ref{lem:XIII.3.1.2} 得出 $f_* \Psi$ 可构.

证明 $(F, f)$ 相对 $S$ 在维数 $\leqslant 1$ 时上同调真. 由 \ref{prop:XIII.1.13} 2) 只需证明：对每个平展于 $Y$ 上的概形 $Y_1$ 和 $X_1 = X \times_Y Y_1$ 上的每个挠子 $Q_1$，若 $f_1 \colon X_1 \to Y_1$ 是典范态射，则 $(Q_1/F_1, f_1)$ 相对 $S$ 在维数 $\leqslant 0$ 时上同调真. 由 $T$ 的定义，$Q_1$ 局部于 $Y_1$ 的平展拓扑是 $Q$ 的逆像，这证明归结.

4) 2. 一般情形.

用\cref{lem:XIII.6.1.1}、4) 1 和 1) a) 看到，把 $S$ 限制到非空开集后可假设 $S(f_* \Phi)$ 可构，$(S_\Phi, f)$ 相对 $S$ 在维数 $\leqslant 0$ 时上同调真. 则可以找到由有限型平展 $Y$-概形 $T$ 表示的层和满态射
\[
a \colon T \longrightarrow S(f_* \Phi).
\]
沿用 4) 1 的记号，再令 $Z = T \times_Y T$，$X_Z = X \times_Y Z$，用 $f_Z \colon X_Z \to Z$ 记典范态射. 可假设 $T$ 选得使 $T(T)$ 的恒等截面经 $a$ 的像 $q$ 由 $(f_* \Phi)_T = \Phi_{X_T}$ 的对象 $p$ 定义. 设 $p_1$ 和 $p_2$ (resp. $q_1$ 和 $q_2$) 是 $p$ (resp. $q$) 经 $Z$ 到 $T$ 的两个投影的逆像. 可假设 (把 $S$ 限制到非空开集后) $(\mathrm{Aut}_{X_T}(p), f_T)$ 相对 $S$ 在维数 $\leqslant 1$ 时上同调真，$(\mathrm{Hom}_{X_Z}(p_1, p_2), f_Z)$ 相对 $S$ 在维数 $\leqslant 0$ 时上同调真，且 $f_{T*}(\mathrm{Aut}_{X_T}(p))$，$f_{Z*}(\mathrm{Hom}_{X_Z}(p_1, p_2))$ 可构 (a) 1) 和 4) 1)).

由此首先得出，对每个平展于 $Y$ 上的概形 $Y_1$ 和 $(f_* \Phi)_{Y_1}$ 的每对对象 $y$，$y_1$，层 $\mathrm{Aut}_{Y_1}(y)$ (resp. $\mathrm{Hom}_{Y_1}(y, y_1)$) 可构，这样的层局部于 $Y_1$ 的平展拓扑是 $f_{T*}(\mathrm{Aut}_{X_T}(p))$ (resp. $f_{Z*}(\mathrm{Hom}_{X_Z}(p_1, p_2))$) 的逆像.

剩下证明 $(\Phi, f)$ 相对 $S$ 在维数 $\leqslant 1$ 时上同调真. 只需证明：对每个 $S'$-概形 $S$ 和 $Y$ 的每个几何点 $y'$，若记 $Y$ 为 $Y$ 在 $y'$ 处的严格局部化，$X = X \times_Y Y$，等等，则典范函子
\[
\phi \colon \Phi(X) \longrightarrow \Phi'(X')
\]
是范畴等价.

证明 $\phi$ 忠实全. 设 $x$，$y \in \Phi(X)$，$x'$，$y'$ 是它们在 $\Phi'(X')$ 中的像，证明典范态射
\[
(\ast) \quad \mathrm{Hom}_X(x, y) \longrightarrow \mathrm{Hom}_{X'}(x', y')
\]
双射. 由 $T$ 的定义，存在两个 $Y$ 到 $T$ 的态射，使 $x$ 和 $y$ 是 $p$ 经这两个态射的逆像. 这等价于有态射 $Y \to Z$，从而有态射 $h \colon X \to X_Z$，使
\[
h^{*}(p_1) = x \quad h^{*}(p_2) = y.
\]
从而有典范同构
\[
\mathrm{Hom}_X(x, y) = h^{*}(\mathrm{Hom}_{X_Z}(p_1, p_2)).
\]
但考虑到 $(\mathrm{Hom}_{X_Z}(p_1, p_2), f_Z)$ 相对 $S$ 在维数 $\leqslant 0$ 时上同调真，看到 $(\mathrm{Hom}_X(x, y), f)$ 也如此，这证明态射 $(\ast)$ 双射.

证明 $\phi$ 本质满射. 设 $x' \in \Phi'(X')$. 由于 $(S_\Phi, f)$ 相对 $S$ 在维数 $\leqslant 0$ 时上同调真，典范态射
\[
H^0(X, S_\Phi) \longrightarrow H^0(X', S_{\Phi'})
\]
双射. 设 $G'$ 是 $x'$ 生成的 $\Phi'$ 的极大子叠丛；则存在 $\Phi$ 的极大子叠丛 $G$，使 $G$ 在 $X'$ 上的逆像是 $G'$. 由 4) 1，$(G, f)$ 相对 $S$ 在维数 $\leqslant 1$ 时上同调真；从而典范函子
\[
G(X) \longrightarrow G'(X')
\]
是范畴等价，这证明存在 $\Phi(X)$ 的元素 $x$，使 $x$ 在 $\Phi'(X')$ 中的像同构于 $x'$，完成定理的证明.
\end{proof}

\begin{lemma}\label{lem:XIII.3.1.1}
设 $S$ 是局部诺特概形，
\[
\Phi \xrightarrow{p} \Phi_1 \xrightarrow[p_2]{p_1} \Phi_2
\]
是 $S$ 上叠的正合图式 (1.10.1). 若 $\Phi_1$ 可构，$\Phi$ 也可构. 若对每个平展于 $S$ 上的概形 $S'$ 和 $(\Phi_1)_{S'}$ 的每对对象 $x_1$，$y_1$，层 $\mathrm{Hom}_{S'}(x_1, y_1)$ 可构，则对 $\Phi_{S'}$ 的每对对象 $x$，$y$，$\mathrm{Hom}_{S'}(x, y)$ 也可构. 假设 $\Phi_1$ 1-可构 (0) 且 $\Phi_2$ 可构，则 $\Phi$ 1-可构.
\end{lemma}

对每个平展于 $S$ 上的概形 $S'$ 和 $\Phi_{S'}$ 的每个对象 $x$，有单态射
\[
\mathrm{Aut}_{S'}(x) \longrightarrow \mathrm{Aut}_{S'}(p(x)).
\]
由 SGA 4 IX 2.9 得出，$\Phi_1$ 可构时 $\Phi$ 也可构，引理的第二个断言同样证明.

现在假设 $\Phi_1$ 1-可构，$\Phi_2$ 可构. 态射 $p$ 在极大子叠丛层上诱导态射
\[
\phi \colon S_\Phi \longrightarrow S_{\Phi_1}.
\]
设 $G$ 是 $S_\Phi$ 经 $\phi$ 的像；由 SGA 4 IX 2.9，$G$ 是可构层. 可以找到由有限型平展 $S$-概形 $T$ 表示的层和满态射
\[
a \colon T \longrightarrow G
\]
(SGA 4 IX 2.7). 且可以选 $T$，使 $T(T)$ 的恒等截面的像 $y$ 由 $(\Phi_1)_T$ 的形如 $x_1 = p(x)$ (其中 $x \in \Phi_T$) 的对象 $x_1$ 定义.

只需证明对 $S$ 的每个点 $s$，存在 $\overline{\{s\}}$ 的非空开集 $U$，使 $S_\Phi|_U$ 局部常可构. 设 $s \in S$，$s$ 是 $s$ 上方的几何点，$y_1, \ldots, y_n$ 是 $G_s$ 的元素. 由 $T$ 的定义存在态射 $h_i \colon s \to T$，使 $h^{*}_i(y) = y_i$. 设 $S'$ 是 $n$ 个同构于 $T$ 的概形在 $S$ 上的纤维积，$h \colon s \to S'$ 是诸 $h_i$ 的纤维积，$y_i$ (resp. $x_i$) 是 $y$ (resp. $x$) 经 $S'$ 到 $T$ 的第 $i$ 个投影的逆像. 设 $F_i$ 是 $S_\Phi|_{S'}$ 的 $y_i$ 的逆像子层，证明诸 $F_i$ 可构. 层 $F_i$ 是层 $F'_i$ 的商，后者对每个平展于 $S'$ 上方的概形 $S''$ 满足
\[
F'_i(S'') = \{\Phi_{S''} \text{ 的配备同构 } i \colon p(z) \simeq p(x_i|_{S''}) \text{ 的对象 } z \text{ 的模同构的类}\}.
\]
只需证明诸 $F'_i$ 可构. 若令 $z_i = p_1 p(x_i)$，$z'_i = p_2 p(x_i)$，有单态射
\[
\Psi_i \colon F'_i \longrightarrow \mathrm{Isom}_{S'}(z_i, z'_i),
\]
它把每个平展于 $S'$ 上的概形 $S''$ 和使有同构 $i \colon p(z) \simeq p(x_i|_{S''})$ 的 $\Phi_{S''}$ 的对象 $z$ 配给由下列图式交换条件定义的 $z_i$ 到 $z'_i$ 的同构
\[
\begin{tikzcd}[row sep=large, column sep=large]
p_1 p(z) \arrow[r,"{p_1(i)}"] \arrow[d,"{j}"] & z_i \arrow[d] \\
p_2 p(z) \arrow[r,"{p_2(i)}"'] & z'_i
\end{tikzcd}
\]
($j$ 是与正合图式关联的典范态射).

态射 $\Psi_i$ 单射，因为说两个使有同构 $i \colon p(z) \xrightarrow{\sim} p(x_i|_{S''})$，$i' \colon p(z') \xrightarrow{\sim} p(x_i|_{S''})$ 的对象 $z$，$z'$ 定义 $\mathrm{Isom}_{S''}(z_i, z'_i)$ 的同一个元素，等价于说 $p_1(i'^{-1} i) = p_2(i'^{-1} i)$，i.e. $i'^{-1} i$ 来自同构 $z \to z'$. 层 $F'_i$ 是 $\mathrm{Isom}_{S'}(z_i, z'_i)$ 的子层，故可构.

由于可以找到 $\overline{\{s\}}$ 的非空开集 $U$，使 $y_1|_U, \ldots, y_n|_U$ 生成 $G$，由下面\cref{lem:XIII.6.1.2} 得出 $S_\Phi|_U$ 可构，故把 $U$ 缩小后 $S_\Phi|_U$ 局部常可构.

\begin{lemma}\label{lem:XIII.3.1.2}
设 $S$ 是局部诺特概形，$f \colon X \to S$ 是态射，$F \to G$ 是 $X$ 上群层的单态射，$\Psi$ (resp. $\Psi_1$) 是 $F$ (resp. $G$) 下的挠子叠，$\Phi = f_* \Psi$，$\Phi_1 = f_* \Psi_1$. 假设给定 $S$ 上由有限型平展 $S$-概形 $T$ 表示的层和满态射
\[
a \colon T \longrightarrow S_{\Phi_1} \simeq R^1 f_* G,
\]
使存在 $X_T = X \times_S T$ 上群 $G|_{X_T}$ 的挠子 $Q$，在 $R^1 f_* G(T)$ 中定义 $T(T)$ 的恒等截面经 $a$ 的像. 设 $f_T \colon X_T \to T$ 是典范态射，令 $F_T = F|_{X_T}$. 假设 $\Phi_1$ 1-可构且 $f_{T*}(Q/F_T)$ 可构. 则 $\Phi$ 1-可构.
\end{lemma}

只需照抄 \ref{lem:XIII.3.1.1} 的证明，把存在态射 $\Psi_i$ 换成存在同构
\[
F'_i \xrightarrow{\sim} f_*(Q/F_T)|_{S'}.
\]

\begin{remark}\label{rem:XIII.3.1.3}
假设 $k = k(s)$ 特征零；则 $k$ 上维数 $\leqslant \dim X$ 的有限型概形强可消解奇点，\ref{thm:XIII.3.1} 的证明可以证明下列结果：
a) 存在 $S$ 的非空开集 $S_1$，使对 $S_1$ 上方的每个极大点特征零的概形 $S'$ 和 $X' = X \times_S S'$ 上每个局部常可构集合层 $F$，$(F, f')$ 相对 $S'$ 在维数 $\leqslant 0$ 时上同调真.
b) 若 $S$ 的所有剩余特征为零，存在 $S$ 的非空开集 $S_1$，使对 $S_1$ 上方的每个概形 $S'$ 和 $X'$ 上每个局部常可构群层 $F$，$(F, f')$ 相对 $S'$ 在维数 $\leqslant 1$ 时上同调真.

只需照抄 \ref{thm:XIII.3.1} 2) a) 的证明. 3) 4 中用的命题 \ref{cor:XIII.2.7} 适用于局部常层 $F$ 的情形，因为沿用 3) 4 的记号，由于 $S$ 的所有剩余特征为零，$F$ 下的每个挠子相对 $S$ 在 $P$ 上温和分歧.
\end{remark}

\begin{corollary}\label{cor:XIII.3.2}
设 $k$ 是特征 $p \geqslant 0$ 的域，$p'$ 是与 $p$ 不同的素数集，$f \colon X \to k$ 是相容态射.
\begin{enumerate}[1)]
\item 对每个集合层 $F$，$(F, f)$ 在维数 $\leqslant 0$ 时上同调真.
\item 假设满足下列条件之一：
\begin{enumerate}[a)]
\item $f$ 有限型，且 $k$ 的代数闭包上维数 $\leqslant \dim X$ 的有限型概形强可消解奇点.
\item $k$ 的代数闭包上有限型概形强可消解奇点.
\end{enumerate}
则对每个 ind-$p'$-群层 $F$，$(F, f)$ 在维数 $\leqslant 1$ 时上同调真.
\end{enumerate}
\end{corollary}

设 $F$ 是集合层 (resp. ind-$p'$-群层). 由 SGA 4 IX 2.7.2 可以把 $F$ 写成滤过归纳极限
\[
F = \varinjlim_i F_i,
\]
其中诸 $F_i$ 是可构集合层 (resp. ind-$p'$-群层). 由于 $f$ 相容，$f_*$ (resp. $R^1 f_*$) 与归纳极限交换 (SGA 4 VII 3.3). 若知道诸 $(F_i, f)$ 在维数 $\leqslant 0$ (resp. 每个集合层在维数 $\leqslant 0$ 且诸 $(F_i, f)$ 在维数 $\leqslant 1$) 时上同调真，则 $(F, f)$ 也如此. 可假设 $F$ 可构.

若假设 $f$ 有限型 (resp. 满足 a))，命题由 \ref{thm:XIII.3.1} 1) b) (resp. \ref{thm:XIII.3.1} 2) b)) 得出. 现在证明不再假设 $f$ 有限型时的 3.2. 对每个 $k$ 上方的概形 $S'$ 和 $S'$ 的每个几何点 $s$，记 $k$ (resp. $S'$) 为 $k$ 在 $s$ 处 (resp. $S'$ 在 $s$ 处) 的严格局部化，$X$ 为 $X$ 在 $k$ 上的逆像，考虑笛卡儿方形
\[
\begin{tikzcd}[row sep=large, column sep=large]
X \arrow[d,"{f}"] & X' \arrow[l,"g"] \arrow[d,"{f'}"] \\
k & S' \arrow[l]
\end{tikzcd}
\]
只需证明对每个 $S'$ 和每个 $s$，典范态射
\[
H^0(X, F) \longrightarrow H^0(X', F') \quad \text{(resp. } H^1(X, F) \longrightarrow H^1(X', F')\text{)}
\]
是同构. 只需证明有关系
\[
(\ast) \quad F \simeq g_* g^{*} F \quad \text{(resp. } R^1 g_*(g^{*} F) = 0\text{)}.
\]
这种形式的问题局部于 $X$ 的平展拓扑，可假设 $X$ 仿射；经极限过渡可假设 $X$ 在 $k$ 上有限型. 已知 $(F, f)$ 在维数 $\leqslant 0$ (resp. $\leqslant 1$) 时上同调真，把 $X$ 换成平展于 $X$ 上的有限型概形时也如此，这证明 $(\ast)$.

\begin{theorem}\label{thm:XIII.3.3}
设 $S$ 是生成点为 $s$ 的不可约概形，$f \colon X \to S$ 是有限表现态射. 假设 $k$ 的代数闭包 $k$ 上维数 $\leqslant \dim X_s$ 的有限型概形可消解奇点 (EGA IV 7.9.1). 则若 $L$ 表示与 $S$ 的剩余特征不同的素数集，可以找到 $S$ 的非空开集 $S_1$，使态射 $f|_{S_1}$ 对 $L$ 万有局部 1-非球.
\end{theorem}

可假设 $S$ 整且 $X$ 既约 (SGA 4 VIII 1.1). 经极限过渡可假设 $S$ 诺特. 而且只需在有限扩张 $S_1 \to S$ 后证明定理，其中 $S_1$ 是整概形，$S_1 \to S$ 由平展扩张和根式满射扩张复合而成.

先证明，把 $S$ 限制到非空开集后 $f$ 万有局部 0-无环. 作 $k(s)$ 的根式扩张后，可假设态射 $(X_s)_{\mathrm{r\acute{e}d}} \to s$ 可分；故可假设 (把 $S$ 限制到非空开集并作 $S$ 的根式满射扩张后) 态射 $f$ 平坦、几何纤维可分 (EGA IV 12.1.1)，这蕴含 $f$ 万有 0-无环 (SGA 4 XV 4.1).

证明，把 $S$ 限制到非空开集后 $f$ 对 $L$ 万有局部 1-非球. 作 $k(s)$ 的有限扩张 (可把它看作平展扩张与根式扩张的复合) 后，可以找到紧合满射态射 $p_s \colon Y_s \to X_s$，其中 $Y_s$ 是与 $X_s$ 同维数、在 $S$ 上光滑的概形，且可假设 $p_s$ 来自紧合满射态射 $p \colon Y \to X$，其中 $Y$ 在 $S$ 上光滑 (EGA IV 9.6.1 和 12.1.6). 只需证明，把 $S$ 限制到非空开集后，对每个方形笛卡儿的图式
\[
\begin{tikzcd}[row sep=large, column sep=large]
S'' \arrow[d,"{i}"] & X'' \arrow[l,"{f''}"] \arrow[d,"{j}"] \\
S' \arrow[d] & X' \arrow[l,"{f'}"] \arrow[d] \\
S & X \arrow[l,"{f}"]
\end{tikzcd}
\]
其中 $i$ 有限表现平展，和 $S''$ 上每个 ind-$L$-群层 $F$，若 $\Phi$ 是 $F$ 下的挠子叠，典范态射
\[
f'^{*} i_* \Phi \longrightarrow j_* f''^{*} \Phi
\]
是等价. 设 $Z = Y \times_X Y$，$T = Y \times_X Y \times_X Y$. 自然地有交换图式
\[
\begin{tikzcd}[row sep=large, column sep=large]
S'' \arrow[d,"{i}"] & X'' \arrow[l,"{f''}"] \arrow[d,"{j}"] & Y'' \arrow[l,"{p''}"] \arrow[d,"{p'}"] & Z'' \arrow[l] \arrow[d] & T'' \arrow[l] \arrow[d] \\
S' \arrow[d] & X' \arrow[l,"{f'}"] \arrow[d,"{f}"] & Y' \arrow[l,"{p'}"] \arrow[d,"{p}"] & Z' \arrow[l] \arrow[d] & T' \arrow[l] \arrow[d] \\
S & X \arrow[l] & Y \arrow[l] & Z \arrow[l] & T
\end{tikzcd}
\]
设 $q \colon Z \to X$ 和 $r \colon T \to X$ 是典范态射，$q'$，$r'$，$q''$，$r''$ 的记号显然. 由 \ref{exa:XIII.1.11} 2)，有行正合的本质上交换图式：
\[
\begin{tikzcd}[row sep=large, column sep=large]
f'^{*} i_* \Phi \arrow[r] \arrow[d,"{a}"] & p'_* p'^{*}(f'^{*} i_* \Phi) \arrow[r, Rightarrow] \arrow[d,"{b}"] & q'_* q'^{*}(f'^{*} i_* \Phi) \arrow[r, Rightarrow] \arrow[d,"{c}"] & r'_* r'^{*}(f'^{*} i_* \Phi) \arrow[d,"{d}"] \\
j_* f''^{*} \Phi \arrow[r] & j_* p''_* p''^{*}(f''^{*} \Phi) \arrow[r, Rightarrow] & j_* q''_* q''^{*}(f''^{*} \Phi) \arrow[r, Rightarrow] & j_* r''_* r''^{*}(f''^{*} \Phi)
\end{tikzcd}
\]
由于 $Y$ 在 $S$ 上光滑，态射 $fp$ 对 $L$ 万有局部 1-非球 (SGA 4 XV 2.1)，由 \cite[VII 2.1.7]{exp13-2} 得出 $b$ 是范畴等价. 另一方面，把 $S$ 限制到非空开集后可假设态射 $Z \to S$ 和 $T \to S$ 万有局部 0-无环. 由此得出函子 $c$ 和 $d$ 忠实全，上面图式于是证明 $a$ 是等价，这完成证明.

\begin{corollary}\label{cor:XIII.3.4}
设 $k$ 是特征 $p \geqslant 0$ 的域，$p'$ 是与 $p$ 不同的素数集，$f \colon X \to k$ 是相容态射. 假设满足下列条件之一：
a) $f$ 有限型，且 $k$ 的代数闭包上维数 $\leqslant \dim X$ 的有限型概形可消解奇点.
b) $k$ 的代数闭包上有限型概形可消解奇点.
则 $f$ 对 $p'$ 万有局部 1-非球.
\end{corollary}

情形 a) 由 \ref{thm:XIII.3.3} 得出. 在情形 b)，问题局部于 $X$，可假设 $X$ 仿射；经极限过渡 (SGA 4 XV 1.3) 归结到 $X$ 在 $k$ 上有限型的情形.

\begin{corollary}\label{cor:XIII.3.5}
设 $S$ 是生成点为 $s$ 的不可约概形，$f \colon X \to S$ 是有限表现态射. 假设 $k(s)$ 的代数闭包 $k$ 上维数 $\leqslant \dim X_s$ 的有限型概形强可消解奇点 (SGA 5 I 3.1.5). 若 $L$ 表示与 $S$ 的剩余特征不同的素数集，可以找到 $S$ 的非空开集 $S_1$，使对 $S_1$ 的几何点的每个特殊化 $s_1 \to s_2$，特殊化态射 (2.10)
\[
\pi^L_1(X_{s_1}) \longrightarrow \pi^L_1(X_{s_2})
\]
双射.
\end{corollary}

由 \ref{thm:XIII.3.1} 和 \ref{thm:XIII.3.3}，可 (把 $S$ 限制到非空开集后) 假设 $f$ 对 $L$ 局部 1-非球，且对 $X$ 上每个有限常 $L$-群层 $F$，$(F, f)$ 在维数 $\leqslant 1$ 时上同调真. 由 \ref{prop:XIII.1.14} 得出，对 $S_1$ 的几何点的每个特殊化 $s_1 \to s_2$，特殊化态射
\[
(R^1 f_* F)_{s_2} \longrightarrow (R^1 f_* F)_{s_1}
\]
双射. 该推论正是前面结果用基本群语言的翻译.

\section{同伦正合序列}

\refstepcounter{theorem}\addtocounter{theorem}{-1}\noindent\textbf{\thesection.0.}\quad 设 $X$ 和 $S$ 是两个连通概形，$f \colon X \to S$ 是态射，$a$ 是 $X$ 的几何点，$L$ 是素数集. 设 $K$ 是典范同态 $\pi_1(X, a) \to \pi_1(S, a)$ 的核，$N$ 是使 $K/N$ 是 pro-$L$-群 $K^L$ 的 $K$ 的最小 pro-正规子群. 则 $N$ 在 $\pi_1(X, a)$ 中正规，记
\[
\pi'_1(X, a)
\]
为 $\pi_1(X, a)$ 除以 $N$ 的商. 若 $a$ 是几何纤维 $X_s$ 的几何点，典范态射
\[
\pi_1(X_s, a) \longrightarrow \pi_1(X, a) \longrightarrow \pi_1(S, a)
\]
可以定义典范态射
\[
\pi^L_1(X_s, a) \xrightarrow{u} \pi'_1(X, a)\index[notations]{$\pi_1'(X, a)$} \xrightarrow{v} \pi_1(S, a)
\]
有 $vu = 0$.

\begin{proposition}\label{prop:XIII.4.1}
设 $S$ 是连通概形，$f \colon X \to S$ 是局部 0-无环 (SGA 4 XV 1.11) 的态射；再假设 $f$ 0-无环 (当 $f$ 相容时，这等价于 $f$ 的几何纤维连通 (SGA 4 XV 1.16)). 设 $L$ 是素数集. 若 $S'$ 是平展于 $S$ 上的概形，记 $X'$，$f'$ 为 $X$，$f$ 在 $S'$ 上的逆像. 假设对 $S$ 的每个平展覆盖 $S'$ 和 $X'$ 的每个平展覆盖 $E$ (它是群为 $L$-群的 Galois 覆盖的商)，$(E, f')$ 相对 $S'$ 在维数 $\leqslant 0$ 时上同调真，且 $f'_* E$ 可构. 则若 $s$ 是 $S$ 的几何点，$a$ 是纤维 $X_s$ 的几何点，群同态序列
\[
(4.1.1) \quad \pi^L_1(X_s, a) \xrightarrow{u} \pi'_1(X, a) \xrightarrow{v} \pi_1(S, a) \longrightarrow 1
\]
正合.
\end{proposition}

这个断言推广 X.1.4，将照抄它的证明.

先证明 $v$ 满射. 只需证明对 $S$ 的每个连通平展覆盖 $S'$，$X'$ 也连通 (\cref{prop:V.6.9}). 设 $C$ 是至少有两个元素的集合. 由 $f$ 0-无环的事实得出典范态射
\[
H^0(S', C_{S'}) \longrightarrow H^0(X', C_{X'})
\]
双射，故 $X'$ 连通，从而 $v$ 满射.

由 $K^L$ 的定义 (4.0)，有正合序列
\[
1 \longrightarrow K^L \longrightarrow \pi'_1(X, a) \longrightarrow \pi_1(S, a) \longrightarrow 1
\]
设 $\widetilde{S}$ 是 $S$ 的万有覆盖，$\widetilde{X} = \widetilde{S} \times_S X$；群 $K^L$ 分类群为 $L$-群的 Galois 覆盖 $P$，使存在 $S$ 的平展覆盖 $S'$ 和 $X' = X \times_S S'$ 的 Galois 覆盖 $Q$，使有同构 $P \simeq Q \times_{X'} \widetilde{X}$. 为使序列 (4.1.1) 正合，必要且充分条件是典范态射
\[
\pi^L_1(X_s, a) \longrightarrow K^L
\]
满射. 由 $K^L$ 的解释，这等价于：对 $S$ 的每个平展覆盖 $S'$ 和 $X'$ 的每个群为 $L$-群的 Galois 覆盖 $Q$，若 $P = Q \times_{X'} \widetilde{X}$ 连通，则 $Q|_{X_s}$ 连通. 证明这最后一个条件满足. 设 $S'$ 是 $S$ 的平展覆盖，$Q$ 是 $X'$ 的群为 $L$-群 $F$ 的 Galois 覆盖，使 $Q|_{X_s}$ 不连通，证明 (把 $S'$ 换成平展覆盖后) $Q$ 变成不连通. 存在 $F$ 的与 $F$ 不同的子群 $G$ 和 $G|_{X_s}$ 下的挠子 $R$，使 $Q|_{X_s}$ 由 $R$ 经结构群扩张 $G \to F$ 得到. $X'$ 的平展覆盖 $E = Q/G$ 使 $E|_{X_s}$ 有截面. 由 \ref{cor:XIII.1.16}，$f_* E$ 局部常可构，把 $S'$ 换成平展覆盖后甚至可假设 $f_* E$ 常. 由于 $(E, f')$ 相对 $S'$ 在维数 $\leqslant 0$ 时上同调真，且 $H^0(X_s, E|_{X_s})$ 非空，看到 $E$ 有截面. 但这证明 $Q$ 不连通，完成证明.

由 \ref{prop:XIII.4.1} 推出下列将在 \ref{prop:XIII.4.6} 使用的引理.

\begin{lemma}\label{lem:XIII.4.2}
设 $S$ 是连通概形，$f \colon X \to S$ 是局部 0-无环且 0-无环的态射，$L$ 是素数集. 假设对 $X$ 上每个有限常 $L$-群层 $F$，$(F, f)$ 相对 $S$ 在维数 $\leqslant 1$ 时上同调真，且对 $S$ 的每个平展覆盖 $S'$ 和 $X'$ 的每个平展覆盖 $E$ (它是群为 $L$-群的 Galois 覆盖的商)，$f'_* E$ 可构. 则若 $s$ 是 $S$ 的几何点，$a$ 是纤维 $X_s$ 的几何点，群同态序列
\[
\pi^L_1(X_s, a) \longrightarrow \pi'_1(X, a) \longrightarrow \pi_1(S, a) \longrightarrow 1
\]
正合.
\end{lemma}

\ref{prop:XIII.4.1} 的假设满足. 确实由 \ref{prop:XIII.1.13} 3) 得出，对每个平展于 $S$ 上的概形 $S'$ 和 $X'$ 的每个平展覆盖 $E$ (它是群为 $L$-群的 Galois 覆盖的商)，$(E, f')$ 相对 $S'$ 在维数 $\leqslant 0$ 时上同调真.

\begin{proposition}\label{prop:XIII.4.3}
设 $S$ 是连通概形，$L$ 是素数集，$f \colon X \to S$ 是 0-无环、对 $L$ 局部 1-非球 (SGA 4 XV 1.11) 的态射，$g \colon S \to X$ 是 $f$ 的截面. 设 $s$ 是 $S$ 的几何点，$a$ 是纤维 $X_s$ 的几何点. 假设对每个有限常 $L$-群层 $F$，$(F, f)$ 在维数 $\leqslant 1$ 时上同调真，$F$ 下的挠子叠经 $f$ 的直像是 \emph{1-可构叠} (0)，且对 $X'$ 的每个平展覆盖 $E$ (它是群为 $L$-群的 Galois 覆盖的商)，$f'_* E$ 可构. 则群同态序列
\[
(4.3.1) \quad 1 \longrightarrow \pi^L_1(X_s, a) \xrightarrow{u} \pi'_1(X, a) \xrightarrow{v} \pi_1(S, a) \longrightarrow 1
\]
正合.
\end{proposition}

考虑到 \ref{lem:XIII.4.2}，只需证明态射 $u$ 单射，i.e. 证明：对 $X_s$ 的每个群为 $L$-群 $C$ 的主覆盖 $Z$，存在 $X$ 的平展覆盖 $\widetilde{Z}$ 和 $Z|_{X_s}$ 的某连通分量到 $Z$ 的态射 (\cref{cor:V.6.8}). 设 $Z$ 是 $X_s$ 的群为 $L$-群 $C$ 的主覆盖，$z$ 是它在 $H^1(X_s, C_{X_s})$ 中的类. 由 \ref{rem:XIII.1.5} d)，有典范同构
\[
R^1 f_* C_X(s) \xrightarrow{\sim} H^1(X_s, C_{X_s}),
\]
且由 \ref{cor:XIII.1.16}，$R^1 f_* C_X$ 是局部常可构层. 可以找到 $S$ 的平展覆盖 $S'$，使 $R^1 f_* C_X|_{S'}$ 常. 若 $s \to S'$ 是位于几何点 $s \to S$ 上方的几何点，存在 $H^0(S', R^1 f_* C_X)$ 的元素 $z$，使它在 $H^1(X_s, C_{X_s})$ 中的像是 $z$. 由下面\cref{lem:XIII.4.3.1}，可以找到 $S'$ 的平展覆盖 $S'_1$ 和 $X'_1$ 上群 $C$ 的挠子 $P$，使它在 $H^0(S'_1, R^1 f_* C)$ 中的像等于 $z$ 的限制. 挠子 $P$ 可由 $X'_1 = X \times_S S'_1$ 的平展覆盖 $Z$ 表示，使 $Z \times_{X'_1} X_s$ 同构于 $Z$. 把 $Z$ 看作 $X$ 的平展覆盖，则有 $Z \times_X X_s$ 到 $Z$ 的态射，完成证明.

\begin{lemma}\label{lem:XIII.4.3.1}
设 $f \colon X \to S$ 是 0-无环且局部 0-无环的态射，$g$ 是 $f$ 的截面. 设 $C$ 是使 $(C_X, f)$ 在维数 $\leqslant 0$ 时上同调真、$C_X$ 下的挠子叠经 $f$ 的直像可构的有限常群. 则对 $H^0(S, R^1 f_* C_X)$ 的每个截面 $z$，可以找到 $S$ 的平展覆盖 $S_1$，且若 $X_1 = X \times_S S_1$，存在 $H^1(X_1, C_{X_1})$ 的元素，使它在典范态射
\[
H^1(X_1, C_{X_1}) \longrightarrow H^0(S_1, R^1 f_* C_{X_1})
\]
下的像等于 $z$ 在 $H^0(S_1, R^1 f_* C_{X_1})$ 中的限制.
\end{lemma}

对每个平展于 $S$ 上的概形 $S'$，令 $X' = X \times_S S'$，用 $g'$ (resp. $F'$，等等) 记 $g$ (resp. $F$，等等) 经态射 $S' \to S$ 的逆像. 由
\[
G(S') = \{X' \text{ 上群 } C_{X'} \text{ 的配备同构 } g'^{*} P \xrightarrow{\sim} C_{S'} \text{ 的挠子 } P \text{ 的模同构的类}\}
\]
定义的 $S$ 上预层 $G$ 是层. 这由下降得出，因为 $X'$ 上挠子 $P$ 的同构由它在 $g'(S')$ 上的限制确定. 而且有满态射
\[
G \longrightarrow R^1 f_* F.
\]
设 $z$ 是 $H^0(S, R^1 f_* C_X)$ 的元素，$H$ 是 $G$ 的 $z$ 的逆像子层. 只需证明 $H$ 是局部常可构层. 该性质局部于 $S$，可假设 $z$ 来自 $H^1(X, C_X)$ 的由使 $g^{*} P$ 同构于 $C_S$ 的挠子 $P$ 表示的元素. 给出同构 $i \colon g^{*} P \xrightarrow{\sim} C_S$ 等价于给出 $\mathrm{Aut}_{C_S}(g^{*} P)$ 的全局截面，且两个同构 $i$ 和 $i'$ 定义 $G(X)$ 的同一个元素，当且仅当 $ii'^{-1}$ 是 $\mathrm{Aut}_{C_X}(P)$ 的某个元素的像. 若考虑典范单射
\[
f_* \mathrm{Aut}_{C_X}(P) \longrightarrow \mathrm{Aut}_{C_S}(g^{*} P) \simeq C,
\]
$H$ 等同于 $\mathrm{Aut}_{C_S}(g^{*} P)$ 除以 $f_* \mathrm{Aut}_{C_X}(P)$ 的商. 由 \ref{cor:XIII.1.16}，$f_* \mathrm{Aut}_{C_X}(P)$ 局部常；故 $H$ 也局部常，完成证明.

\begin{example}\label{exa:XIII.4.4}
注意，若 $S$ 是连通概形，态射 $f$ 紧合平坦有限表现、几何纤维可分连通时 \ref{prop:XIII.4.1} 的假设满足，$L$ 任意 (参见 X 1.3). 若 $f$ 还光滑且有截面，$L$ 为与 $S$ 的剩余特征不同的素数集 (SGA 4 XV 2.1 和 XVI 5.2) 时，\ref{prop:XIII.4.3} 的假设满足.

若 $S$ 连通，有紧合有限表现概形 $Z$ 在 $S$ 上平坦、几何纤维可分连通，$X$ 是 $Z$ 中相对 $S$ 正规交叉除子的余集，$L$ 为与 $S$ 的剩余特征不同的素数集 (\cref{cor:XIII.2.9}) 时，\ref{prop:XIII.4.1} 的假设也满足. 若 $f$ 还光滑且有截面，\ref{prop:XIII.4.3} 的假设满足.
\end{example}

\refstepcounter{theorem}\noindent\textbf{\thesection.\arabic{theorem}.}\quad 沿用 \ref{prop:XIII.4.3} 的记号和假设. 若 $s$ 是 $S$ 的几何点，$a = g(s)$，截面 $g$ 可以定义态射
\[
w \colon \pi_1(S, a) \longrightarrow \pi'_1(X, a),
\]
使 $\pi'_1(X, a)$ 等同于 $\pi_1(S, a)$ 被 $\pi_1(X_s, a)$ 的半直积. pro-有限群 $\pi_1(S, a)$ 于是作用在 $\pi_1(X_s, a)$ 上. 由于 $\pi^L_1(X_s, a)$ 是 $\pi_1(S, a)$ 的作用不变群的严格投射极限，配备该作用的 $\pi^L_1(X_s, a)$ 的给出等价于 $S$ 上有限平展群概形的严格投射系统的给出，记为
\[
\pi^L_1(X/S, g, s) \quad \text{或简记}\index[notations]{$\pi_1^L(X/S, g, s)$} \quad \pi^L_1(X/S, g).
\]
有下列性质：

4.5.1. 对 $S$ 上每个纤维为 $L$-群的有限平展群概形 $G$，配备同构 $g^{*} P \xrightarrow{\sim} G$ 的 $G$ 在 $X$ 上的逆像 $G_X$ 下的挠子 $P$ 的类的集 $E$，典范同构于集
\[
\mathrm{Hom}_S(\pi^L_1(X/S, g, s), G) \quad \text{模 } G \text{ 的内自同构}.
\]

4.5.2. 对 $S$ 上每个纤维为 $L$-群的有限平展群概形 $G$，层 $R^1 f_* G_X$ 典范同构于与预层
\[
S' \longmapsto \mathrm{Hom}_{S'}(\pi^L_1(X/S, g, s), G) \text{ 模 } G \text{ 的内自同构}
\]
关联的层. ($S'$ 表示平展于 $S$ 上的概形).

4.5.3. 设 $S'$ 是连通 $S$-概形，$s$ 是 $S'$ 的几何点，$X'$，$g'$ 是 $X$，$g$ 在 $S'$ 上的逆像. 则 $\pi^L_1(X'/S', g', s)$ 典范同构于 $\pi^L_1(X/S, g, s)$ 在 $S'$ 上的逆像. 对 $S$ 的每个几何点 $\xi$，纤维 $\pi^L_1(X/S, g, s)_\xi$ 同构于 $\pi^L_1(X_\xi)$.

$G$ 的给出确实等价于配备 $\pi_1(S, a)$ 的作用的抽象 $L$-群 $G$ 的给出，从而有 $\pi'_1(X, a)$ 在 $G$ 上的作用. 4.5.1 中定义的同构由典范态射
\[
H^1(\pi'_1(X, a), G) \longrightarrow H^1(\pi^L_1(X_s, a), G) = \mathrm{Hom}(\pi^L_1(X_s, a), G)/\mathrm{aut.\,int.}\, G
\]
限制到子集 $E$ 得到，$E$ 双射到与 $\pi_1(S, a)$ 的作用相容的 $\pi^L_1(X_s, a)$ 到 $G$ 的态射子集. 断言 4.5.3 由 $\pi^L_1(X/S, g, s)$ 的定义，考虑到同伦正合序列 (\cref{lem:XIII.4.3.1}) 得出，4.5.2 由 4.5.1 和 4.5.3 推出.

\begin{proposition}\label{prop:XIII.4.6}
[Künneth 公式] 设 $k$ 是特征 $p \geqslant 0$ 的可分闭域，$X$ 和 $Y$ 是两个连通 $k$-概形，$a$ 是 $X$ 的几何点，$b$ 是 $Y$ 的几何点，$c$ 是 $X \times_k Y$ 的位于 $a$ 和 $b$ 上方的几何点. 假设满足下列条件之一：
a) $X$ 在 $k$ 上有限型，且 $k$ 的代数闭包 $k$ 上维数 $\leqslant \dim X$ 的有限型概形强可消解奇点 (SGA 5 I 3.1.5).
b) $X$ 拟紧拟分离，且 $k$ 上每个有限型概形强可消解奇点.
则若 $p'$ 是与 $p$ 不同的素数集，态射
\[
(4.6.0) \quad \pi^{p'}_1(X \times_k Y, c) \longrightarrow \pi^{p'}_1(X, a) \times \pi^{p'}_1(Y, b),
\]
由与投影
\[
X \times_k Y \longrightarrow X \quad \text{和} \quad X \times_k Y \longrightarrow Y
\]
关联的基本群同态得出，是同构.
\end{proposition}

可假设 $k$ 代数闭且 $X$ 既约 (SGA 4 VIII 1.1). 设 $Z = X \times_k Y$，$g \colon X \to k$ 和 $f \colon Z \to Y$ 是典范态射. 由 \ref{cor:XIII.3.5}，态射 $g$ 从而 $f$ 对 $p'$ 万有局部 1-非球. 由于 $X$ 连通，$f$ 0-无环 (SGA 4 XV 1.16). 另一方面由 \ref{cor:XIII.3.2} 得出，对每个 $p'$-有限群 $C$，$(C_X, g)$ 相对 $k$ 在维数 $\leqslant 1$ 时上同调真. 由此得出 $(C_Z, f)$ 相对 $Y$ 在维数 $\leqslant 1$ 时上同调真 (1.5 c))，且 $f_* C_Z$ 和 $R^1 f_* C_Z$ 是常层. 从而 $g$ 满足 \ref{lem:XIII.4.2} 的所有假设. 有正合序列
\[
\pi^{p'}_1(X_b, c) \longrightarrow \pi^{p'}_1(Z, c) \longrightarrow \pi^{p'}_1(Y, b) \longrightarrow 1.
\]
而且复合态射
\[
\pi^{p'}_1(X_b, c) \longrightarrow \pi^{p'}_1(Z, c) \longrightarrow \pi^{p'}_1(Y, b)
\]
是同构，故有正合序列
\[
1 \longrightarrow \pi^{p'}_1(X, a) \longrightarrow \pi^{p'}_1(Z, c) \longrightarrow \pi^{p'}_1(Y, b) \longrightarrow 1.
\]
另一方面态射 (4.6.0) 可以定义该正合序列到正合序列
\[
1 \longrightarrow \pi^{p'}_1(X, a) \longrightarrow \pi^{p'}_1(X, a) \times \pi^{p'}_1(Y, b) \longrightarrow \pi^{p'}_1(Y, b) \longrightarrow 1
\]
的态射，由此得出态射 (4.6.0) 是同构.

\refstepcounter{theorem}\noindent\textbf{\thesection.\arabic{theorem}.}\quad 设
\[
\begin{tikzcd}[row sep=large, column sep=large]
X \arrow[d,"{f}"] & X_U \arrow[l] \arrow[d,"{f_U}"] & X_s \arrow[l] \arrow[d] \\
S & U \arrow[l] & s \arrow[l]
\end{tikzcd}
\]
是所有方形笛卡儿的图式，其中 $S$ 是弧连通概形 (SGA 4 IX 2.12)，$U$ 是 $S$ 的连通开集，$s$ 是 $U$ 的几何点. 设 $a$ 是 $X_s$ 的几何点，$L$ 是素数集. 设 $g$ 是 $f$ 的截面，假设满足下列条件：
a) 态射 $f$ 0-无环、局部 0-无环，且对 $S$ 的每个平展覆盖 $S'$ 和 $X \times_S S'$ 的每个平展覆盖 $E$ (它是群为 $L$-群的 Galois 覆盖的商)，$(E, f_{(S')})$ 相对 $S'$ 在维数 $\leqslant 0$ 时上同调真.
b) 态射 $f_U$ 对 $L$ 局部 1-非球，且对 $X_U$ 上每个有限常 $L$-群层 $F$，$(F, f_U)$ 在维数 $\leqslant 1$ 时上同调真，$R^1 f_{U*} F$ 的纤维有限.
由 \ref{prop:XIII.4.1} 和 \ref{prop:XIII.4.3} 得出下列行正合的交换图式：
\begin{equation}
\begin{tikzcd}[row sep=large, column sep=large]
1 \arrow[r] & \pi^L_1(X_s, a) \arrow[r] \arrow[d] & \pi'_1(X_U, a) \arrow[r] \arrow[d] & \pi_1(U, a) \arrow[r] \arrow[d] & 1 \\
& \pi^L_1(X_s, a) \arrow[r] & \pi'_1(X, a) \arrow[r] & \pi_1(S, a) \arrow[r] & 1
\end{tikzcd}
\tag{4.7.0}
\label{eq:XIII.4.7.0}
\end{equation}
借助截面 $g$，有态射
\[
\pi_1(U, a) \longrightarrow \pi'_1(X_U, a), \quad \pi_1(S, a) \longrightarrow \pi'_1(X, a)
\]
由此得到 $\pi_1(S, a)$ 和 $\pi'_1(X_U, a)$ 在 $\pi_1(U, a)$ 上方的余直和到 $\pi'_1(X, a)$ 的态射：
\[
(4.7.1) \quad \phi \colon \pi = \pi_1(S, a) \amalg_{\pi_1(U, a)} \pi'_1(X_U, a) \longrightarrow \pi'_1(X, a).
\]
假设满足下列条件：
c) 若 $T = S - U$，有 $\mathrm{prof}^{\mathrm{ét}}_T(S) \geqslant 2$ (SGA 2 XIV 1.1).
则把 $S$ 的平展覆盖配给它在 $U$ 上的限制的函子忠实全 (SGA 2 XVI 1.4). 由此得出态射 $\pi_1(U, a) \to \pi_1(S, a)$ 满射 (\cref{prop:V.6.9})，由 \eqref{eq:XIII.4.7.0} 图式得出态射 $\pi'_1(X_U, a) \to \pi'_1(X, a)$ 也满射；更不用说 $\phi$ 是满态射. 设
\[
(4.7.2) \quad K = \mathrm{Ker}(\pi_1(U, a) \longrightarrow \pi_1(S, a)).
\]
(4.7.1) 中的群 $\pi$ 等同于 $\pi'_1(X_U, a)$ 除以 $K$ 的像 $L$ 在 $\pi'_1(X_U, a)$ 中生成的闭不变子群的商. 把 $\pi'_1(X_U, a)$ 看作 $\pi_1(U, a)$ 被 $\pi^L_1(X_s, a)$ 的半直积. 群 $K$ 于是经内自同构作用在 $\pi^L_1(X_s, a)$ 上，商 $\pi = \pi'_1(X_U, a)/L$ 等同于 $\pi_1(U, a)/K = \pi_1(S, a)$ 被 $\pi^L_1(X_s, a)$ 在 $K$ 下的余不变量群 $\pi^L_1(X_s, a)_K$ 的半直积. 最终有满态射
\[
(4.7.3) \quad \phi \colon \pi = \pi^L_1(X_s, a)_K\index[notations]{$\pi_1^L(X_s, a)_K$} \cdot \pi_1(S, a) \longrightarrow \pi'_1(X, a).
\]

下列命题给出 $\phi$ 是同构的条件.

\begin{proposition}\label{prop:XIII.4.7.4}
记号同 4.7. 假设除条件 a)，b)，c) 外还满足下列条件：
d) 对 $T = S - U$ 的每个点 $t$，态射 $f$ 在 $g(t)$ 处对 $L$ 局部 1-非球.
e) 对每个点 $t \in T$，纤维 $X_t$ 的每个不可约分量含 $g(t)$，且对 $X_t - \{g(t)\}$ 的每个非极大点 $x$，有
\[
\mathrm{prof}^{\mathrm{hop}}_x(X) \geqslant 3 \quad \text{(SGA 2 XIV 1.2)}
\]
且环 $\mathcal{O}_{X, x}$ 诺特.
则态射 (4.7.3) 是同构.
\end{proposition}

如前所说，群 $\pi$ 等同于 $\pi'_1(X_U, a)$ 除以 $K$ (4.7.2) 的像在 $\pi'_1(X_U, a)$ 中生成的闭不变子群 $L$ 的商. 这等价于说 $\pi$ 分类 $X_U$ 的主覆盖 $Z$ (使 $g^{*}_U(Z)$ 延拓成 $S$ 的平展覆盖，且 $Z$ 在 $X_s$ 上诱导的覆盖由群为 $L$-群的主覆盖经结构群扩张得到). 为证明 $\phi$ 是同构，只需证明这样的覆盖 $Z$ 延拓到整个 $X$.

先证明 $Z$ 延拓成含 $X_U$ 和 $g(S)$ 的开集的平展覆盖. 设 $W$ 是平展于 $X$ 上的概形，它的像含 $X_U$ 和 $g(S)$，令 $W_U = W \times_S U$. 由态射 $W \to S$ 0-无环和关系 $\mathrm{prof}^{\mathrm{ét}}_T S \geqslant 2$ 得出
\[
\mathrm{prof}^{\mathrm{ét}}_{(W - W_U)} W \geqslant 2
\]
(SGA 2 XIV 1.13)；从而若 $Z|_{W_U}$ 延拓成 $W$ 的平展覆盖，该延拓在唯一同构意义下唯一. 由此得出把 $Z$ 延拓到 $g(S) \cup X_U$ 的邻域的问题局部于 $g(T)$ 的点处的平展拓扑. 若 $t$ 是 $T$ 的点，令 $x = g(t)$，记 $X$ (resp. $S$) 为 $X$ 在 $x$ 处 (resp. $S$ 在 $t$ 处) 的严格局部化，$U = U \times_S S$，$X_U = X \times_X X_U$，$g \colon S \to X$ 为由 $g$ 导出的态射. 只需证明对 $T$ 的每个点 $t$，$Z$ 在 $X_U$ 上的逆像 $Z$ 延拓到 $X$ (或等价地平凡). 由 $Z$ 的定义，$Z$ 在 $U$ 上的逆像平凡. 为证明 $Z$ 平凡，只需证明它形如 $f^{*}_U E$，其中 $E$ 是 $U$ 的主覆盖；将有 $E \simeq g^{*}_U f^{*}_U E \simeq g^{*}_U Z$，从而由 $g^{*}_U Z$ 平凡得出结果. 由于态射 $f_U$ 0-无环且局部 0-无环，为证明 $Z$ 来自 $U$，只需证明对位于 $U$ 的某点上方 (可假设是点 $s$) 的每个几何点，$Z|_{X_s}$ 平凡 (SGA 4 XV 1.15). 但 $Z|_{X_s}$ 由群为 $L$-群的主覆盖经结构群扩张得到，由态射 $f$ 对 $L$ 1-非球得出结果.

于是证明存在 $g(S) \cup X_U$ 的开邻域 $V$，使 $Z$ 延拓成 $V$ 的平展覆盖 $Z_V$. 证明 $Z_V$ 延拓到整个 $X$. 只需看对 $X - V$ 的每个点 $x$，有
\[
\mathrm{prof}^{\mathrm{hop}}_x X \geqslant 3.
\]
这由假设 e) 和 $X - V$ 的点 $x$ 不能是它的纤维 $X_t$ 的极大点的事实得出，因为 $X_t$ 的每个不可约分量含 $g(t)$，$X_t$ 的每个极大点属于 $V$.

\begin{corollary}\label{cor:XIII.4.8}
假设同 \ref{prop:XIII.4.7.4}，但再假设 $\pi_1(S, a) = 1$. 则有同构
\[
(\ast) \quad \pi^L_1(X, a) \xrightarrow{\sim} \pi^L_1(X_s, a)_K.
\]
特别，若 $\pi^L_1(X_s, a)$ 拓扑有限表现，且 $K$ 经由有限型群作用在 $\pi^L_1(X_s, a)$ 上，则 $\pi^L_1(X, a)$ 拓扑有限表现.
\end{corollary}

同构 $(\ast)$ 已在 \ref{prop:XIII.4.7.4} 中证明. 假设 $\pi^L_1(X_s, a)$ 是 $n$ 个生成元的自由 pro-$L$-群 $L(x_1, \ldots, x_n)$ 除以由 $L(x_1, \ldots, x_n)$ 的元素 $y_1, \ldots, y_p$ 生成的闭不变子群的商，且 $K$ 经由由元素 $k_1, \ldots, k_q$ 生成的群作用. 若对每个 $i \in [1, n]$，$j \in [1, q]$，记 $z_{ij}$ 为 $L(x_1, \ldots, x_n)$ 的提升元素 $(k_j \cdot x_i)x^{-1}_i$，则 $\pi^L_1(X_s, a)_K$ 是 $L(x_1, \ldots, x_n)$ 除以由元素 $(y_i)_{i \in [1,p]}$，$(z_{ij})_{i \in [1,n], j \in [1,q]}$ 生成的闭不变子群的商.

\begin{remark}\label{rem:XIII.4.9}
a) 当 $S$ 是连通正规概形，$U$ 是 $S$ 的稠密拟紧开集，$f$ 是有限表现紧合态射、在 $T$ 的每个点 $t$ 处几何纤维连通且不可约，$f$ 还在 $X_U \cup g(T)$ 的点处可分光滑，$L$ 为与 $S$ 的剩余特征不同的素数集，$X$ 在 $X_t$ 的每个点处正则时，4.7 的条件 a) 至 e) 满足. 条件 a) 确实由 SGA 4 XV 4.1 和 \ref{exa:XIII.1.4} 得出；条件 b) 和 d) 由 \ref{exa:XIII.1.4} 和 SGA 4 XV 2.1 及 XVI 5.2 得出. 最后 e) 由 SGA 2 XIV 1.11 得出.

b) 推论 4.5 可以用来计算特征 $p$ 的可分闭域 $k$ 上紧合光滑曲面 $X$ 的基本群 $\pi^{p'}_1(X)$ ($p'$ 表示与 $p$ 不同的素数集). 该方法由 J.P. Murre 告诉我们；它归结到 (爆破 $X$ 后) 有纤维化 $X \to \mathbf{P}^1_k$ 和满足 4.7 的假设的 $\mathbf{P}^1_k$ 的开集 $U$ 的情形 (参见 SGA 7 的细节). 同样方法可以更一般地 (loc. cit.) 用来证明：若 $X$ 是连通有限型 $k$-概形，且 $k$ 的代数闭包上维数 $\leqslant \dim X$ 的有限型概形强可消解奇点 (SGA 5 I 3.1.5)，则 $\pi^{p'}_1(X)$ 拓扑有限表现.
\end{remark}

\section{附录 I：Abhyankar 引理的变体}

本附录含 Abhyankar 引理的不同变体.

\begin{proposition}\label{prop:XIII.5.1}
设 $X = \mathrm{Spec}\, A$ 是正则局部概形，$D = \sum_{1 \leqslant i \leqslant r} \mathrm{div}\, f_i$ 是正规交叉除子，其中诸 $f_i$ 是 $A$ 的极大理想的属于正则参数系的元素. 设诸 $n_i$ ($1 \leqslant i \leqslant r$) 是 $\geqslant 0$ 的整数，令
\[
X' = X[T_1, \ldots, T_r]/(T_1^{n_1} - f_1, \ldots, T_r^{n_r} - f_r),
\]
$U' = U \times_X X'$. 则 $X'$ 正则，且 $U'$ 是 $X'$ 中正规交叉除子 $\sum_{1 \leqslant i \leqslant r} \mathrm{div}\, T_i$ 的余集. 若诸整数 $n_i$ 与 $X$ 的剩余特征 $p$ 互素，$U'$ 是 $U$ 的相对 $D$ 温和分歧 (2.3 c)) 的连通平展覆盖.
\end{proposition}

确实 $X'$ 是极大理想由 $T_1, \ldots, T_r$ 生成的局部环 $A'$ 的谱. (可假设 $r = \dim(A)$.) 由于 $A'$ 在 $A$ 上有限且平坦，故维数 $r$，$A'$ 正则 (EGA IV$_0$ 17.1.1)，诸 $T_i$ 构成 $A'$ 的正则参数系. 假设诸 $n_i$ 与 $p$ 互素. 由于诸 $f_i$ 在 $U$ 上可逆，$U'$ 在 $U$ 上平展的事实由 I 7.4 得出. 而且 $U'$ 相对 $D$ 温和分歧；设 $x_i$ 是 $V(f_i)$ 的生成点，$R$ 是 $R = \mathcal{O}_{X, x_i}$ 的严格局部化，$K$ 是 $R$ 的分式域. 则表示 $U'|_K$ 的 $K$-代数由域 $K[T_i]/(T_i^{n_i} - f_i)$ 经非分歧扩张得到；故它在 $R$ 上温和分歧.

\begin{proposition}\label{prop:XIII.5.2}
[\emph{绝对 Abhyankar 引理}]\index[term]{juedui Abhyankar yinli@绝对 Abhyankar 引理} 设 $X$ 是正则局部概形，
\[
D = \sum_{1 \leqslant i \leqslant r} \mathrm{div}\, f_i
\]
是如 \ref{prop:XIII.5.1} 的正规交叉除子，$Y = \mathrm{Supp}\, D$，$U = X - Y$. 设 $V$ 是 $U$ 的相对 $D$ 温和分歧的平展覆盖. 若 $x_i$ 是闭集 $V(f_i)$ 的生成点，$\mathcal{O}_{X, x_i}$ 是分式域 $K_i$ 的离散赋值环，且 $V|_{K_i} = \mathrm{Spec}(\prod_{j \in J_i} L_j)$，其中诸 $L_j$ 是 $K_i$ 的有限可分扩张；记 $n_j$ 为 $L_j$ 生成的 伽罗瓦扩张的惯性群的阶，$n_i$ 为 $j$ 遍历 $J_i$ 时诸 $n_j$ 的最小公倍数. 若令
\[
X' = X[T_1, \ldots, T_r]/(T_1^{n_1} - f_1, \ldots, T_r^{n_r} - f_r),
\]
$U' = U_{(X')}$，$V' = V_{(X')}$，等等，$U'$ 的平展覆盖 $V'$ 唯一地 (模唯一同构) 延拓成 $X'$ 的平展覆盖，且诸 $n_i$ 与 $X$ 的剩余特征 $p$ 互素.
\end{proposition}

唯一性由 $X'$ 正规 (\cref{prop:XIII.5.1}) 得出；确实，延拓 $V'$ 的 $X'$ 的平展覆盖同构于 $X'$ 在 $V'$ 在 $X'$ 的生成点处的纤维中的正规化 (\cref{cor:I.10.2}). 若 $x'$ 是 $Y'$ 的几何点，记 $X'$ 为 $X'$ 在 $x'$ 处的严格局部化，$V' = V'_{(X')}$，等等. 经下降，考虑到唯一性，只需证明对 $Y'$ 的每个几何点 $x'$，$U'$ 的平展覆盖 $V'$ 延拓到 $X'$. 给定正则概形 $X'$ 的含所有满足 $\dim \mathcal{O}_{X', x'} \leqslant 1$ 的点 $x'$ 的开集的平展覆盖延拓到整个 $X'$ (SGA 2 XIV 1.11)，可限于 $x'$ 在 $Y'$ 的极大点上的投影. 在这样的点 $x'$，$V'$ 延拓成 $X'$ 的平展覆盖的事实由 3.6 得出.

证明诸 $n_i$ 与 $p$ 互素. 若不然，有例如 $p|n_1$. 把 $X$ 换成
\[
X[T_1, \ldots, T_r]/(T_1^{n_1/p} - f_1, T_2^{n_2} - f_2, \ldots, T_r^{n_r} - f_r),
\]
归结到 $X' = X[T_1]/(T_1^p - f_1)$ 的情形. 只需证明 $V$ 延拓成 $X$ 的平展覆盖，因为此时将有 $n_1 = 1$，与假设矛盾. 为此可假设 $X$ 严格局部. 设 $Z$ 是 $X$ 的方程 $p = 0$ 的闭子概形，$Z_1 = Z \cap X_{f_1}$；$Z_1$ 是 $Z$ 的非空开集. 由前面，$U'$ 的平展覆盖 $V'$ 延拓成 $X'$ 的平展覆盖 $W'$. 设 $W''_1$ 和 $W''_2$ 是 $W'$ 经两个投影 $X'' = X' \times_X X' \rightrightarrows X'$ 的逆像，证明下降同构 $u \colon W''_1|_{U''} \to W''_2|_{U''}$ 延拓成 $X''$-态射 $W''_1 \to W''_2$，它必然是 $W'$ 上相对 $X' \to X$ 的下降数据. 设 $Z''$ (resp. $Z''_1$) 是 $Z$ (resp. $Z_1$) 在 $X''$ 中的逆像. 由于态射 $Z'' \to Z$ 根式，存在延拓同构 $u|_{Z''_1}$ 的同构 $v \colon W''_1|_{Z''} \to W''_2|_{Z''}$. 但 $X$ Hensel 局部，有双射
\[
\mathrm{Hom}_{X''}(W''_1, W''_2) \simeq \mathrm{Hom}_{Z''}(W''_1|_{Z''}, W''_2|_{Z''}),
\]
从而有提升 $v$ 的态射 $w \colon W''_1 \to W''_2$. $X''$ 的 $u$ 和 $w$ 重合的同时开且闭子概形含 $Z''_1$，故等于整个 $X''$，由此得出 $V$ 延拓到 $X$.

5.3.0. 沿用 \ref{prop:XIII.5.2} 的假设和记号，再假设 $X$ 严格局部. 由 loc. cit. 得出 $U$ 的每个相对 $D$ 温和分歧的连通平展覆盖是形如
\[
U' = U[T_1, \ldots, T_r]/(T_1^{n_1} - f_1, \ldots, T_r^{n_r} - f_r)
\]
的温和分歧覆盖的商，其中诸 $n_i$ 是与 $p$ 互素的整数. 设 $\mu_n$ 是 $U$ 的 $n$ 次单位根群. $U'$ 的 $U$-自同构群正是 $\mu_{n_1} \times \cdots \times \mu_{n_r}$，其中 $n_i$ 次单位根 $\xi_i$ 把 $T_i$ 变成 $\xi_i T_i$ 地作用在 $U'$ 上. 故有下列断言：

\begin{corollary}\label{cor:XIII.5.3}
设 $X$ 是剩余特征 $p \geqslant 0$ 的严格局部正则概形，$D = \sum_{1 \leqslant i \leqslant n} \mathrm{div}\, f_i$ 是 $X$ 上的正规交叉除子，$U = X - \mathrm{Supp}\, D$. 令
\[
\widetilde{U} = \varprojlim_{(n_i)} U[T_1, \ldots, T_r]/(T_1^{n_1} - f_1, \ldots, T_r^{n_r} - f_r),
\]
投射极限按与 $p$ 互素的整数族 $n_i > 0$ (对整除关系) 的滤过有序集取. 则 $\widetilde{U}$ 是 $U$ 的温和分歧万有覆盖. 从而 $U$ 的温和分歧基本群为
\[
\pi^t_1(U) \simeq \prod_{\ell \neq p} \mathbb{Z}_\ell[1]^r \quad \text{(典范同构)}
\]
其中令 $\mathbb{Z}_\ell[1] = \varprojlim_{n > 0} \mu_{\ell^n}$\index[notations]{$\mathbb{Z}_\ell[1]$}. 群 $\pi^t_1(U)$ 非典范同构于 $\prod_{\ell \neq p} \mathbb{Z}_\ell^r$.
\end{corollary}

\begin{proposition}\label{prop:XIII.5.4}
设 $f \colon X \to S$ 是概形态射，$D = \sum_{1 \leqslant i \leqslant r} \mathrm{div}\, f_i$ 是相对 $S$ 正规交叉 (2.1) 的除子，其中对 $Y = \mathrm{Supp}\, D$ 的每个点 $x$，若 $I(x) \subset [1, r]$ 是使 $f_i(x) = 0$ 的 $i$ 的集，子概形 $V((f_i)_{i \in I(x)})$ 在 $S$ 上光滑且在 $X$ 中余维数 $\mathrm{card}\, I(x)$. 设 $U = X - Y$. 设 $x$ 是 $Y$ 的点，$X_1 = \mathrm{Spec}\, \mathcal{O}_{X, x}$，$U_1 = U \times_X X_1$，诸 $n_i$ ($i \in I(x)$) 是整数，
\[
X' = X[T_i]_{i \in I(x)}/(T_i^{n_i} - f_i).
\]
则若 $x'$ 是 $X'$ 的位于 $x$ 上方的点，$X'$ 在 $x'$ 处平展于 $S$ 上. 若诸整数 $n_i$ 与 $k(x)$ 的特征 $p$ 互素，$U'_1 = U_1 \times_X X'$ 是 $U_1$ 的相对 $S_1$ 在 $X_1$ 上温和分歧 (\cref{def:XIII.2.1.1}) 的连通平展覆盖.
\end{proposition}

若 $s = f(x)$，几何纤维 $X'_s$ 在点 $x'$ 处正则 (\cref{prop:XIII.5.1})；由于 $X'$ 在 $Y$ 的邻域内在 $S$ 上平坦，这证明 $X'$ 在 $x'$ 处平展于 $S$ 上 (EGA IV 12.1.6). 若诸整数 $n_i$ 与 $p$ 互素，$U'_1$ 是 $U_1$ 的平展覆盖 (\cref{cor:I.7.4})；它相对 $S$ 在 $X_1$ 上温和分歧，因为在 $S$ 的每个点处的几何纤维上如此 (\cref{prop:XIII.5.1}). 最后 $U'_1$ 连通的事实由 (SGA 4 XVI 3.2) 得出.

\begin{proposition}\label{prop:XIII.5.5}
[\emph{相对 Abhyankar 引理}]\index[term]{xiangdui Abhyankar yinli@相对 Abhyankar 引理} 设 $X$ 是 $S$-概形，$D$ 是如 (5.4) 的相对 $S$ 正规交叉除子. 设 $Y = X - \mathrm{Supp}\, D$，$U = X - Y$，$x$ 是 $Y$ 的点，$X_1$ 是 $X$ 在 $x$ 上方的几何点处的严格局部化，$U_1 = U \times_X X_1$，$V_1$ 是 $U_1$ 的平展覆盖. 假设对 $S$ 的每个极大点 $s$，$V_{1s}$ 相对 $s$ 在 $X_{1s}$ 上温和分歧. 则可以找到与 $k(x)$ 的特征 $p$ 互素、$i \in I(x)$ 的整数 $n_i$，使若令
\[
X'_1 = X_1[T_i]_{i \in I(x)}/(T_i^{n_i} - f_i),
\]
$U'_1 = U_1 \times_{X_1} X'_1$，等等，$U'_1$ 的平展覆盖 $V'_1$ 唯一地 (模唯一同构) 延拓成 $X'_1$ 的平展覆盖. 特别 $V_1$ 相对 $S$ 在 $X_1$ 上温和分歧.
\end{proposition}

可假设 $S$ 是闭点 $f(x)$ 的局部诺特局部环的谱. 对 $S$ 的每个极大点 $s$ 和每个 $i \in I(x)$，设 $x_i$ 是纤维 $X_{1s}$ 的闭集 $V(f_i)$ 的生成点. 局部环 $(\mathcal{O}_{X_1, x_i})_{\mathrm{r\acute{e}d}}$ 是分式域 $K_i$ 的离散赋值环，且 $V_1|_{K_i} = \mathrm{Spec}(\prod_{j \in I(x_i)} L_j)$，其中 $L_j$ 是 $K$ 的有限可分扩张；记 $n_j$ 为 $L_j$ 生成的 伽罗瓦扩张的惯性群的阶，$n_i$ 为 $s$ 遍历 $S$ 的极大点、$j \in I(x_i)$ 时诸 $n_j$ 的最小公倍数.

$n_i$ 这样选定后，证明 $V'_1$ 唯一地延拓成 $X'_1$ 的平展覆盖. 唯一性由 $X'$ 在 $Y$ 的点处平展于 $S$ 上 (故 $\mathrm{prof}^{\mathrm{ét}}_{Y'_1}(X'_1) \geqslant 2$ (SGA 4 XVI 3.2 或 SGA 2 XIV 1.19)) 得出. 设 $x'_1$ 是 $Y'_1$ 的点，$x'_1$ 是 $x'_1$ 上方的几何点，记 $X'_1$ 为 $X'_1$ 在 $x'_1$ 处的严格局部化，$U'_1 = U'_{1(X'_1)}$，等等. 经下降，考虑到唯一性，只需证明 $V'_1$ 延拓到 $X'_1$. 而且可以限于 $x'_1$ 为 $Y'_1$ 的极大点；确实此时在 $X'_1$ 的含 $Y'_1$ 的极大点的开集 $W'_1$ 上有 $V'_1$ 的延拓；若 $Z'_1 = X'_1 - W'_1$，当 $s$ 是 $S$ 的极大点时 $\mathrm{codim}(Z'_{1s}, X'_{1s}) \geqslant 2$，当 $s$ 是非极大的 $S$ 的点时 $\mathrm{codim}(Z'_{1s}, X'_{1s}) \geqslant 1$，$\mathrm{prof}^{\mathrm{ét}}_s(S) \geqslant 1$；$V'_1$ 延拓到整个 $X'_1$ 的事实由 (SGA 2 XIV 1.20) 得出. 而在位于 $Y'_1$ 的极大点上方的几何点 $x'_1$ 处，$X'_{1\,\mathrm{r\acute{e}d}}$ 是离散赋值环的谱，$V'_1$ 延拓到 $X'_1$ 的事实由 X 3.6 得出.

证明诸 $n_i$ 与 $p$ 互素. 若不然，有使 $p$ 整除 $n_{i_0}$ 的指标 $i_0 \in I(x)$. 把 $X$ 换成
\[
X_1[T_{i_0}, T_i]_{i \in I(x)}/(T_{i_0}^{n_{i_0}/p} - f_{i_0}, T_i^{n_i} - f_i),
\]
归结到 $X'_1 = X_1[T]/T^p - f_{i_0}$ 的情形. 由前面，$U'_1$ 的平展覆盖 $V'_1$ 延拓成 $X'_1$ 的平展覆盖 $E'_1$. 设 $\eta$ 是 $S$ 的闭点；由于态射 $X'_{1\eta} \to X_{1\eta}$ 根式，$V_{1\eta}$ 延拓成 $X_{1\eta}$ 的平展覆盖 $E_{1\eta}$. 由此如 \ref{prop:XIII.5.2} 得出，$E'_1$ 配备相对态射 $X'_1 \to X_1$ 的下降数据，它延拓 $E'_1|_{U'_1}$ 上的自然下降数据. 由此得出 $V_1$ 延拓到 $X_1$；但这蕴含 $n_{i_0} = 1$，与假设 $n_{i_0} = p$ 矛盾.

\begin{corollary}\label{cor:XIII.5.6}
设 $X$ 是 $S$-概形，$D = \sum_{1 \leqslant i \leqslant r} \mathrm{div}\, f_i$ 是如 \ref{prop:XIII.5.4} 的相对 $S$ 正规交叉除子. 设 $x$ 是 $X$ 的几何点，$X$ 是 $X$ 在 $x$ 处的严格局部化，$Y = Y_{(X)}$，$U = X - Y$，
\[
\widetilde{U} = \varprojlim_{(n_i)} U[T_i]_{i \in I(x)}/(T_i^{n_i} - f_i),
\]
投射极限按与 $k(x)$ 的特征 $p$ 互素的整数族 $n_i > 0$ 的滤过集取. 则 $\widetilde{U}$ 是 $U$ 的相对 $S$ 温和分歧万有覆盖. 从而 $U$ 的温和分歧基本群为
\[
\pi^t_1(U) \simeq \prod_{\ell \neq p} \mathbb{Z}_\ell[1]^{I(x)} \quad \text{(典范同构)}.
\]
群 $\pi^t_1(U)$ 非典范同构于 $\prod_{\ell \neq p} \mathbb{Z}_\ell^{I(x)}$.
\end{corollary}

\begin{remark}\label{rem:XIII.5.6.1}
设 $X$ 是 $S$-概形，$D = \sum_{1 \leqslant i \leqslant r} \mathrm{div}\, f_i$ 是如 \ref{prop:XIII.5.4} 的相对 $S$ 正规交叉除子，$U = X - \mathrm{Supp}\, D$. 对每个部分 $I \subset [1, r]$，设
\[
X_I = \Big(\bigcap_{i \in I} V(f_i)\Big) \cap \Big(\bigcap_{i \notin I} X_{f_i}\Big).
\]
设 $p$ 是素数或零，$Z$ 是 $X_I$ 的诸点特征都为 $p$ 的子集. 设
\[
\widetilde{U}_I = \varprojlim_{(n_i)} U[T_i]_{i \in I}/(T_i^{n_i} - f_i),
\]
其中投射极限按与 $p$ 互素的整数族 $n_i > 0$ 的滤过集取. 则对 $Z$ 的每个几何点 $x$，$\widetilde{U}_I$ 在 $U$ 上的逆像等同于 $U$ 的温和分歧万有覆盖.
\end{remark}

\begin{corollary}\label{cor:XIII.5.7}
记号同 5.6. 设 $S$ 是 $S$ 在 $x$ 处的严格局部化，
\[
g \colon U \longrightarrow S \quad \text{和} \quad \widetilde{g} \colon \widetilde{U} \longrightarrow \widetilde{S}
\]
是典范态射. 则态射 $g$ 和 $\widetilde{g}$ 0-无环 (SGA 4 XV 1.3). 设 $G$ 是 $S$ 上的可构群层，$F = g^{*} G$，$P$ 是 $F$ 下的挠子. 则 $P$ 相对 $S$ 在 $X$ 上温和分歧，当且仅当它在 $\widetilde{U}$ 上的逆像 $\widetilde{P}$ 平凡.
\end{corollary}

确实，对每个概形 $X' = X[T_i]_{i \in I(x)}/(T_i^{n_i} - f_i)$ (其中诸 $n_i$ 是与 $p$ 互素的整数 $> 0$)，态射 $f' \colon X' \to S$ 0-无环. $f'$ 在 $S$ 的各点处的几何纤维因此连通甚至不可约. 态射 $g' \colon U' \to S$ 的几何纤维也如此，这证明诸 $g'$，从而 $\widetilde{g}$，0-无环 (SGA 4 XV 1.16).

$U$ 上群 $F$ 的在 $\widetilde{U}$ 上的逆像平凡的挠子 $P$ 显然相对 $S$ 在 $X$ 上温和分歧. 证明反之，若 $P$ 相对 $S$ 在 $X$ 上温和分歧，它在 $\widetilde{U}$ 上的逆像平凡.

由 (SGA 4 IX 2.14 (ii)) 得出可以找到有限态射 $n \colon S_1 \to S$、$S_1$ 上的常群层 $C$ 和单态射 $G \to n_* C$. 考虑下列所有方形笛卡儿的交换图式：
\[
\begin{tikzcd}[row sep=large, column sep=large]
\widetilde{U}_1 \arrow[r] \arrow[d,"{r}"] & U_1 \arrow[r,"{g_1}"] \arrow[d,"{q}"] & S_1 \arrow[d,"{n}"] \\
\widetilde{U} \arrow[r] & U \arrow[r,"{g}"] & S.
\end{tikzcd}
\]
设 $C_1$ (resp. $\widetilde{C}_1$) 是 $C$ 在 $U_1$ (resp. $\widetilde{U}_1$) 上的逆像. 有 $i$ 和 $j$ 为同构 (SGA 4 VIII 5.8) 的交换图式：
\begin{equation}
\begin{tikzcd}[row sep=large, column sep=large]
H^1(U, q_* C_1) \arrow[r,"{i}"] \arrow[d] & H^1(U_1, C_1) \arrow[d,"{j}"] \\
H^1(\widetilde{U}, r_* \widetilde{C}_1) \arrow[r] & H^1(\widetilde{U}_1, \widetilde{C}_1).
\end{tikzcd}
\tag{$\ast$}
\end{equation}
设 $Q$ 是由 $P$ 经结构群扩张 $F \to q_* C_1$ 得到的 $q_* C_1$ 下挠子. 由 2.1.4，$Q$ 相对 $S$ 在 $X$ 上温和分歧. 与挠子 $Q$ 对应 (借助 $i$) $C_1$ 下的挠子 $Q_1$，且显然 $Q_1$ 相对 $S_1$ 在 $X_1 = X \times_S S_1$ 上温和分歧. 由 \ref{cor:XIII.5.6} 得出 $Q_1$ 在 $\widetilde{U}_1$ 上的逆像 $\widetilde{Q}_1$ 平凡，图式 ($\ast$) 于是表明 $Q$ 在 $\widetilde{U}$ 上的逆像 $\widetilde{Q}$ 平凡.

考虑下列第二行正合 (SGA 4 XII 3.1) 的交换图式：
\begin{equation}
\begin{tikzcd}[row sep=large, column sep=large]
H^0(S, n_* C/G) \arrow[r] \arrow[d,"{k}"] & H^1(S, G) \arrow[r, equal] & 1 \arrow[d] \\
H^0(\widetilde{U}, r_* \widetilde{C}_1/\widetilde{F}) \arrow[r] & H^1(\widetilde{U}, \widetilde{F}) \arrow[r] & H^1(\widetilde{U}, r_* \widetilde{C}_1).
\end{tikzcd}
\tag{$\ast\ast$}
\end{equation}
由于态射 $\widetilde{U} \to S$ 0-无环，$k$ 是同构. $\widetilde{P}$ 平凡的事实于是由 ($\ast\ast$) 得出.

\section{附录 II：叠的直像的有限性定理}

\begin{proposition}\label{prop:XIII.6.1}
设 $S$ 是局部诺特概形，$f \colon X \to S$ 是态射. 若 $S'$ 是 $S$-概形，记 $X'$ (resp. $f'$，等等) 为 $X$ (resp. $f$，等等) 经态射 $S' \to S$ 的逆像. 假设对每个平展于 $S$ 上的概形 $S'$ 和 $X'$ 上每个可构集合层 $F$，$f'_* F$ 可构，且对 $X'$ 上每个可构群层 $F$，$R^1 f'_* F$ 可构. 设 $\Phi$ 是 $X$ 上 (0) 的 1-可构叠. 则 $f_* \Phi$ 1-可构.
\end{proposition}

对每个平展于 $S$ 上的概形 $S'$ 和 $(f_* \Phi)_{S'}$ 的每个对象 $x$，有同构
\[
\mathrm{Aut}_{S'}(x) \simeq f'_* \mathrm{Aut}_{X'}(x),
\]
其中第二个成员中 $x$ 看作 $\Phi_{X'}$ 的对象. 所做假设于是蕴含 $f_* \Phi$ 可构. 设 $S_\Phi$ 是 $\Phi$ \cite[III 2.1.7]{exp13-1} 的极大子叠丛层. 由于 $f_*(S_\Phi)$ 可构，可以给它应用 (SGA 4 IX 2.7)，$f_* \Phi$ 1-可构的事实于是由下列引理得出.

\begin{lemma}\label{lem:XIII.6.1.1}
设 $S$ 是局部诺特概形，$f \colon X \to S$ 是态射，$\Phi$ 是 $X$ 上的叠. 假设给定 $S$ 上由有限型平展 $S$-概形 $T$ 表示的层、满态射
\[
a \colon T \longrightarrow f_*(S_\Phi)
\]
和纤维 $\Phi_{X_T}$ (其中 $X_T = X \times_S T$) 的对象 $p$，它在 $f_*(S_\Phi)(T) = S_\Phi(X_T)$ 中定义等于 $T(T)$ 的恒等截面经 $a$ 的像 $q$ 的元素. 设 $f_T \colon X_T \to T$ 是典范态射，假设层 $R^1 f_{T*}(\mathrm{Aut}_{X_T}(p))$ 可构；则 $S(f_* \Phi)$ 也可构.
\end{lemma}

典范态射 $f^{*} f_* \Phi \to \Phi$ 给出态射
\[
S(f^{*} f_* \Phi) \simeq f^{*}(S(f_* \Phi)) \longrightarrow S_\Phi,
\]
从而有典范态射
\[
\phi \colon S(f_* \Phi) \longrightarrow f_*(S_\Phi).
\]
设 $F = S(f_* \Phi)$，$G$ 是 $F$ 经 $\phi$ 的像；由 (SGA 4 IX 2.9)，$G$ 是可构层.

只需证明对 $S$ 的每个点 $s$，存在 $s$ 的非空开集 $U$，使 $F|_U$ 局部常可构. 设 $s \in S$，$s$ 是 $s$ 上方的几何点，$q_1, \ldots, q_n$ 是 $G_s$ 的元素. 由 $T$ 的定义，存在 $S$-态射 $h_i \colon s \to S'$，使 $q_i = h^{*}_i(q)$. 设 $S'$ 是 $n$ 个同构于 $T$ 的概形在 $S$ 上的纤维积，$s \to S'$ 是诸 $h_i$ 的纤维积，$X' = X \times_S S'$，$q_i$ (resp. $p_i$) 是 $q$ (resp. $p$) 经 $S'$ 到 $T$ 的第 $i$ 个投影的逆像. 若 $\Psi_i$ 是 $p_i$ 生成的 $\Phi|_{X'}$ 的极大子叠丛，层 $F_i = R^1 f'_*(\mathrm{Aut}_{X'}(p_i))$ 正是 $f'_* \Psi_i$ 的极大子叠丛层 $S(f'_* \Psi_i)$. 特别典范单射 $\Psi_i \to \Phi|_{X'}$ 给出态射
\[
\alpha_i \colon F_i \longrightarrow F|_{S'}.
\]
证明 $\alpha_i$ 是 $F_i$ 到 $q_i$ 在 $F|_{S'}$ 中的逆像的双射. 对每个平展于 $S'$ 上的概形 $S''$，$F_i(S'')$ 的每个截面 $y$ 在 $F(S'')$ 中的像是 $q_i|_{S''}$，因为局部于 $S''$ 的平展拓扑，$y$ 由同构于 $p_i|_{X''}$ 的 $\Phi_{X''}$ 的对象 $x$ 定义. 反之，若 $y \in F(S'')$ 在 $F(S'')$ 中的像是 $q_i|_{S''}$，局部于 $S''$ 的平展拓扑，$y$ 由同构于 $p_i$ 的 $\Phi_{X''}$ 的对象 $x$ 定义；从而 $x$ 是 $\Psi_{iX''}$ 的对象，故 $y \in F_i(S'')$.

用下面 \ref{lem:XIII.6.1.2} 完成证明. 可以找到 $s$ 的开邻域 $U'$，使 $q_1|_{U'}, \ldots, q_n|_{U'}$ 是 $G(U')$ 的截面并生成该层. 由于诸 $F_i|_{U'}$ 和 $G|_{U'}$ 可构，由 \ref{lem:XIII.6.1.2} 得出 $F|_{U'}$ 也可构；把 $U$ 换成更小的开集后 $F|_U$ 局部常，完成证明.

\begin{lemma}\label{lem:XIII.6.1.2}
设 $S$ 是局部诺特概形，$F \to G$ 是 $S$ 上群层的满态射. 设 $q_i$ 是 $X$ 上生成 $G$ 的 $G$ 截面的有限族，对每个 $i$，设 $F_i$ 是 $F$ 的 $q_i$ 的逆像子层. 则若 $G$ 和诸 $F_i$ 可构，$F$ 也可构.
\end{lemma}

为证明 $F$ 可构，只需证明对 $S$ 的每个点 $s$，存在 $s$ 的开邻域 $U$，使 $F|_U$ 局部常可构. 设 $s$ 是 $S$ 的点. 由于层 $F_i$ 和 $G$ 可构，可以找到 $s$ 的开邻域 $U$，使 $F_i|_U$ 和 $G|_U$ 局部常. 证明 $F|_U$ 局部常. 由 (SGA 4 IX 2.13 (i))，只需看：若 $s$ 是 $s$ 上方的几何点，$\widetilde{s}$ 是 $U$ 的几何点，$s \to \widetilde{s}$ 是特殊化态射，典范态射
\[
F_{\widetilde{s}} \longrightarrow F_s
\]
双射.

考虑交换图式
\[
\begin{tikzcd}[row sep=large, column sep=large]
(F_i)_{\widetilde{s}} \arrow[r,"{\widetilde{a}_i}"] \arrow[d,"{\simeq}"] & F_{\widetilde{s}} \arrow[r,"{\widetilde{a}}"] \arrow[d,"{b}"] & G_{\widetilde{s}} \arrow[d,"{\simeq}"] \\
(F_i)_s \arrow[r,"{a_i}"'] & F_s \arrow[r,"{a}"'] & G_s
\end{tikzcd}
\]
设 $q_i$ (resp. $\widetilde{q}_i$) 是 $q_i$ 在 $G_s$ (resp. $G_{\widetilde{s}}$) 中的逆像；态射 $a$ 和 $\widetilde{a}$ 满射，态射 $a_i$ (resp. $\widetilde{a}_i$) 诱导 $(F_i)_s$ (resp. $(F_i)_{\widetilde{s}}$) 到 $a^{-1}(q_i)$ (resp. $\widetilde{a}^{-1}(q_i)$) 的双射. 由上面图式得出 $b$ 是同构.

\begin{corollary}\label{cor:XIII.6.2}
设 $S$ 是局部诺特概形，$f \colon X \to S$ 是紧合态射. 设 $\Phi$ 是 $X$ 上 1-可构叠，则 $f_* \Phi$ 是 1-可构叠.
\end{corollary}

\ref{prop:XIII.6.1} 的证明也证明下列结果，考虑到 \ref{thm:XIII.2.4} 2)：

\begin{corollary}\label{cor:XIII.6.3}
设 $S$ 是局部诺特概形，$f \colon X \to S$ 是态射，$D$ 是 $X$ 上相对 $S$ 正规交叉 (2.1) 的除子，$Y = \mathrm{Supp}\, D$，$U = X - Y$，$i \colon U \to X$ 是典范浸入. 设 $\Phi$ 是 $U$ 上的叠，局部于 $X$ 和 $S$ 的平展拓扑是 $S$ 上 1-可构叠 $\Psi$ 的逆像. 则叠 $i_{*t} \Phi$ 1-可构.
\end{corollary}

\section{文献}\printbibliography[heading=none]


